\documentclass[a4paper,UKenglish,numberwithinsect,pdfa]{lipics-v2021}

%\pdfoutput=1 %uncomment to ensure pdflatex processing (mandatatory e.g. to submit to arXiv)
%\hideLIPIcs  %uncomment to remove references to LIPIcs series (logo, DOI, ...), e.g. when preparing a pre-final version to be uploaded to arXiv or another public repository

% Our packages
\usepackage[inline]{enumitem}
\usepackage[b]{esvect} % variantes: [a], [b], [c], [e]...
\usepackage{proof}
\usepackage[only=llparenthesis,rrparenthesis,llbracket,rrbracket]{stmaryrd}
\SetSymbolFont{stmry}{bold}{U}{stmry}{m}{n}
\usepackage{tikz}
\usetikzlibrary{cd,positioning,decorations.text,decorations.pathmorphing}
\tikzcdset{scale cd/.style={every label/.append style={scale=#1},cells={nodes={scale=#1}}}}
\usepackage{tikz-cd}
\usepackage{anyfontsize}
\makeatletter
\newtheoremstyle{restatestyle}
  {0pt} % Espacio superior (se ajusta al entorno de LIPIcs)
  {0pt} % Espacio inferior
  {\itshape} % Cuerpo en itálica como pediste
  {} % Sangría
  {\bfseries} % Fuente del título
  {.} % Puntuación después del título
  {.5em} % Espacio después del título
  {#3} % <--- Esto hace que solo se use el argumento opcional (el 3er parámetro)
\theoremstyle{restatestyle}
\newtheorem*{restatetheorem}{}
\newtheorem*{restatecorollary}{}
\newtheorem*{restatelemma}{}
\makeatother

\allowdisplaybreaks

\bibliographystyle{plainurl}% the mandatory bibstyle

\title{Basis-Sensitive Quantum Typing via Realisability}

\author{Alejandro Díaz-Caro}
{
  Université de Lorraine, CNRS, Inria, LORIA, France
  \and
  Universidad de la República, CENUR Suroeste, Facultad de Ingeniería, InCo, Uruguay
}
{alejandro.diaz-caro@inria.fr}
{https://orcid.org/0000-0002-5175-6882}
{}
%{Supported by the Plan France 2030 through the PEPR integrated project EPiQ (ANR-21-PETQ-0007)}

\author{Octavio Malherbe}
{Universidad de la República, Facultad de Ingeniería, IMERL, Uruguay}
{malherbe@fing.edu.uy}
{https://orcid.org/0009-0004-0624-9285}
{}


\author{Rafael Romero}
{Universidad Nacional de Quilmes, Departamento de Ciencia y Tecnología, Argentina
  \and 
ICC, CONICET-Universidad de Buenos Aires, Argentina}
{lromero@dc.uba.ar}
{https://orcid.org/0000-0003-3186-1172}	
{}

\authorrunning{A. Díaz-Caro, O. Malherbe, and R. Romero}

\Copyright{Alejandro Díaz-Caro, Octavio Malherbe, and Rafael Romero}

\ccsdesc[500]{Theory of computation~Quantum computation theory}
\ccsdesc[500]{Theory of computation~Lambda calculus}
\ccsdesc[500]{Theory of computation~Program semantics}

\keywords{Quantum computing, Realisability semantics, \texorpdfstring{$\lambda$}{lambda}-calculus}

\relatedversion{} %optional, e.g. full version hosted on arXiv, HAL, or other respository/website
%\relatedversiondetails[linktext={opt. text shown instead of the URL}, cite=DBLP:books/mk/GrayR93]{Classification (e.g. Full Version, Extended Version, Previous Version}{URL to related version} %linktext and cite are optional

%\supplement{}%optional, e.g. related research data, source code, ... hosted on a repository like zenodo, figshare, GitHub, ...
%\supplementdetails[linktext={opt. text shown instead of the URL}, cite=DBLP:books/mk/GrayR93, subcategory={Description, Subcategory}, swhid={Software Heritage Identifier}]{General Classification (e.g. Software, Dataset, Model, ...)}{URL to related version} %linktext, cite, and subcategory are optional

%\funding{Supported by the European Union through the MSCA SE project QCOMICAL (Grant Agreement ID: 101182521) and by the Uruguayan CSIC grant 22520220100073UD.}%optional, to capture a funding statement, which applies to all authors. Please enter author specific funding statements as fifth argument of the \author macro.

\acknowledgements{}%optional

%\nolinenumbers %uncomment to disable line numbering



%Editor-only macros:: begin (do not touch as author)%%%%%%%%%%%%%%%%%%%%%%%%%%%%%%%%%%
\EventEditors{Michal Koucký and Daniela Petrisan}
\EventNoEds{1}
\EventLongTitle{51st International Symposium on Mathematical Foundations of Computer Science (MFCS 2026)}
\EventShortTitle{MFCS 2026}
\EventAcronym{MFCS}
\EventYear{2026}
\EventDate{August 24--28, 2026}
\EventLocation{Paris, France}
\EventLogo{}
\SeriesVolume{51}
\ArticleNo{}
%%%%%%%%%%%%%%%%%%%%%%%%%%%%%%%%%%%%%%%%%%%%%%%%%%%%%%

\newcommand{\lambdaSh}{\ensuremath{\text{Lambda-}{\sharp}}}
\newcommand{\lambdaSI}{\ensuremath{\text{Lambda-}\mathcal S_1}}

\newcommand\cmark{\tikz\fill[scale=0.4](0,.35) -- (.25,0) -- (1,.7) -- (.25,.15) -- cycle;} 

\newcommand\braces[1]{\left\{#1\right\}}
\newcommand\So{\ensuremath{{\mathcal S}_1}}

\newcommand\ket[1]{\ensuremath{|#1\rangle}}
\newcommand\parts[1]{\ensuremath{{\mathcal P}_{\!\!*}({#1})}}
\newcommand\Span[1]{\ensuremath{{\mathsf{Span}}{#1}}}
\newcommand\Forg[1]{\ensuremath{{U}{#1}}}
\newcommand\Definible{\mathsf{Def}}
\newcommand\comp[2][]{#2^{\bot^{#1}}}
\newcommand\Rpart[1]{\mathsf{Re(#1)}}
\newcommand\Ipart[1]{\mathsf{Im(#1)}}

\newcommand\s[1]{\ensuremath{\mathsf{#1}}}
\newcommand\Op[1]{\ensuremath{#1^{\mathsf{op}}}}
\newcommand\SSet{\s{Set}}
\newcommand\Var{\ensuremath{\mathsf{Var}}}
\newcommand\Val{{\s V}}
\DeclareRobustCommand{\ValD}{\ensuremath{\vv{\mathsf{V}}}}
\newcommand\VecV{\s{SVec}_{\ValD}}
\newcommand\SetV{\Set_{\ValD}}

\newcommand\interp[1]{\ensuremath{\llbracket #1 \rrbracket}}
\newcommand\Hom{\s{Hom}}
\newcommand\HomS[1]{\Hom_{\SetV}(#1)}
\newcommand\HomV[1]{\Hom_{\VecV}(#1)}
\newcommand\Ob[1]{\s{Ob}(#1)}
\newcommand\lra{\longrightarrow}
\newcommand\xlra[1]{\xrightarrow{#1}}
\newcommand\Hil{\mathcal H_{\ValD}}
\newcommand\SpFun[1]{\ensuremath{{S}{#1}}}
\newcommand\Id{\mathsf{Id}}

\newcommand\lambdaQ{\ensuremath{\lambda^Q}}
\newcommand\lambdaS{\ensuremath{\lambda_{\mathcal{S1}}}}
\newcommand\lambdaH{\ensuremath{\lambda_{\sharp}}}
\newcommand\inl[1]{\ensuremath{\mathsf{inl}(#1)}}
\newcommand\inr[1]{\ensuremath{\mathsf{inr}(#1)}}
\newcommand\qlet[3]{\ensuremath{\mathsf{let }#1 = #2\ \mathsf{ in }\ #3}}
\newcommand\qmatch[5]{\ensuremath{\mathsf{match}\ #1\ \{\inl{#2}\mapsto #3\ |\ \inr{#4} \mapsto #5 \}}}
\newcommand\sspan[1]{\ensuremath{\mathsf{span}(#1)}}
\newcommand\ansubst[2]{\ensuremath{\langle #1 \rangle_{#2}}}
\newcommand\sqsubst[1]{\ensuremath{\left[ #1 \right]}}
\newcommand\AbsBasis{\ensuremath{\mathsf{H\!O}}}
\newcommand\dom[1]{\mathrm{dom}(#1)}
\newcommand\sdom[1]{\mathrm{dom}^{\sharp}(#1)}
\newcommand\FV[1]{\mathrm{FV}(#1)}
%%% Logic
\def\limp{\Rightarrow}
\def\liff{\Leftrightarrow}
%%% Sets
\def\N{\mathbb{N}}            % set of natural numbers
\def\R{\mathbb{R}}            % set of real numbers
\def\C{\mathbb{C}}            % set of complex numbers
\def\Pow{\mathfrak{P}}        % powerset
\def\Powfin{\Pow_{\text{fin}}}  % set of finite subsets
\def\X{\mathcal{X}}           % set of variables
\def\Val{\mathrm{V}}          % set of basic values
\def\pto{\rightharpoonup}     % set of partial functions
\def\Cone{\mathrm{cone}}      % cone of a set of vectors
\def\Sph{\mathcal{S}_1}       % unit sphere
%%% Scalar product
\def\scal#1#2{\langle{#1}~|~{#2}\rangle}
\def\bigscal#1#2{\bigl\langle{#1}~\bigm|~{#2}\bigr\rangle}
\def\Bigscal#1#2{\Bigl\langle{#1}~\Bigm|~{#2}\Bigr\rangle}
\def\valscal#1#2{\left\langle{#1}~\middle|~{#2}\right\rangle}
%%% Syntax
\def\<{\langle}
\def\>{\rangle}
\def\Void{*} % void object
\def\Pair#1#2{(#1,#2)} % pairing construct
\def\Lam#1#2#3{\lambda#1^{#2}{.}#3} % lambda abstraction
%%% Let construct
\def\letkeyword{\mathsf{let}}
\def\inkeyword{\mathsf{in}}
\def\LetV#1#2{\letkeyword~\Void=#1~\inkeyword~#2}
\def\LetP#1#2#3#4#5#6{\letkeyword\,\Pair{#1^{#2}}{#3^{#4}}=#5~\inkeyword~#6}
%%% Left injection
\def\inleftkeyword{\texttt{inl}}
\def\Inl#1{\inleftkeyword(#1)}
\def\bigInl#1{\inleftkeyword\bigl(#1\bigr)}
\def\BigInl#1{\inleftkeyword\Bigl(#1\Bigr)}
%%% Right injection
\def\inrightkeyword{\texttt{inr}}
\def\Inr#1{\inrightkeyword(#1)}
\def\bigInr#1{\inrightkeyword\bigl(#1\bigr)}
\def\BigInr#1{\inrightkeyword\Bigl(#1\Bigr)}
%%% Match construct
\def\matchkeyword{\texttt{match}}
\def\Match#1#2#3#4#5{\matchkeyword~#1~%
  \{\Inl{#2}\mapsto#3~|~\Inr{#4}\mapsto#5\}}
\def\bigMatch#1#2#3#4#5{\matchkeyword~#1~%
  \bigl\{\Inl{#2}\mapsto#3~\bigm|~\Inr{#4}\mapsto#5\bigr\}}
\def\BigMatch#1#2#3#4#5{\matchkeyword~#1~%
  \Bigl\{\Inl{#2}\mapsto#3~\Bigm|~\Inr{#4}\mapsto#5\Bigr\}}
%%% Match construct for lists
\def\Nil{\ensuremath{\mathtt{nil}}}
\def\MatchL#1#2#3#4#5{\matchkeyword~#1~%
  \{\Nil\mapsto#2~|~#3::#4\mapsto#5\}}
\def\bigMatchL#1#2#3#4#5{\matchkeyword~#1~%
  \bigl\{\Nil\mapsto#2~\bigm|~#3::#4\mapsto#5\bigr\}}
\def\BigMatchL#1#2#3#4#5{\matchkeyword~#1~%
  \Bigl\{\Nil\mapsto#2~\Bigm|~#3::#4\mapsto#5\Bigr\}}
%%% If construct
\def\tt{\ensuremath{\mathtt{t\!t}}}
\def\ff{\ensuremath{\mathtt{f\!f}}}
\def\ifkeyword{\ensuremath{\mathtt{i{\mskip-1mu}f}}}
\def\If#1#2#3{\ifkeyword~#1~\{#2\mid#3\}}
\def\bigIf#1#2#3{\ifkeyword~#1~\bigl\{#2\bigm|#3\bigr\}}
\def\BigIf#1#2#3{\ifkeyword~#1~\Bigl\{#2\Bigm|#3\Bigr\}}
%%% Case construct
\def\case#1#2#3#4#5{\ensuremath{\mathsf{case}~#1~\mathsf{of}~\{#2\mapsto #4 \mid #3\mapsto #5\}}}
\def\gencase#1#2#3#4#5{\ensuremath{\mathsf{case}~#1~\mathsf{of}~\{#2\mapsto #4 \mid \dotsb \mid #3\mapsto #5\}}}
\def\gencasei#1#2#3{\ensuremath{\mathsf{case}~#1~\mathsf{of}~\{#2\mapsto #3\}}}
%%% Syntax (misc.)
\def\supp{\mathrm{supp}}
\def\weight{\varpi}
\def\Kron#1#2{\ensuremath{\delta_{#1,#2}}}
%%% Evaluation
\def\evalat{\mathrel{\triangleright}}
\def\nevalat{\mathrel{\not\triangleright}}
\def\lraneq{\rightsquigarrow}
\def\eval{\lra^*}
\def\lave{\mathrel{\reflectbox{$\eval$}}}
%%% Types
\def\Unit{\mathbb{U}}
\def\Bool{\mathbb{B}}
\def\arr{\rightarrow}
\def\Arr{\Rightarrow}
\def\Type{\mathbb{T}}
\def\BasisType{\Type_\flat}
%%% Semantics
\def\sem#1{\llbracket#1\rrbracket}
\def\semr#1{\{{\real}~#1\}}
\def\SUB#1#2{#1\le#2}
\def\NSUB#1#2{#1\not\le#2}
\def\EQV#1#2{#1\simeq#2}
\def\TYP#1#2#3{#1~{\vdash}~#2~{:}~#3}
\def\SORTH#1#2#3#4{#1~{\vdash}~#2\perp#3~{:}~#4}
\def\ORTH#1#2#3#4#5#6{#1~{\vdash}~(#2~{\vdash}~#3)\perp(#4~{\vdash}~#5)~{:}~#6}
\def\rnam#1{\textsc{\small\upshape(#1)}}
\def\snam#1{\textsc{\scriptsize\upshape(#1)}}
\def\real{\Vdash}
\def\ureal{\Vvdash}
%%% Misc.
\def\ds{\displaystyle}
\outer\long\def\COUIC#1{}
\def\sqrthalf{{\textstyle\frac{1}{\sqrt{2}}}}
\def\minsqrthalf{{\textstyle\bigl(-\frac{1}{\sqrt{2}}\bigr)}}
\def\isqrthalf{{\textstyle\frac{i}{\sqrt{2}}}}
\def\minisqrthalf{{\textstyle\bigl(-\frac{i}{\sqrt{2}}\bigr)}}

%%% Even more macros -- to be merged

\newcommand\pair[1]{\langle #1 \rangle}
\newcommand\Let[3]{\mathsf{let}\ {#1}={#2}\ \mathsf{in}\ {#3}}
\newcommand{\ttrue}{\ensuremath{\mathtt{t\!t}}}
\newcommand{\ffalse}{\ensuremath{\mathtt{f\!f}}}
\newcommand{\tif}[3]{\mathsf{if}\left(#1\right)\left(#2\right)\left(#3\right)}
\newcommand{\pif}[2]{\ensuremath{\mathtt{if}\left(#1\right)\left(#2\right)}}
\newcommand\trad[1]{\llparenthesis{#1}\rrparenthesis}
\newcommand\B{\mathbb Z}
\newcommand\XB{\mathbb X}
\newcommand\Hd{\mathsf{H}}
\newcommand\Q{\sharp\B}
\newcommand{\True}{\mathbb{T}}
\newcommand{\False}{\mathbb{F}}
\newcommand{\Cx}{\mathbb{C}}
\newcommand{\cnot}[2]{\mathsf{CNOT}\ #1\ #2}
\newcommand{\tcnot}{\mathsf{CNOT}}
\newcommand{\pauliX}[1]{\mathsf{NOT}\ #1}
\newcommand{\pauliZXB}{\mathsf{Z}_{\XB}}
\newcommand{\cnotXB}[2]{\mathsf{CNOT}_{\XB}\ #1\ #2}
\newcommand{\tcnotXB}{\mathsf{CNOT}_{\XB}}
\newcommand{\pauliXXB}[1]{\mathsf{NOT}_{\XB}\ #1}
\newcommand{\tpauliXXB}{\mathsf{NOT}_{\XB}}
\newcommand{\Bell}{\mathsf{Bell}}
\newcommand{\lambdaB}{\lambda_B}
%%% If construct
%\newcommand\basis[1]{\ensuremath{\mathsf{B}_{#1}}}
\newcommand\basis[1]{\ensuremath{\flat_{#1}}}
\newcommand\genbasis[3]{\ensuremath{\basis{\{#1\}_{#2}^{#3}}}}

%%% Teleportation algorithm

\newcommand{\bellcase}[5]{\ensuremath{\mathsf{case}~#1~\mathsf{of}~ 
\{\Phi^+ \mapsto \Pair{\Phi^+}{#2} \mid
\Phi^- \mapsto \Pair{\Phi^-}{#3} \mid
\Psi^+ \mapsto \Pair{\Psi^+}{#4} \mid
\Psi^- \mapsto \Pair{\Psi^-}{#5} \}}}


\begin{document}

\setlength{\abovedisplayskip}{3pt}
\setlength{\belowdisplayskip}{3pt}
\setlength{\abovedisplayshortskip}{3pt}
\setlength{\belowdisplayshortskip}{3pt}

\maketitle
\begin{abstract}
  We present $\lambdaB$, a quantum‐control $\lambda$‐calculus that refines
  previous basis-sensitive systems by allowing abstractions to be expressed with
  respect to arbitrary---possibly entangled---bases.  Each abstraction and
  $\mathsf{let}$-construct is annotated with a basis, and a new basis-dependent
  substitution governs the decomposition of value combinations.  These
  extensions preserve the expressive power of earlier calculi while enabling
  finer reasoning about programs under basis changes.  A realisability semantics
  connects the reduction system with the type system, yielding a direct
  characterisation of unitary operators and ensuring safety by construction.
  From this semantics we derive a validated family of typing rules, forming the
  foundation of a type-safe quantum programming language.  We illustrate the
  expressive benefits of $\lambdaB$ through examples such as Deutsch’s algorithm
  and quantum teleportation, where basis-aware typing captures classical
  determinism and deferred-measurement behaviour within a uniform framework.
\end{abstract}
\newpage

\section{Introduction}
The no-cloning theorem \cite{WoottersZurek1982} and the no-deleting theorem
\cite{PatiBraunstein2000} are two well-known results in quantum mechanics that
state it is impossible to copy or delete an arbitrary qubit. 
There is, however, a subtlety: although arbitrary qubits
cannot be copied or deleted, this is possible for known---and, in the case of
deletion, separable---qubits. This implies that qubits with known values behave
as classical data and can be treated accordingly. Moreover, it suffices to know
the basis to which a qubit belongs in order to copy and delete it.

In most quantum programming languages, qubits are defined with respect to a
canonical basis—often referred to as the computational basis. In this setting,
classical bits correspond to the basis vectors, whereas qubits are unit-norm
linear combinations of them. Classical bits can be copied and deleted freely,
while such operations on arbitrary qubits are restricted.

In this paper, we introduce a quantum $\lambda$-calculus in the
quantum-control paradigm—by contrast with the classical-control one, which
cannot be superposed. Our approach is inspired by a line of work on basis-sensitive
quantum typing. The earliest system, Lambda-S \cite{DiazcaroDowekRinaldiBIO19},
identified qubits in the computational basis, allowing duplication and erasure
only in that case. Later, Lambda-S$_1$ \cite{DiazcaroGuillermoMiquelValironLICS19,DiazCaroMalherbe2022}
refined this approach by restricting to unit-norm vectors, introducing higher-order
abstractions, and ensuring that terms of type qubit-to-qubit denote unitary maps.
More recently, the Lambda-SX calculus \cite{DiazcaroMonzonAPLAS25} generalised
Lambda-S to arbitrary non-entangled single-qubit bases, to first order. The
$\lambda_B$ calculus presented here extends these ideas to a unit-norm, higher-order
setting over multiple-qubit bases, grounded in a realisability semantics; a comparison
with previous systems is provided in Table~\ref{tab:related}.
\begin{table}
  \centering
  \begin{tabular}{|l|c|c|c|c|c|}\hline
    & Lambda-S \cite{DiazcaroDowekRinaldiBIO19} 
    & Lambda-SX \cite{DiazcaroMonzonAPLAS25} 
    & Lambda-S$_1$ \cite{DiazcaroGuillermoMiquelValironLICS19,DiazCaroMalherbe2022} 
    & $\lambdaB$ 
    \\\hline\hline
    Canonical basis 
    & \cmark 
    & \cmark 
    & \cmark 
    & \cmark 
    \\\hline
    Unit-norm 
    &  
    &  
    & \cmark 
    & \cmark 
    \\\hline
    Realisability semantics 
    &  
    &  
    & \cmark 
    & \cmark 
    \\\hline
    Higher-order 
    &  
    &  
    & \cmark (only in the model)
    & \cmark 
    \\\hline
    Multi-qubit bases 
    &  
    & \cmark 
    &  
    & \cmark 
    \\\hline
  \end{tabular}
  \caption{Comparison of basis-sensitive quantum $\lambda$-calculi.}
  \label{tab:related}
\end{table}

The realisability methodology provides a constructive framework connecting
operational semantics with type systems. In our context, it allows extracting a
sound type system directly from the reduction semantics, ensuring safety by
construction. The approach is as follows: define a calculus with a
deterministic evaluation; then, define types as sets of closed values; finally, 
define typing so that an assertion equates to the term reducing to a value of
that type.  Rather than
building ad hoc rules, the system is derived from the calculus's computational
content. While the untyped calculus allows for non-physical terms, the
extracted type system acts as a logical filter: only unit-norm, unitary terms
are typable, enabling static certification of unitary behaviour across basis
changes and capturing basis-dependent invariants beyond fixed-basis systems.

Following this approach, we enrich abstractions with explicit basis
decorations. Intuitively, the reduction system treats values expressed in the
chosen basis as classical data, while linear combinations of these values
represent quantum data and reduce linearly within the term. This refinement
enables duplication and erasure for qubits in known bases while maintaining
linear handling for unknown qubits.

This work pursues two goals: first, to provide a more exact program description
via the extracted type system; and second, to express quantum algorithms compositionally,
rather than as mere circuit translations.

%\paragraph{Related work.}
The idea of tracking basis information within quantum programming languages has
appeared in several forms across the literature.  
The work in~\cite{Perdrix2008} proposed an abstract model for static
entanglement analysis, where the basis of qubits is used to control duplication
and track non-entanglement properties.  
The approach in~\cite{DiazcaroMonzonAPLAS25} introduced a typed
$\lambda$-calculus that incorporates basis annotations into the type system,
ensuring linear handling of quantum data while providing strong
meta-theoretical guarantees.  
Our system generalises these ideas by allowing abstractions to range over
arbitrary---possibly entangled---bases, extending basis awareness beyond the
single-qubit setting and recovering higher-order computation.

A number of other frameworks support reasoning about multiple bases through
different paradigms. The ZX-calculus~\cite{CoeckeRossICALP08} provides a
graphical language where spiders correspond to computations relative to the
computational and Hadamard bases, while the Many-Worlds Calculus
\cite{ChardonnetdeVismeValironVilmartLMCS25} also accommodates
superpositions of programs but within a diagrammatic semantics.  
By contrast, our approach operates directly within a typed term calculus,
allowing basis transitions to be tracked syntactically.

From a categorical standpoint, the framework
in~\cite{HeunenKaarsgaard2021} models measurement and decoherence through
\emph{quantum information effects}, capturing basis change as a semantic
effect.  
Similarly,~\cite{CaretteJeunenKaarsgaardSabry2024} introduces
\textsc{Quantum}$\Pi$, a language combining two interpretations of a reversible
classical calculus—one per basis—via an effect construction.  
In both cases, basis sensitivity arises semantically.  
In contrast, $\lambdaB$ integrates it directly at the syntactic and typing
levels, ensuring that basis transformations are explicit in program structure.

Other languages aim to unify classical and quantum computation.
For instance,~\cite{VoichickLiRandHicks2023} generalises classical
control constructs and interprets duplication and discarding semantically as
entanglement and partial trace.  
Our calculus follows the opposite philosophy: linearity and duplicability are
syntactically restricted by the basis-sensitive type system, guaranteeing
coherent quantum control without relying on external semantic constraints.

Finally, related efforts in semantic characterisation—such as the fully
abstract model of~\cite{ClairambaultdeVisme2019} or the dual calculi
described in~\cite{ChoudhuryGay2025}—approach the logical foundations of
quantum computation from complementary angles.  
In contrast, $\lambdaB$ focuses on the syntactic distinction of bases and
superpositions within a unified quantum $\lambda$-calculus, providing a concrete
basis-sensitive typing and higher-order quantum control.

%\paragraph{Contributions.}
The contributions of this paper can be summarised as follows.
\begin{itemize}[nosep]
  \item We introduce $\lambdaB$, a quantum-control $\lambda$-calculus that 
  generalises the basis-sensitive paradigm to arbitrary, possibly entangled 
  multi-qubit bases in a higher-order setting.
  \item We provide a full \emph{realisability-based semantics} that 
  logically internalises the semantic requirement of unitarity. This allows 
  the systematic extraction of a sound type system where the physical consistency 
  of programs (unit-norm preservation and isometries) is guaranteed by 
  construction.
  \item We achieve a direct \emph{semantic characterisation of unitary operators} 
  of arbitrary finite dimension $2^n$ as abstractions in our type algebra. 
  This result generalises previous characterisations which were restricted 
  to the computational basis.
  \item We demonstrate the \emph{expressive benefits} of basis-awareness 
  through a series of examples. Specifically, we show how the system captures 
  classical determinism and the deferred-measurement principle via static 
  typing, yielding more precise program descriptions than earlier calculi.
\end{itemize}

The paper is organised as follows.  
Section~\ref{sec:calculus} introduces the core language, its syntax, congruence rules,
and basis-dependent substitution.  
Section~\ref{sec:reduction} defines the operational semantics, and
Section~\ref{sec:model} presents the realisability model with unitary type semantics,
the characterisation of unitary operators, and the typing rules.  
Section~\ref{sec:examples} presents some representative examples.
Section~\ref{sec:conclusion} concludes.
All omitted proofs appear in the appendices.

\section{Core language}\label{sec:calculus}
%\subsection{Syntax and congruence}
This section presents the calculus on which our realisability model will be
built. It is a $\lambda$-calculus extended with linear combinations of terms, 
along the lines of~\cite{ArrighiDowekLMCS17}.

%\begin{table}[t]
%  \begin{align*}
%    v &::= x \mid \Lam{x}B{\vec{t}} \mid (v, v) \mid \ket{0} \mid \ket{1} &
%    (\Val)\\
%    t &::= v \mid tt \mid (t,t) \mid
%    \LetP{x}{B_1}{y}{B_2}{\vec{t}}{\vec{t}}\mid
%    \gencase{\vec{t}}{\vec{v}}{\vec{v}}{\vec{t}}{\vec{t}} &
%    (\Lambda) \\
%    \vec{v} &::= v \mid \vec{v}+\vec{v} \mid \alpha \vec{v} \mid \vec 0
%    \qquad\hfill(\alpha\in\C) & (\ValD) \\
%    \vec{t} &::= t \mid \vec{t}+\vec{t} \mid \alpha \vec{t} \mid \vec 0
%    \qquad\hfill(\alpha\in\C) & (\vec \Lambda)
%  \end{align*}
%  \caption{Syntax of the calculus, where $B, B_1, B_2 \subseteq \ValD$.}
%  \label{tab:Syntax}
%\end{table}
\begin{table}[t]
  \centering
  \begin{tabular}{rcll}
    $v$ & $::=$ & $x \mid \Lam{x}B{\vec{t}} \mid (v, v) \mid \ket{0} \mid \ket{1}$ & $(\Val)$ \\
    $t$ & $::=$ & $v \mid tt \mid (t,t) \mid \LetP{x}{B_1}{y}{B_2}{\vec{t}}{\vec{t}} \mid \gencase{\vec{t}}{\vec{v}}{\vec{v}}{\vec{t}}{\vec{t}}$ & $(\Lambda)$ \\
    $\vec{v}$ & $::=$ & $v \mid \vec{v}+\vec{v} \mid \alpha \vec{v} \mid \vec 0 \qquad (\alpha\in\C)$ & $(\ValD)$ \\
    $\vec{t}$ & $::=$ & $t \mid \vec{t}+\vec{t} \mid \alpha \vec{t} \mid \vec 0 \qquad (\alpha\in\C)$ & $(\vec \Lambda)$
  \end{tabular}
  \caption{Syntax of the calculus, where $B, B_1, B_2 \subseteq \ValD$.}
  \label{tab:Syntax}
\end{table}
The syntax of the calculus is described in
Table~\ref{tab:Syntax}; it is divided into four
distinct syntactic categories: \emph{basic values} ($\Val$), \emph{basic terms}
($\Lambda$), \emph{value combinations} ($\ValD$), and \emph{term
combinations} ($\vec\Lambda$). Values are composed of variables, decorated
lambda abstractions, and two basis values representing orthogonal vectors,
$\ket 0$ and $\ket 1$. A pair of values is also considered a value.  Terms
include values, applications, pair constructors and destructors, and pattern
matching for orthogonal vectors, represented by the $\mathsf{case}$ operator.
Both term and value combinations are built $\C$-linear combinations of
basic terms and values, respectively. We use $v,u,w$ to denote values and
$t,s,r$ for terms, writing $\vec{\cdot}$ when they are combinations.

The symbols $B$, $B_1$, and $B_2$ appearing in the abstractions and pair
destructors denote the bases of the vector spaces in which the terms are
expressed; their precise nature is made explicit below.
Before formally defining the notion of bases, we first establish the
algebraic structure underlying the space of value combinations.
This is achieved by introducing the congruence relation defined in
Table~\ref{tab:Congruence}.
The congruence captures the intended behaviour of addition and scalar
multiplication, allowing us to treat linear combinations of terms and values
as genuine algebraic entities rather than mere syntactic constructions.
In particular, it enforces the commutativity and associativity of addition and
the distributivity of scalar multiplication, thereby justifying summation
notation~$\sum$.
%\begin{table}[t]
%  \begin{align*}
%    0\vec t &\equiv \vec 0 & \vec t+\vec 0 &\equiv \vec t \\
%    1\vec{t} &\equiv \vec{t} &
%    \alpha(\beta\vec{t}) &\equiv (\alpha\beta)\vec{t} \\
%    \vec{t}_{1}+\vec{t}_{2} &\equiv \vec{t}_{2}+\vec{t}_{1} &
%    (\vec{t}_{1}+\vec{t}_{2})+\vec{t}_{3} &\equiv \vec{t}_{1}+(\vec{t}_{2}+\vec{t}_{3}) \\
%    (\alpha+\beta)\vec{t} &\equiv \alpha\vec{t}+\beta\vec{t} &
%    \alpha(\vec{t}_{1}+\vec{t}_{2}) &\equiv \alpha\vec{t}_{1}+\alpha\vec{t}_{2} \\
%    \vec{t}(\alpha\vec{s}) &\equiv \alpha(\vec{t}\vec{s}) &
%    (\alpha\vec{t})\vec{s} &\equiv \alpha(\vec{t}\vec{s}) \\
%    (\vec{t}+\vec{s})\vec{r} &\equiv \vec{t}\vec{r}+\vec{s}\vec{r} &
%    \vec{t}(\vec{s}+\vec{r}) &\equiv \vec{t}\vec{s}+\vec{t}\vec{r} &&
%  \end{align*}
%  \vspace{-2\baselineskip}
%  \begin{align*}
%    &\LetP{x_1}{X}{x_2}{Y}{(\alpha\vec{t})}{\vec{s}}
%    \equiv \alpha(\LetP{x_1}{X}{x_2}{Y}{\vec{t}}{\vec{s}}) \\
%    &\LetP{x_1}{X}{x_2}{Y}{\vec{t}+\vec{s}}{\vec{r}}
%    \equiv
%      (\LetP{x_1}{X}{x_2}{Y}{\vec{t}}{\vec{r}})
%      +(\LetP{x_1}{X}{x_2}{Y}{\vec{s}}{\vec{r}}) 
%      \\
%    &\gencase{\alpha \vec{t}}{\vec{v}_{1}}{\vec{v}_{k}}{\vec{s}_{1}}{\vec{s}_{k}}
%    \equiv \alpha(\gencase{\vec{t}}{\vec{v}_{1}}{\vec{v}_{k}}{\vec{s}_{1}}{\vec{s}_{k}}) \\
%    &\gencase{(\vec{t}+\vec{s})}{\vec{v}_{1}}{\vec{v}_{k}}{\vec{r}_{1}}{\vec{r}_{k}}
%    \equiv\begin{aligned}[t]
%      &\gencase{\vec{t}}{\vec{v}}{\vec{w}}{\vec{r}_{1}}{\vec{r}_{k}}\\
%      &+\gencase{\vec{s}}{\vec{v}}{\vec{w}}{\vec{r}_{1}}{\vec{r}_{k}}
%    \end{aligned}
%  \end{align*}
%  \caption{Term congruence}
%  \label{tab:Congruence}
%\end{table}
\begin{table}[t]
  \centering
  \begin{tabular}{r@{~}c@{~}l@{\qquad}r@{~}c@{~}l}
    \multicolumn{6}{c}{$
      \vec{t}_1+\vec{t}_2 \equiv \vec{t}_2+\vec{t}_1,
      \quad
      (\vec{t}_1+\vec{t}_2)+\vec{t}_3 \equiv \vec{t}_1+(\vec{t}_2+\vec{t}_3),
      \quad
      0\vec{t} \equiv \vec{0},
      \quad
      1\vec{t} \equiv \vec{t},
      \quad 
      \vec{t}+\vec{0} \equiv \vec{t},
    $} \\
    \multicolumn{6}{c}{$
      \alpha(\beta\vec{t})\equiv(\alpha\beta)\vec t,
      \quad 
      (\alpha+\beta)\vec{t} \equiv \alpha\vec{t}+\beta\vec{t},
      \quad
      \alpha(\vec{t}_1+\vec{t}_2) \equiv \alpha\vec{t}_1+\alpha\vec{t}_2,
    $} \\
    \multicolumn{6}{c}{$
      \vec{t}(\alpha\vec{s}) \equiv \alpha(\vec{t}\vec{s}),
      \qquad
      (\alpha\vec{t})\vec{s} \equiv \alpha(\vec{t}\vec{s}),
      \qquad
      (\vec{t}+\vec{s})\vec{r} \equiv \vec{t}\vec{r}+\vec{s}\vec{r},
      \qquad
      \vec{t}(\vec{s}+\vec{r}) \equiv \vec{t}\vec{s}+\vec{t}\vec{r},
    $} \\
    \multicolumn{3}{l}{$\mathsf{let} (x^X, y^Y) = (\alpha\vec{t}) \text{ in } \vec{s} \equiv \alpha\mathsf{let} (x^X, y^Y) = \vec{t} \text{ in } \vec{s}$} &
    \multicolumn{3}{l}{$\text{case } \alpha \vec{t} \text{ of } \{\vec{v}_i \mapsto \vec{s}_i\} \equiv \alpha\text{case } \vec{t} \text{ of } \{\vec{v}_i \mapsto \vec{s}_i\}$} \\
    \noalign{\smallskip}
    \multicolumn{3}{l}{$\mathsf{let} (x^X, y^Y) = (\vec{t}+\vec{s}) \text{ in } \vec{r} \equiv \sum_{z \in \{\vec{t}, \vec{s}\}} \mathsf{let} \dots z \dots$} &
    \multicolumn{3}{l}{$\text{case } (\vec{t}+\vec{s}) \text{ of } \dots \equiv \sum_{z \in \{\vec{t}, \vec{s}\}} \text{case } z \dots$} 
  \end{tabular}
  \caption{Term congruence}
  \label{tab:Congruence}
\end{table}

The rules in Table~\ref{tab:Congruence} ensure that the set of value combinations
satisfies the axioms of a vector space. The notion of basis in this calculus builds
on this algebraic foundation. Once the vector structure is established, we can define
which subsets $B\subseteq \ValD$ qualify as \emph{bases}, thereby justifying the
decorations appearing in the syntax of Table~\ref{tab:Syntax}.

We will also consider the following notation for linear combinations of
pairs. This notation does not appear in the syntax, but is rather useful to describe
specific states.
\[
  \Pair{\alpha  t+\vec{t}_{1}}{\vec{t}_{2}} 
  := \alpha\Pair{t}{\vec{t}_{2}} + \Pair{\vec{t}_{1}}{\vec{t}_{2}}
  \qquad\qquad
  \Pair{r}{\alpha  t+\vec{t}_{1}} 
  := \alpha\Pair{r}{t} + \Pair{r}{\vec{t}_{1}}
\]

Before proceeding further, let us briefly illustrate the intuition behind the
congruence in Table~\ref{tab:Congruence}.  
The key idea is that arguments are decomposed over the bases associated with
their corresponding abstractions.  
As in linear algebra, a vector can be rewritten as a linear combination of
basis elements; for instance,
\(
  (1,0)
  = \tfrac{1}{\sqrt{2}}\!\left(\tfrac{(1,1)}{\sqrt{2}}
  + \tfrac{(1,-1)}{\sqrt{2}}\right)
\)
expresses $(1,0)$ in the basis
$\{\tfrac{(1,1)}{\sqrt{2}}, \tfrac{(1,-1)}{\sqrt{2}}\}$.
By identifying $(1,0)$ with $\ket{0}$ and $(0,1)$ with $\ket{1}$, we see that
the calculus allows every vector (or value combination) to be expressed as a
superposition of elements of the basis attached to each abstraction.
The congruence ensures that these linear combinations behave algebraically as
in a complex vector space, supporting the standard operations of addition,
scalar multiplication, and the zero vector.

To properly characterise the sets that decorate the lambda abstractions, we
must first specify the kind of values they contain.
\begin{definition}[Qubits]\label{def:Qubit}
  A \emph{1-qubit state} is a value combination of the form
  $\alpha\ket{0} + \beta\ket{1}$, where $|\alpha|^2 + |\beta|^2 = 1$. An
  \emph{$n$-qubit state} is a value combination of the form
  $\alpha\Pair{\ket{0}}{\vec{w}_{1}} + \beta\Pair{\ket{1}}{\vec{w}_{2}}$, where
  $\vec{w}_{1}$ and $\vec{w}_{2}$ are $(n-1)$-qubit states,
  where $|\alpha|^2 + |\beta|^2 = 1$.
\end{definition}
We use the usual Dirac shorthand $\ket{xy}$ for $(\ket{x},\ket{y})$, and
extend it to longer tuples by associating to the right.

From now on, we call the value combinations, i.e. the elements of $\ValD$,
\emph{vectors}. The vector space $\ValD$ is equipped with an inner
product defined by,  
\(
  \scal{\vec{v}}{\vec{w}} := \sum_{i=1}^k\sum_{j=1}^k
  \overline{\alpha_i}\,\beta_j\,\delta_{v_i,w_j},
\)
and an $\ell_2$-norm 
\(
  \|\vec{v}\| := \sqrt{\scal{\vec{v}}{\vec{v}}}
  = \sqrt{\sum_{i=1}^k|\alpha_i|^2}
\),
where $\vec{v}=\sum_{i=1}^k\alpha_i v_i$ and
$\vec{w}=\sum_{j=1}^k\beta_j w_j$, and $\delta_{v_i,w_j}$ is the Kronecker
delta, equal to $1$ if $v_i=w_j$ and $0$ otherwise.

With this notion of inner product, we can complete the description of the
calculus syntax. As expected, two values are \emph{orthogonal} when their inner
product is equal to zero. We can now formally define the sets that
decorate abstractions.
\begin{definition}[Basis]\label{def:NthDimensionalBasis}
  A set of value combinations $B$ is a \emph{$2^n$-dimensional orthonormal
  basis} if it satisfies:
  \begin{enumerate*}
    \item $B$ has cardinality $2^n$.
    \item Each element of $B$ is an $n$-qubit state \emph{(cf.~Definition~\ref{def:Qubit})}.
    \item Distinct elements of $B$ are pairwise orthogonal.
  \end{enumerate*}
\end{definition}

As a first step, the syntax introduced in Table~\ref{tab:Syntax}
can be restricted to sets $B$ forming $2^n$-dimensional
orthonormal bases. These sets indicate the basis in which a term is expressed:
qubits in the decorating basis are treated in call-by-value, while others are decomposed as
$\C$-linear combinations of basis elements, with the function applied linearly
to each component.  If a term cannot be expressed in the decorating basis,
evaluation becomes stuck.

The basis analysis will be restricted to $n$-qubit states. In other words, we will
not be considering bases for $\lambda$-abstractions, this is due to the fact that a syntactical
approach to higher-order bases is insufficient. For example, two different $\lambda$-terms
which encode the same function $f$, are considered orthogonal to each other. This itself is
an issue since the norm-$1$ linear combination of these terms does not preserve unitarity.
Solving this problem is not a trivial endeavor and we leave it as a possible avenue for future
research.

As expected, no non-trivial linear combination of basis elements yields the null vector;
otherwise, some basis vector would violate pairwise orthogonality. Consequently, each
decomposition over a basis is unique.

\begin{theorem}[Unique decomposition]\label{thm:UniqueDecomposition}
  If $B$ is a $2^n$-dimensional basis, then every $n$-qubit state has a
  unique decomposition over $B$.
  \qed
\end{theorem}

\begin{corollary}[Preservation under congruence]\label{cor:EquivalentDecomposition}
  If $\vec{v} \equiv \vec{w}$, then they share the same decomposition over any
  basis $B$.
  \qed
\end{corollary}

We denote the computational basis by $\B=\{\ket{0},\ket{1}\}$ and the Hadamard
basis by $\XB=\{\ket{+},\ket{-}\}$, where
$\ket{\pm}=\tfrac{1}{\sqrt{2}}(\ket{0}\pm\ket{1})$.  

%\subsection{Basis-dependent substitutions}
As usual, the expression $\vec t[\vec v/x]$ denotes the
capture-avoiding substitution of $\vec v$ for $x$ in $\vec t$.
However, beta-reduction depends on the basis chosen for the abstraction, so we must
define a substitution that takes this mechanism into account. Intuitively, this
operation substitutes variables with vectors expressed in the chosen basis; the
accompanying coefficients are those of the value combination being
substituted.

Alongside this substitution, we expand the set of orthogonal basis with a 
special symbol, denoted $\AbsBasis$. Since we only consider bases for
first-order terms, we need a method to deal with combinations of $\lambda$-abstractions.
This symbol marks that the argument of a function should be treated linearly over its
basic values as defined in the syntax. This is the only instance of higher-order
functions in the calculus. So, from now on, the symbols $B, B_1, B_2$ in Table
\ref{tab:Syntax} are either $2^n$-dimensional orthonormal basis, or the special
symbol $\AbsBasis$.

%\begin{definition}[Basis-dependent substitution]
%  Let $\vec t$ be a term combination, $\vec v$ a value combination, $x$ a
%  variable, and $B$ an orthonormal basis. We define the substitution
%  $\vec t\ansubst{\vec v/x}{B}$ as follows:
%  \[
%    \vec t\ansubst{\vec v/x}{B} \;=\;
%    \begin{cases}
%      \ds\sum_{i\in I}\alpha_i\,\vec t\,[\vec{b}_{i}/x] &
%        \text{if } B=\{\vec{b}_{i}\}_{i\in I}\ \text{and}\
%        \vec v \equiv \ds\sum_{i\in I}\alpha_i\,\vec{b}_{i},\\[6pt]
%      \ds\sum_{i\in I}\alpha_i\,\vec t\,[v_i/x] &
%        \text{if } B=\AbsBasis\ \text{and}\
%        \vec v = \ds\sum_{i\in I}\alpha_i\,v_i,\\[2pt]
%      \text{undefined} & \text{otherwise.}
%    \end{cases}
%  \]
%  Notice that the uniqueness of the decomposition over $B$ is guaranteed by Theorem~\ref{thm:UniqueDecomposition}.
%\end{definition}
\begin{definition}[Basis-dependent substitution]
  Let $\vec t$ be a term combination, $\vec v$ a value combination, $x$ a
  variable, and $B$ an orthonormal basis or $\AbsBasis$. The substitution 
  $\vec t\ansubst{\vec v/x}{B}$ is defined as $\sum_{i\in I}\alpha_i\vec{t}[\vec{b}_i/x]$ 
  if $B=\{\vec{b}_i\}_{i\in I}$ and $\vec{v}\equiv\sum_{i\in I} \alpha_i\vec{b}_i$; 
  or as $\sum_{i\in I}\alpha_i\vec{t}[v_i/x]$ if $B=\AbsBasis$ and 
  $\vec{v}=\sum_{i\in I} \alpha_i v_i$; otherwise it is undefined. The 
  uniqueness of the decomposition over $B$ is guaranteed by 
  Theorem~\ref{thm:UniqueDecomposition}.
\end{definition}

This definition captures two substitution modes: linear decomposition of $\vec v$ over an explicit basis $B$, or over its basic values when $B=\AbsBasis$ (recovering the mechanism of~\cite{ArrighiDowekLMCS17}). The latter is the only case applicable to $\lambda$-abstractions. The operation extends to pairs: if $\vec v\equiv\sum_{i\in I}\alpha_i(\vec{v}_i, \vec{w}_i)$, then $\vec t\ansubst{\vec v/x\otimes y}{B_1\otimes B_2} = \sum_{i\in I}\alpha_i (\vec t\ansubst{\vec{v}_i/x}{B_1}\ansubst{\vec{w}_i/y}{B_2})$, where $B_1, B_2$ are the bases associated with $x, y$.

%The definition captures two distinct substitution modes.  
%In the first, substitution proceeds linearly using the decomposition of
%$\vec v$ over the explicit basis $B$. In the second, when $B=\AbsBasis$,
%substitution proceeds linearly over the basic values that form $\vec v$.  
%This case recovers the mechanism of non base-sensitive calculi
%---as first defined in \cite{ArrighiDowekLMCS17}. This special case is
%also the only one applicable to $\lambda$-abstractions, since abstractions
%do not belong to orthonormal bases.
%
%The definition also extends to pairs of values. If
%$\vec v\equiv\sum_{i\in I}\alpha_i\,\Pair{\vec{v}_{i}}{\vec{w}_{i}}$,
%then:
%\[
%  \vec t\ansubst{\vec v/x\otimes y}{B_1\otimes B_2}
%    \;=\; \sum_{i\in I}\alpha_i\,
%    \bigl(\vec t\ansubst{\vec{v}_{i}/x}{B_1}\ansubst{\vec{w}_{i}/y}{B_2}\bigr),
%    \text{ if defined},
%\]
%where $B_1$ and $B_2$ are (orthonormal) bases---or $\AbsBasis$---associated
%with $x$ and $y$, respectively; the symbol $\otimes$ in $B_1\otimes B_2$ is
%purely notational.

%\begin{example}
%  Let
%  $\vec v = \alpha\ket{01}
%           + \beta\ket{10}$.
%  Then 
%  \[
%    (y,x)\ansubst{\vec v/x\otimes y}{\B\otimes\B}
%    = \alpha\,(y,x)[\ket{0}/x][\ket{1}/y]
%    + \beta\,(y,x)[\ket{1}/x][\ket{0}/y] 
%    = 
%    \alpha\ket{10}
%    + \beta\ket{01}
%  \]
%\end{example}

With this substitution in place, we can establish certain needed properties.
First, basis-dependent substitution distributes over linear combinations.

\begin{lemma}[Distributivity over linear combinations]\label{lem:distributiveSubstitution}
  For term combinations $\vec{t}_{i}$, a value combination $\vec{v}$, a
  variable $x$, coefficients $\alpha_i\in\C$, and a basis $B$ such that
  $\ansubst{\vec v/x}{B}$ is defined:
  \(
    \Bigl(\sum_i \alpha_i\vec{t}_{i}\Bigr)\ansubst{\vec v/x}{B}
    \equiv
    \sum_i \alpha_i\vec{t}_{i} \ansubst{\vec v/x}{B}
  \).
  \qed
\end{lemma}

The next result states that substitution behaves consistently within each
equivalence class induced by the congruence $\equiv$.

\begin{lemma}[Compatibility with congruence]\label{lem:EquivSubstitutions}
  For value combinations $\vec{v},\vec{w}$, a term combination $\vec{t}$, and
  an orthonormal basis $B$ such that both
  $\ansubst{\vec{v}/x}{B}$ and $\ansubst{\vec{w}/x}{B}$ are defined:
  if $\vec{v}\equiv\vec{w}$, then
  $\vec{t}\ansubst{\vec{v}/x}{B}
  =\vec{t}\ansubst{\vec{w}/x}{B}$.
  \qed
\end{lemma}

We remark that the result in the previous lemma uses $=$ instead of $\equiv$
since $\vec{t}\ansubst{\vec{v}/x}{B}$ and $\vec{t}\ansubst{\vec{w}/x}{B}$
are both substitution with the exact same outcome. Note that this property
does not extend across different bases; that is, $\vec{t}\ansubst{\vec{v}/x}{A}
\not\equiv\vec{t}\ansubst{\vec{v}/x}{B}$.
For example,
\[
  (\Lam{x}{C}{y})\ansubst{\ket{+}/y}{\XB}
  = \Lam{x}{C}{\ket{+}} 
  \not\equiv
  \tfrac{1}{\sqrt{2}}\big((\Lam{x}{C}{\ket{0}})
  + (\Lam{x}{C}{\ket{1}})\big)
  = (\Lam{x}{C}{y})\ansubst{\ket{+}/y}{\B}.
\]
This difference arises because the relation $\equiv$ does not commute with
lambda abstraction, nor with the case construct (indeed, the calculus is
intensional with respect to functions, since it is issued from a realizability
model).  Although the two terms are operationally equivalent, the calculus
distinguishes between the superposition of results, $\Lam{x}{B}{\alpha\vec{v}_{1}
+ \beta\vec{v}_{2}}$, and the superposition of functions,
$\alpha(\Lam{x}{B}{\vec{v}_{1}}) + \beta(\Lam{x}{B}{\vec{v}_{2}})$.  This distinction
reflects a physical intuition: the former corresponds to a single experiment
producing a superposition of outcomes, while the latter represents a
superposition of distinct experiments.

Finally, we introduce a convenient notation for generalised substitutions over
a term by closed values. A substitution $\sigma$ can be seen as a finite set of
individual substitutions applied consecutively to a term. Formally, for a term
$\vec{t}$, closed value combinations $\vec{v}_{1},\dots,\vec{v}_{n}$, variables
$x_1,\dots,x_n$, and bases $B_1,\dots,B_n$:
\(
  \vec{t}\ansubst{\sigma}{}
  := \vec{t}\ansubst{\vec{v}_{1}/x_1}{B_1}\dotsb\ansubst{\vec{v}_{n}/x_n}{B_n}.
\)
Since each $\vec{v}_{i}$ is closed, the order of substitutions is irrelevant. We
regard $\sigma$ as a partial function from variables to pairs of closed value
combinations and bases, and write $\dom{\sigma}$ for its domain. The operation
extends naturally: for a term $\vec{t}$, substitution $\sigma$, a new variable
$x\notin\dom{\sigma}$, value combination $\vec{v}$, and basis $B$,
\(
  \vec{t}\ansubst{\sigma}{}\ansubst{\vec{v}/x}{B}
  = \vec{t}\ansubst{\sigma'}{}
\),
where $\sigma'$ extends $\sigma$ by mapping $x$ to $(\vec v,B)$. Likewise, two
disjoint substitutions $\sigma_1$ and $\sigma_2$ can be merged:
\(
  \vec{t}\ansubst{\sigma_1}{}\ansubst{\sigma_2}{}
  = \vec{t}\ansubst{\sigma'}{}
\),
where $\dom{\sigma_1}\cap\dom{\sigma_2}=\emptyset$ and $\sigma'$ coincides with
$\sigma_i$ on $\dom{\sigma_i}$ for $i=1,2$.


\section{Operational semantics}\label{sec:reduction}
%\begin{table}[t]
%  \begin{align*}
%    \text{If }\vec{t}\ansubst{\vec v/x}{X}\text{ is defined}:\quad&
%    (\Lam{x}{X}{\vec{t}})\vec{v} \lraneq \vec{t}\ansubst{\vec v/x}{X}\\
%    \text{If }\vec{t}\ansubst{\vec v/x}{X\otimes Y}\text{ is defined}:\quad&
%    \LetP{x}{X}{y}{Y}{\vec v}{\vec{t}} \lraneq \vec{t}\ansubst{\vec{v}/x\otimes y}{X\otimes Y}\\
%    \text{If for all } i\neq j\ \vec{v}_{i}\perp\vec{v}_{j}:\quad&
%    \gencase{\sum_{i=1}^{n}\alpha_i \vec{v}_{i}}{\vec{v}_{1}}{\vec{v}_{n}}{\vec{t}_{1}}{\vec{t}_{n}} \lraneq \sum_{i=1}^{n}\alpha_i \vec{t}_{i}
%  \end{align*}
%  If $t\lraneq \vec r$, then
%  $\left\{\begin{aligned}
%    st &\lraneq s\vec r \qquad\qquad tv \lraneq \vec rv \\
%    \text{If }\alpha\neq 0:\quad\vec v+\alpha t+\vec s &\lraneq \vec v+\alpha \vec r+\vec s & (*)\\
%    \LetP{x}{A}{y}{B}{t}{\vec{s}} &\lraneq \LetP{x}{A}{y}{B}{\vec r}{\vec{s}} \\
%    \gencase{t}{\vec{v}_{1}}{\vec {v_n}}{\vec{s}_{1}}{\vec{s}_{n}} &\lraneq \gencase{\vec r}{\vec{v}_{1}}{\vec {v_n}}{\vec{s}_{1}}{\vec{s}_{n}}
%\end{aligned}\right.$
%
%   \noindent In $(*)$ we use the convention that the subterms in the sum are ordered as follows: first we place all the values, if any, followed by the non-value terms in a lexicographical order.
%  \caption{Reduction system.}
%  \label{tab:Reduction}
%\end{table}
\begin{table}[t]
  \centering
  \begin{tabular}{r@{~}c@{~}l}
    $(\lambda x^X.\vec{t})\vec{v}$ & $\lraneq$ & $\vec{t}\langle\vec{v}/x\rangle_X$ \hfill (if defined) \\
    $\mathsf{let} (x^X, y^Y) = \vec{v} \text{ in } \vec{t}$ & $\lraneq$ & $\vec{t}\langle\vec{v}/x \otimes y\rangle_{X \otimes Y}$ \hfill (if defined) \\
    %$\mathsf{case} \sum \alpha_i \vec{v}_i \text{ of } \{\vec{v}_1 \mapsto \vec{t}_1 \mid \dots\}$ & $\rightsquigarrow$ & $\sum \alpha_i \vec{t}_i$ \hfill (if $\vec{v}_i \perp \vec{v}_j$) \\
    $\gencasei{\sum_{i}\alpha_i \vec{v_i}}{\vec{v_i}}{\vec{t_i}}$ & $\lraneq$ & $\sum_{i}\alpha_i \vec{t_i}$ \hfill (if $\vec{v}_i \perp \vec{v}_j$) \\
    \multicolumn{3}{l}{
      If $t\lraneq \vec r$, then
      $\left\{\begin{aligned}
	  &st \lraneq s\vec r, \qquad\qquad tv, \lraneq \vec rv, \qquad\qquad
	  \vec v+\alpha t+\vec s \lraneq \vec v+\alpha \vec r+\vec s\quad  (\text{if }\alpha\neq 0)^*,\\
	  &\LetP{x}{A}{y}{B}{t}{\vec{s}} \lraneq \LetP{x}{A}{y}{B}{\vec r}{\vec{s}},\\
	  &\gencasei{t}{\vec{v_i}}{\vec{s_i}} \lraneq \gencasei{\vec r}{\vec{v_i}}{\vec{s_i}}
      \end{aligned}\right.$
    }
  \end{tabular}
  \noindent\footnotesize $^*$Subterms are ordered: values first, then non-values in lexicographical order.
  \caption{Reduction system.}
  \label{tab:Reduction}
\end{table}
The reduction system interprets every vector relative to the basis attached to
its abstraction, allowing a step only when the argument can be decomposed on
that basis.  
Table~\ref{tab:Reduction} defines the elementary relation~$\lraneq$, while terms are
considered modulo the congruence of Table~\ref{tab:Congruence}.  
Hence the effective reduction, written~$\lra$, is defined modulo~$\equiv$:
a step $\vec t\lra\vec r$ abbreviates $\vec t\equiv\vec t'\lraneq\vec r'\equiv\vec r$.
The side condition~$\alpha\neq0$ avoids vacuous steps such as
$0t\lraneq0r$, preventing spurious nondeterminism. On that note,
the condition on the contextual rule of the sum, imposes a strategy
preserving the determinism of~$\lraneq$.

The main reduction rules are $\beta$-reduction, $\mathsf{let}$ binding, and
$\mathsf{case}$ pattern matching.  
Both $\lambda$ and $\mathsf{let}$ bind variables decorated with an orthonormal
basis, indicating which vectors are treated as classical data.  Linear
combinations of these vectors are handled as quantum data, reducing linearly by
the congruence rules, so that
$t(\alpha\vec s+\beta\vec r)$ is equivalent to
$\alpha\,t\vec s+\beta\,t\vec r$.

The only exception occurs for higher-order terms: since no orthogonal bases are
defined for abstractions, the symbol $\AbsBasis$ allows for linear reduction
over basis values. For instance, 
$(\Lam{x}{\AbsBasis}{\vec t}) \sum_i
\alpha_i(\Lam{y}{X}{\vec{s}_{i}}) \lra \sum_i\alpha_i\vec
t[\Lam{y}{X}{\vec{s}_{i}}/x]$.
%\[
%  (\Lam{x}{\AbsBasis}{\vec t})
%  \sum_{i=1}^{n}\alpha_i(\Lam{y}{X}{\vec{s}_{i}})
%  \lra
%  \sum_{i=1}^n\alpha_i\vec t[\Lam{y}{X}{\vec{s}_{i}}/x].
%\]

The $\mathsf{case}$ pattern matching controls program flow.  
Each operator keeps track of a set of orthogonal values and tests whether the
argument equals each vector, selecting the matching branch.  
If the argument is a linear combination of several vectors, the result is the
corresponding linear combination of the branches.  For example:
\[
  \case{\ket{-}}{\ket{0}}{\ket{1}}{\vec{t}_{1}}{\vec{t}_{2}} \lra
  \tfrac{1}{\sqrt{2}}\,\vec{t}_{1} - \tfrac{1}{\sqrt{2}}\,\vec{t}_{2}.
\]

The advantage over a binary conditional branching is the ability to match against
several vectors simultaneously. This characteristic comes into play when dealing
with vectors from entangled bases, in other words, bases which cannot be expressed
as products of smaller ones.  This general $\mathsf{case}$ matches directly
against those vectors.  For example, using the \emph{Bell basis}\footnote{ The four 
Bell states are $\Phi^{\pm}=\sqrthalf(\ket{00}\pm\ket{11})$ and $\Psi^{\pm}=
\sqrthalf(\ket{01}\pm\ket{10})$.}:
\[
  \mathsf{case}\;\vec{v}\;\mathsf{of}\;\{
  \Phi^+ \mapsto \vec{t}_{1}\ \mid
  \Phi^- \mapsto \vec{t}_{2}\ \mid
  \Psi^+ \mapsto \vec{t}_{3}\ \mid
  \Psi^- \mapsto \vec{t}_{4}\ \}.
\]

The Bell basis is central in quantum communication. In
Section~\ref{sec:teleportation}, we revisit the quantum teleportation protocol,
which relies heavily on these states.

Defining the system in this way yields a strategy within the
\emph{call-by-value} family, namely a generalisation of the
\emph{call-by-basis} strategy introduced in~\cite{ArrighiDowekLMCS17} and
further analysed in~\cite{AssafDiazcaroPerdrixTassonValironLMCS14}. While
call-by-basis fixes the canonical basis for evaluation, our variant,
which we call \emph{call-by-arbitrary-basis}, allows each abstraction to attach
its own orthonormal basis to its argument. Evaluation remains weak: no
reduction takes place under $\lambda$, pairs, $\mathsf{let}$, or
$\mathsf{case}$ constructors.

The congruence relation on terms gives rise to different redexes. We write
$\eval$ for the reflexive–transitive closure of $\lra$. We can show that the
equivalence relation $\equiv$ commutes with $\eval$; in other words,
equivalence is preserved by reduction modulo~$\equiv$.

\begin{theorem}[Reduction preserves equivalence]\label{thm:confluence}
  Let $\vec{t}$ and $\vec{s}$ be closed term combinations with
  $\vec{t}\equiv\vec{s}$. If $\vec{t}\lraneq\vv{t'}$ and $\vec{s}\lraneq\vv{s'}$,
  then there exist term combinations $\vec{r}_{1}$ and $\vec{r}_{2}$ such that
  $\vv{t'}\eval\vec{r}_{1}$, $\vv{s'}\eval\vec{r}_{2}$, and
  $\vec{r}_{1}\equiv\vec{r}_{2}$.
%  Diagrammatically:
%  \[
%    \begin{tikzcd}[row sep=1pt,baseline=(current bounding box.south)]
%      & \vec{t}
%        \arrow[ld,decorate,decoration={snake, amplitude=0.8, segment length=6pt}, ->]
%        &[-3em] \equiv
%        &[-3em] \vec{s}
%        \arrow[rd,decorate,decoration={snake, amplitude=0.8, segment length=6pt}, ->]
%        &\\
%      \vv{t'}\arrow[dr,"*",pos=0.9] & & & &
%      \vv{s'}\arrow[ld,"*"',pos=0.9] \\
%      & \vec{r}_{1} & \equiv & \vec{r}_{2} & 
%    \end{tikzcd}
%    \tag*{\smash{\raisebox{.6\baselineskip}{\qed}}}
%  \]
  \qed
\end{theorem}

\begin{remark}\label{rmk:determinism}
  Since the reduction relation $\eval$ is defined modulo~$\equiv$, the result
  above is equivalent to stating that there exists a single combination
  $\vec{r}$ such that $\vv{t'}\eval\vec{r}$ and $\vv{s'}\eval\vec{r}$.
  Moreover, as the elementary reduction~$\lraneq$ is deterministic, the
  reduction relation~$\lra$ is also deterministic; that is, if
  $\vec{t}\lra\vec{r}_{1}$ and $\vec{t}\lra\vec{r}_{2}$, then
  $\vec{r}_{1}\equiv\vec{r}_{2}$.
\end{remark}


\section{Realizability model}\label{sec:model}
\subsection{Unitary type semantics}
Given the deterministic machine presented in the previous section
(see~Remark~\ref{rmk:determinism}), the next
step towards extracting a typing system is to define the sets of values that
characterise its types. To achieve this, we first need to identify the notion
of a type.

Our aim is to define types exclusively inhabited by values of norm~1. The
vectors we wish to study belong to the \emph{unit sphere}. We denote $\Sph 
:= \{\vec v \in \vv{\Val} \mid \|\vec v\| = 1\}$, which
corresponds to the mathematical representation of quantum data as unit vectors
in a Hilbert space.

\begin{definition}[Unitary value combination]
  A value combination $\vec{v}$ is said to be \emph{unitary} if its norm equals~1,
  that is, if $\vec{v}\in\Sph$.
\end{definition}

\begin{definition}[Unitary type]
  A \emph{unitary type} (or simply a \emph{type}) is an expression~$A$
  with a set of unitary value combinations, denoted~$\sem{A}$, called the
  \emph{unitary semantics} of~$A$.
\end{definition}


We now turn to type realizers.  
Since the global phase of a quantum state has no physical significance, terms
that differ only by a phase should share the same type.  
Thus, we assign identical types to $\vec t$ and $e^{i\theta}\vec t$, a principle
that guides the definition of realizers.

\begin{definition}[Type realizer]\label{def:Realizer}
  Given a term combination~$\vec t$ and type~$A$, we say that~$\vec t$
  \emph{realizes}~$A$ (not.~$\vec t \real A$) if there exists a value
  combination~$\vec v\in\sem A$ such that
  $\vec{t}\eval e^{i\theta}\vec{v}$ for some $\theta\in\R$,
\end{definition}

\begin{table}[t]
  \centering
  $A := \basis{X} \mid A\Arr A \mid A\times A \mid \sharp A$, with $X$ an orthonormal basis.
  \vspace{2pt}
  \hrule
  \vspace{4pt}
  \begin{tabular}{l}
    $\sem{\basis{X}}:= X \qquad \sem{A\times B}:= \{(\vec v, \vec w) \mid \vec v\in\sem{A}, \vec w\in\sem{B}\}$ \\
    $\sem{A\Arr B}:= \{\sum_i \alpha_i(\Lam{x}{X}{\vec{t}_{i}})\in\Sph \mid \forall\vec{w}\in\sem{A}, (\sum_i \alpha_i \vec{t}_{i})\ansubst{\vec{w}/x}{X}\real B\}$ \\
    $\sem{\sharp Y}:= {(\sem{Y}^\bot)}^\bot$, with $Y^\bot = \{\vec{v}\in\Sph \mid \forall y\in Y, \scal{\vec{v}}{y} = 0\}$
  \end{tabular}
  \caption{Type notations and semantics.}
  \label{tab:UnitaryTypes}
\end{table}
With the notions of unitary types and their realizers in place, we can now
define a concrete approach for our language. We begin with the type grammar
given in~Table~\ref{tab:UnitaryTypes} and build a simple algebra from the sets of
values we aim to represent.

The types $\basis{X}$ serve as atomic types. Each of them represents a finite
set~$X$ of orthogonal vectors forming an orthonormal basis. For instance, a
boolean type can be represented by a basis of size 2, yet we are not restricted
to a single one, as infinitely many bases exist.

Pairs are standard, and written using the notation introduced earlier.
The arrow type~$A\Arr B$ consists of combinations of $\lambda$-abstractions
mapping elements of~$\sem{A}$ to realizers of~$B$.  
Finally, $\sharp A$ denotes the double orthogonal complement of~$A$ (within the unit sphere),
representing quantum data. Since these are linear resources, they cannot be erased or
duplicated.
Importantly, linear combinations of abstractions are not, in general, functions:
only those inhabiting $A \Arr B$ behave as maps, while general combinations
inhabit $\sharp(A \Arr B)$ and are interpreted as vectors.
Intuitively, applying the operator $\sharp$ captures the
linear combinations of values in the unitary semantics of $A$, as 
stated in the following theorem.

\begin{theorem}\label{thm:SharpCharacterization}
  The interpretation of a type~$\sharp A$ contains precisely the
  norm-$1$ linear combinations of values in~$\sem{A}$:
  \(
    \sem{\sharp A}
    = (\sem{A}^\bot)^\bot
    = \Span(\sem{A}) \cap \Sph
  \).
  \qed
\end{theorem}

As expected for a span, the $\sharp$ operator is idempotent.
\begin{theorem}\label{thm:IdempotentSharp}
  The~$\sharp$ operator is idempotent; that is,
  $\sem{\sharp A} = \sem{\sharp(\sharp A)}$.
  \qed
\end{theorem}

\begin{remark}
  A basis type~$\basis{X}$ may consist of value combinations of pairs and can
  therefore be written as the product type of smaller bases. For example, if
  $X=\{\ket{00},\ket{01},\ket{10},\ket{11}\}$, then $\sem{\basis{X}}=\sem{\basis{\B}\times\basis{\B}}$.
  However, entangled bases, such as the Bell basis, do not allow this.
\end{remark}

Our type algebra captures the intended sets of
value combinations. The following theorem shows that every member of a type
interpretation has norm~$1$.

\begin{theorem}\label{prop:UnitaryTypes}
  For every type~$A$, $\sem{A}\subseteq\Sph$.
  \qed
\end{theorem}

Defining types as sets of values naturally induces a semantic notion of
subtyping.
\begin{definition}[Subtyping and isomorphism]\label{def:Subtyping}
We say that a type~$A$ is a subtype of a type~$B$
(written~$A\leq B$) when $\sem A\subseteq\sem B$. 
If the two sets coincide, we say that $A$ and $B$ are \emph{isomorphic}
(written~$A\cong B$).
\end{definition}

\begin{example}
  For every type~$A$, we have $A\leq\sharp A$.
  For the base types $\basis{\B}$ and $\basis{\XB}$, 
  neither inclusion holds:
  $\basis{\B}\not\leq\basis{\XB}$ and
  $\basis{\XB}\not\leq\basis{\B}$. However, if they have the same dimension,
  $\sharp\basis{\B}\cong\sharp\basis{\XB}$.
\end{example}

Although every type is defined by norm-1 value combinations, not every
norm-1 combination is contained in the interpretation of a type.
For example, consider the combination
$\tfrac{1}{\sqrt{2}}(\ket{0} + \ket{00})$.
Another example is a linear combination of abstractions defined over different
bases. For instance, the term
\(
  \tfrac{1}{\sqrt{2}}(\Lam{x}{{\B}}{\pauliX{x}})
  + \tfrac{1}{\sqrt{2}}(\Lam{x}{{\XB}}{x})
\)
is not a member of an arrow type, since the bases decorating each abstraction
do not match. However, it is computationally equivalent to the abstraction
$(\Lam{x}{{\B}}{\ket{+}})$, which in turn belongs to the set
$\sem{\basis{\B}\Arr\basis{\XB}}$.


We denote by~$\Type$ the set of all types and by~$\BasisType$ the set of all
basis types~$\basis{X}$.

\subsection{Characterisation of unitary operators}

One of the main results of~\cite{DiazcaroGuillermoMiquelValironLICS19}
is characterising $\C^2\to\C^2$ unitary operators as values in
$\sem{\sharp\basis{\B}\Arr\sharp\basis{\B}}$~
\cite[Theorem IV.12]{DiazcaroGuillermoMiquelValironLICS19}.
Our goal is to extend the characterisation for unitary operators
$\C^{2^n}\to\C^{2^n}$ as abstractions of type $\sharp\basis{X}
\Arr\sharp\basis{Y}$, where both bases have size~$2^n$.

Unitary operators are isomorphisms between Hilbert spaces, as they preserve the
structure of the space. With this in mind, the first step is to show that the
members of $\sem{\sharp\basis{X}\Arr\sharp\basis{Y}}$ map basis vectors from
$\sem{\basis{X}}$ onto orthogonal vectors in $\sem{\sharp\basis{Y}}$. In other
words, these abstractions preserve both norm and orthogonality.

\begin{lemma}\label{lem:BasesIso}
  Let $X$ and $Y$ be orthonormal bases of the same finite
  dimension, and let $\Lam{x}{{X}}{\vec t}$ be a closed $\lambda$-abstraction.
  Then $\Lam{x}{{X}}{\vec t}\in\sem{\sharp\basis{X}\Arr\sharp\basis{Y}}$
  if and only if 
  for all $\vec{v}_{i},\vec{v}_{j}\in\sem{\basis{X}}$,
  there exist value combinations
  $\vec{w}_{i},\vec{w}_{j}\in\sem{\sharp\basis{Y}}$ such that,
    $\vec{t}[\vec{v}_{i}/x]\eval\vec{w}_{i}$
    and
    $\vec{t}[\vec{v}_{j}/x]\eval\vec{w}_{j}$,
    with
    $\vec{w}_{i}\perp\vec{w}_{j}$ whenever $i\neq j$.
    \qed
\end{lemma}

Terms such as $\ket{0}$ and $\ket{1}$ are syntactic objects of the
calculus, not vectors of~$\C^2$. Nevertheless, when discussing the behaviour of
terms on Hilbert spaces, we shall occasionally abuse notation and identify
value combinations representing qubits with their corresponding canonical
basis vectors in~$\C^{2^n}$. This identification applies only to those
combinations denoting quantum data, not to general values such as
$\lambda$-abstractions---in particular, to all elements of
$\sem{\sharp\basis{X}}$ for any orthonormal basis~$X$. This identification can
be made precise as follows.

\begin{definition}
  Let $X$ be an orthonormal basis of size~$2^n$. For every $\vec{v}\in X$, we can
  write
  \(
    \vec{v}\equiv \sum_{i=1}^{2^n}\alpha_i\ket{i}
  \),
  where $\ket{i}$ denotes the $n$-qubit tuple of $\ket{0}$ and $\ket{1}$
  corresponding to the binary representation of~$i$, and
  $\sum_{i=1}^{n}|\alpha_i|^2=1$. For example, for $n=4$, $\ket{3}$ is
  $\ket{0011}$. We then define $\pi_n:X\to\C^{2^n}$ by
  \(
    \pi_n(\vec{v}) = (\alpha_1,\dotsc,\alpha_{2^n})
  \).
\end{definition}

To lighten notation, we shall henceforth omit~$\pi_n$ and use~$\vec v$
directly. We now establish a correspondence between $\lambda$-abstractions and operators
on~$\C^{2^n}$. Intuitively, an abstraction represents a linear operator when its
operational behaviour coincides with the action of that operator on vectors.
Formally:

\begin{definition}
  A $\lambda$-abstraction $\Lam{x}{{X}}{\vec{t}}$ is said to \emph{represent}
  an operator $F:\C^{2^n}\to\C^{2^n}$ if
    $(\Lam{x}{{X}}{\vec{t}})\vec{v} \eval \vec{w}$
    if and only if
    $F(\vec{v}) = \vec{w}$.
\end{definition}

This definition, together with Lemma~\ref{lem:BasesIso}, allows us to build a
characterisation of unitary operators as values in
$\sem{\sharp\basis{X}\Arr\sharp\basis{Y}}$.

\begin{theorem}[Characterisation of Unitary Operators]\label{thm:characterisation}
  Let $X$ and $Y$ be orthonormal bases of size~$2^n$.
  A closed $\lambda$-abstraction
  $\Lam{x}{{X}}{\vec{t}}\in\sem{\sharp\basis{X}\Arr\sharp\basis{Y}}$
  if and only if it represents a unitary operator
  $F:\C^{2^n}\to\C^{2^n}$.
  \qed
\end{theorem}

  This result naturally extends to unitary combinations of
  $\lambda$-ab\-strac\-tions, since a term of the form
  $\Lam{x}{{X}}{\sum_{i=1}^{n}\alpha_i \vec{t}_{i}}$ is syntactically different
  but computationally equivalent to $\sum_{i=1}^{n}\alpha_i
  \Lam{x}{{X}}{\vec{t}_{i}}$.  Hence, the characterisation of unitary operators
  also applies to superpositions of abstractions sharing the same basis~$X$.

\subsection{Typing rules}    
In this section, we focus on enumerating and proving the validity of various
typing rules. The objective is to extract a reasonable set of rules that can
constitute a type system. We first need to lay the groundwork to properly define
what it means for a typing rule to be valid.

\begin{definition}
  A \emph{context} (denoted by capital Greek letters $\Gamma$, $\Delta$) is a
  finite mapping $\Gamma:\Var\to\Type\times\BasisType$ assigning a type and a
  basis to each variable in its domain. We write
  \(
    \Gamma = {x_1}^{{X_1}}:A_1,\dotsb, {x_n}^{{X_n}}:A_n
  \)
  to indicate that $\Gamma(x_i)=(A_i,\basis{X_i})$ for each~$i$.
\end{definition}

As with standard typing judgements, the context records the types of the free
variables in a term. However, since substitution in our calculus depends on a basis,
we also wish to record that information. This is not strictly necessary, as the
basis with respect to which a variable is interpreted should not affect its
type. For instance, consider the following substitutions:
\begin{align*}
(\Lam{x}{\B}{\Pair{x}{y}})\ansubst{\ket{0}/y}{\B}
  &= \Lam{x}{\B}{\Pair{x}{\ket{0}}},\\
\text{and}\quad
(\Lam{x}{\B}{(x, y)})\ansubst{\ket{0}/y}{\XB}
  &= \tfrac{1}{\sqrt{2}}
    \bigl((\Lam{x}{\B}{\Pair{x}{\ket{+}}})
         + (\Lam{x}{\B}{\Pair{x}{\ket{-}}})\bigr).
\end{align*}
These terms are not syntactically identical, yet they are computationally
equivalent. 
Since typing via realisability captures computational behaviour,
their types coincide. Nevertheless, we retain basis information in contexts, as
it will simplify the forthcoming proofs.

\begin{definition}
  Given a context~$\Gamma$, its \emph{unitary semantics},
  denoted~$\sem{\Gamma}$, is the set of substitutions defined by
  \begin{align*}
    \sem{\Gamma}
    := 
    \{&\sigma~\text{substitution} \mid 
      \dom{\sigma} = \dom{\Gamma}
      \text{ and } \forall {x_i} \in \dom{\Gamma},\\
      &\Gamma(x_i) = (A_i, \basis{X_i})
      \text{ implies }
      \sigma(x_i) = \ansubst{\vec{v}_{i}/x_i}{{X_i}}
      \text{ for some }\vec{v}_{i} \in \sem{A_i}\}.
  \end{align*}
\end{definition}


To ensure a coherent treatment of quantum data, we must guarantee that qubits
are handled linearly. The first step is to identify which variables in the
context represent quantum data---those associated with a type of the
form~$\sharp A$. We call the subset of~$\Gamma$ composed of such variables its
\emph{strict domain}.

\begin{definition}
  The \emph{strict domain} of a context~$\Gamma$, denoted~$\sdom{\Gamma}$, is
  defined as
  \(
    \sdom{\Gamma} :=
    \{x \in \dom{\Gamma} \mid
    \sem{\Gamma(x)} = \sem{\sharp(\Gamma(x))}\}
  \).\footnote{This definition relies on the idempotence of the~$\sharp$~operator
  (Theorem~\ref{thm:IdempotentSharp}).
  }
\end{definition}

A typing judgement $\Gamma\vdash \vec{t}:A$ is valid if it satisfies two
conditions.  First, every free variable of~$\vec{t}$ must belong to the domain
of~$\Gamma$, and every variable in the strict domain~$\sdom{\Gamma}$ must occur
in~$\vec{t}$.  This ensures that no information is erased and that all variables
are properly accounted for.  The linear treatment of quantum data is thus
enforced by the substitution.

Second, for every substitution in the unitary semantics of~$\Gamma$, applying
it to~$\vec{t}$ must yield a term that reduces to a realizer of type~$A$.  This
condition ensures that the operational behaviour of the term within the context
is faithfully captured by the type.  Formally:

\begin{definition}[Typing judgement]
  A \emph{typing judgement} $\TYP{\Gamma}{\vec{t}}{A}$ is valid when:
    $\sdom{\Gamma}\subseteq\FV{\vec{t}}\subseteq\dom{\Gamma}$
    and
    for all $\sigma\in\sem{\Gamma}$, 
          $\vec{t}\ansubst{\sigma}{}\real A$.
\end{definition}

We are also interested in \emph{orthogonal terms}, that is, terms that reduce
to orthogonal values.  
We therefore introduce the following notion.

\begin{definition}
  An \emph{orthogonality judgement}
  $\ORTH{\Gamma}{\Delta_1}{\vec{t}}{\Delta_2}{\vec{s}}{A}$
  is said to be \emph{valid} when the judgements
      $\TYP{\Gamma,\Delta_1}{\vec{t}}{A}$ and
      $\TYP{\Gamma,\Delta_2}{\vec{s}}{A}$ are valid, and
      for every
      $\sigma\in\sem{\Gamma,\Delta_1}$ and
      $\tau\in\sem{\Gamma,\Delta_2}$,
      there exist value combinations $\vec{v},\vec{w}$ such that
      $\vec{t}\ansubst{\sigma}{}\eval\vec{v}$,
      $\vec{s}\ansubst{\tau}{}\eval\vec{w}$,
      and $\vec{v}\perp\vec{w}$.
When both $\Delta_1$ and $\Delta_2$ are empty,
we just write
$\SORTH{\Gamma}{\vec{t}}{\vec{s}}{A}$.
\end{definition}



With these definitions in mind, a typing rule is \emph{valid} when valid
premises entail a valid conclusion.  
Although there are infinitely many valid rules (each corresponding to a theorem),
Table~\ref{tab:TypingRules} presents a representative subset forming a reasonable
core typing system for the calculus, whose validity is stated below.
This is the main result of the paper, as the validity of these rules ensures all the safety properties established in the realisability model.

%\begin{table}[t]
%  \[
%    \begin{array}{c}
%      \infer[\snam{Axiom}]{\TYP{x^{X}:A}{x}{A}}{\basis{X}\leq A \text{ or }X=\AbsBasis}
%      \qquad
%	\infer[\snam{UnitLam}]{\TYP{\Gamma}{\sum_{i=1}^n \alpha_i (\Lam{x}{{X}}{\vec{t}_{i}})}{A\Arr B}}
%	{\TYP{\Gamma,x^{X}:A}{\sum_{i=1}^{n}\alpha_i\vec{t}_{i}}{B}}
%      \\
%      \noalign{\medskip}
%      \infer[\snam{App}]{\TYP{\Gamma,\Delta}{\vec{s}\,\vec{t}}{B}}
%      {\TYP{\Gamma}{\vec{s}}{A\Arr B} & \TYP{\Delta}{\vec{t}}{A}}
%      \qquad
%      \infer[\snam{Pair}]{\TYP{\Gamma,\Delta}
%      {\Pair{\vec{t}}{\vec{s}}}{A\times B}}{
%	\TYP{\Gamma}{\vec{t}}{A}&\TYP{\Delta}{\vec{s}}{B}
%      }
%      \\
%      \noalign{\medskip}
%      \infer[\snam{LetPair}]{\TYP{\Gamma,\Delta} 
%      {\LetP{x}{{X}}{y}{{Y}}{\vec{t}}{\vec{s}}}{C}}{
%	\TYP{\Gamma}{\vec{t}}{A\times B}&
%	\TYP{\Delta,x^{{X}}:A,y^{Y}:B}{\vec{s}}{C}
%      }\\
%      \noalign{\medskip}
%      \infer[\snam{LetTens}]{\TYP{\Gamma,\Delta}
%      {\LetP{x}{{X}}{y}{{Y}}{\vec{t}}{\vec{s}}}{\sharp C}}{
%	\TYP{\Gamma}{\vec{t}}{\sharp(A\times B)}&
%	\TYP{\Delta,x^{X}:\sharp A,y^{Y}:\sharp B}{\vec{s}}{C}
%      }\\
%      \noalign{\medskip}
%      \infer[\snam{Case}]{\TYP{\Gamma,\Delta}
%      {\gencase{\vec{t}}{\vec{v}_{1}}{\vec{v}_{n}}{\vec{s}_{1}}{\vec{s}_{n}}}{A}}{
%	\TYP{\Gamma}{\vec{t}}{\genbasis{\vec{v}_{i}}{i=1}{n}}&
%	\forall i,\ \TYP{\Delta}{\vec{s}_{i}}{A}
%      }\\
%      \noalign{\medskip}
%      \infer[\snam{UnitCase}]{\TYP{\Gamma,\Delta}
%      {\gencase{\vec{t}}{\vec{v}_{1}}{\vec{v}_{n}}{\vec{s}_{1}}{\vec{s}_{n}}}{\sharp A}}{
%	\TYP{\Gamma}{\vec{t}}{\sharp \genbasis{\vec{v}_{i}}{i=1}{n}}&
%	\forall i\neq j,\ \SORTH{\Delta}{\vec{s}_{i}}{\vec{s}_{j}}{A}
%      }\\
%      \noalign{\medskip}
%      \infer[\snam{Sum}]
%      {\TYP{\Gamma}{\sum_{i=1}^{n}\alpha_i \vec{t}_{i}}{\sharp A}}
%      {\forall i\neq j,\, \SORTH{\Gamma}{\vec{t}_{i}}{\vec{t}_{j}}{A} &
%      \sum_{i=1}^{n}|\alpha_i|^2 = 1}
%      \\
%      \noalign{\medskip}
%      \infer[\snam{Contr}]{\TYP{\Gamma,x^{X}:{\basis{X}}}{\vec{t}\,[y:=x]}{B}}{
%	\TYP{\Gamma,x^{X}:{\basis{X}},y^{X}:{\basis{X}}}{\vec{t}}{B}
%      } 
%      \qquad
%	\infer[\snam{Weak}]{\TYP{\Gamma,x^{X}:\basis{X}}{\vec{t}}{C}}{
%	  \TYP{\Gamma}{\vec{t}}{C}
%	}
%      \\
%      \noalign{\medskip}
%      \infer[\snam{Sub}]{\TYP{\Gamma}{\vec{t}}{B}}{\TYP{\Gamma}{\vec{t}}{A} & \SUB{A}{B}}
%      \quad
%      \infer[\snam{Equiv}]{\TYP{\Gamma}{\vec{s}}{A}}{
%	\TYP{\Gamma}{\vec{t}}{A}& \vec t\equiv \vec s
%      }
%      \quad 
%      \infer[\snam{Phase}]{\TYP{\Gamma}{e^{i\theta}\vec{t}}{A}}
%      {\TYP{\Gamma}{\vec{t}}{A}}
%    \end{array}
%  \]
%  \caption{Some valid typing rules}
%  \label{tab:TypingRules}
%\end{table}

\begin{table}[t]
  \centering
  \begin{tabular}{c@{\quad}c}
    \infer[\snam{Axiom}]{\TYP{x^{X}:A}{x}{A}}{\basis{X}\leq A \text{ or }X=\AbsBasis} &
    \infer[\snam{UnitLam}]{\TYP{\Gamma}{\sum_i \alpha_i (\lambda x^X.\vec{t}_i)}{A\Arr B}}{\TYP{\Gamma,x^{X}:A}{\sum_i \alpha_i \vec{t}_i}{B}} 
    \\ \noalign{\small\smallskip}
    \infer[\snam{App}]{\TYP{\Gamma,\Delta}{\vec{s}\,\vec{t}}{B}}{\TYP{\Gamma}{\vec{s}}{A\Arr B} & \TYP{\Delta}{\vec{t}}{A}} &
    \infer[\snam{Pair}]{\TYP{\Gamma,\Delta}{\Pair{\vec{t}}{\vec{s}}}{A\times B}}{\TYP{\Gamma}{\vec{t}}{A}&\TYP{\Delta}{\vec{s}}{B}} 
    \\ \noalign{\small\smallskip}
    \multicolumn{2}{c}{
      \infer[\snam{LetPair}]{\TYP{\Gamma,\Delta}{\LetP{x}{X}{y}{Y}{\vec{t}}{\vec{s}}}{C}}{\TYP{\Gamma}{\vec{t}}{A\times B} & \TYP{\Delta,x^X:A,y^Y:B}{\vec{s}}{C}} \quad
      \infer[\snam{LetTens}]{\TYP{\Gamma,\Delta}{\LetP{x}{X}{y}{Y}{\vec{t}}{\vec{s}}}{\sharp C}}{\TYP{\Gamma}{\vec{t}}{\sharp(A\times B)} & \TYP{\Delta,x^X:\sharp A,y^Y:\sharp B}{\vec{s}}{C}}
    } \\ \noalign{\small\smallskip}
    \multicolumn{2}{c}{
      \infer[\snam{Case}]{\TYP{\Gamma,\Delta}{\gencasei{\vec{t}}{\vec{v}_i}{\vec{s}_i}}{A}}{\TYP{\Gamma}{\vec{t}}{\basis{\{\vec{v}_i\}}} & \forall i, \TYP{\Delta}{\vec{s}_i}{A}} \quad
      \infer[\snam{UnitCase}]{\TYP{\Gamma,\Delta}{\gencasei{\vec{t}}{\vec{v}_i}{\vec{s}_i}}{\sharp A}}{\TYP{\Gamma}{\vec{t}}{\sharp\basis{\{\vec{v}_i\}}} & \forall i\neq j, \SORTH{\Delta}{\vec{s}_i}{\vec{s}_j}{A}}
    } \\ \noalign{\small\smallskip}
    \infer[\snam{Sum}]{\TYP{\Gamma}{\sum_i \alpha_i \vec{t}_i}{\sharp A}}{\forall i\neq j, \SORTH{\Gamma}{\vec{t}_i}{\vec{t}_j}{A} & \sum |\alpha_i|^2 = 1} &
    \infer[\snam{Weak}]{\TYP{\Gamma,x^{X}:\basis{X}}{\vec{t}}{C}}{\TYP{\Gamma}{\vec{t}}{C}} 
    \\ \noalign{\small\smallskip}
    \infer[\snam{Contr}]{\TYP{\Gamma,x^{X}:\basis{X}}{\vec{t}[y:=x]}{B}}{\TYP{\Gamma,x^X:\basis{X},y^X:\basis{X}}{\vec{t}}{B}} &
    \infer[\snam{Sub}]{\TYP{\Gamma}{\vec{t}}{B}}{\TYP{\Gamma}{\vec{t}}{A} & A \leq B} 
    \\ \noalign{\small\smallskip}
    \infer[\snam{Equiv}]{\TYP{\Gamma}{\vec{s}}{A}}{\TYP{\Gamma}{\vec{t}}{A} & \vec{t} \equiv \vec{s}} &
    \infer[\snam{Phase}]{\TYP{\Gamma}{e^{i\theta}\vec{t}}{A}}{\TYP{\Gamma}{\vec{t}}{A}}
  \end{tabular}
  \caption{Some valid typing rules.}
  \label{tab:TypingRules}
\end{table}

\begin{theorem}\label{thm:TypingRulesValidity}
  All the typing rules in Table~\ref{tab:TypingRules} are valid.
  \qed
\end{theorem}


The usual safety properties follow straightforwardly in this framework.
\emph{Confluence} is an immediate consequence of the reduction being
deterministic (cf.~Remark~\ref{rmk:determinism}).  
\emph{Strong normalisation} follows directly from the definition of a
realizer (cf.~Definition~\ref{def:Realizer}).  
\emph{Subject reduction} is also immediate: indeed, if
$\TYP{\Gamma}{\vec{t}}{A}$ and $\vec{t}\to\vec{u}$, then
$\TYP{\Gamma}{\vec{u}}{A}$ by definition.
However, if we restrict ourselves to a subset of typing rules—such as those in Table~\ref{tab:TypingRules}—we must
ensure that this restricted system still suffices to type all reducts of a
term, once the underlying realisability semantics is abstracted away.  In our
case, the rules proven valid in Theorem~\ref{thm:TypingRulesValidity}
suffice to guarantee subject reduction, as
stated below.


\begin{theorem}[Subject reduction]\label{thm:SubjectReduction}
  If $\TYP{\Gamma}{\vec{t}}{A}$ can be derived using the set of rules in
  Table~\ref{tab:TypingRules} and $\vec{t}\to\vec{u}$, then
  $\TYP{\Gamma}{\vec{u}}{A}$ can also be derived by the same set of rules.
  \qed
\end{theorem}

\section{Examples}\label{sec:examples}
\subsection{Deutsch's algorithm}\label{subsec:deutsch}
Deutsch's algorithm decides if an unknown Boolean function $f:\{0,1\}\to\{0,1\}$ is constant or balanced with a single query by exploiting quantum superposition. Operationally, an oracle $U_f$ acts as $U_f:\ket{xy}\mapsto \ket{x}\otimes\ket{y\oplus f(x)}$. The standard circuit prepares $\ket{+-}$, applies $U_f$, and then applying a Hadamard gate ($H$) on the first qubit yields $\ket{0}$ if $f$ is constant and $\ket{1}$ if $f$ is balanced.

In $\lambdaB$, we first encode the standard gates in $\B$: 
$H := \Lam{x}{\B}{\case{x}{\ket{0}}{\ket{1}}{\ket{+}}{\ket{-}}}$, 
$\mathsf{NOT} := \Lam{x}{\B}{\case{x}{\ket{0}}{\ket{1}}{\ket{1}}{\ket{0}}}$, and 
$\mathsf{CNOT} := \Lam{x}{\B}{\Lam{y}{\B}{\case{x}{\ket{0}}{\ket{1}}{\Pair{\ket{0}}{y}}{\Pair{\ket{1}}{\mathsf{NOT}\,y}}}}$.
The four possible oracles can be modeled as abstractions of type $\basis{\B}\Arr\basis{\B}\Arr(\basis{\B}\times\basis{\B})$; for instance, 
$O_{\mathrm{const}\,0} := \Lam{x}{\B}{\Lam{y}{\B}{\Pair{x}{y}}}$ and 
$O_{\mathrm{id}} := \Lam{x}{\B}{\Lam{y}{\B}{\mathsf{CNOT}\,x\,y}}$. 
The standard Deutsch term is then:
\[
  \mathsf{Deutsch}_{\mathrm{std}} :=
  \Lam{f}{\AbsBasis}{
    \LetP{x}{\B}{y}{\B}{(f\,(H\ket{0})\,(H\ket{1}))}{\Pair{H\,x}{y}}
  }.
\]
Since the oracle is applied to $\ket{+}$ and $\ket{-}$, the derived type is 
$(\sharp\basis{\B}\Arr\sharp\basis{\B}\Arr \sharp(\basis{\B}\times\basis{\B})) \Arr \sharp(\basis{\B}\times\basis{\B})$. 
This type is coarse: it only ensures the result is a unitary combination of pairs, failing to capture the deterministic nature of the outcome.

By contrast, a basis-aware implementation in $\XB=\{\ket{+},\ket{-}\}$ provides tighter typing. We redefine the gates and oracles in $\XB$ (details in Appendix~\ref{sec:appendixD}); for example, 
$O_{\mathrm{const}\,0}^{\XB} := \Lam{x}{\XB}{\Lam{y}{\XB}{\Pair{x}{y}}}$ and 
$O_{\mathrm{id}}^{\XB} := \Lam{x}{\XB}{\Lam{y}{\XB}{\cnotXB{x}{y}}}$. The basis-aware program is:
\[
\mathsf{Deutsch} := \Lam{f}{\AbsBasis}{ \LetP{x}{\XB}{y}{\XB}{(f\,\ket{+}\,\ket{-})}
{\case{x}{\ket{+}}{\ket{-}}{\ket{0}}{\ket{1}}}}.
\]
This version is typed as $(\basis{\XB}\Arr\basis{\XB}\Arr(\basis{\XB}\times\basis{\XB}))\Arr\basis{\B}$. 
The gain is significant: the type system tracks that the oracle, when fed with basis states in $\XB$, produces a pair of basis states. This allows us to extract a classical bit ($\basis{\B}$) deterministically. Unlike the standard version, no variable in the derivation requires a $\sharp$-type, reflecting that the algorithm's result is classical and its deterministic behaviour is captured entirely at the type level.

%\subsection{Deutsch's algorithm}\label{subsec:deutsch}
%We begin with \emph{Deutsch's algorithm}, a canonical example that highlights
%how basis types in the $\lambdaB$ calculus yield more informative typings for
%quantum programs.
%The algorithm is as follows.
%We are given black-box access to an \emph{oracle}~$U_f$ that implements an
%unknown Boolean function $f:\{0,1\}\to\{0,1\}$.  
%The oracle can only be either \emph{constant} (both inputs map to the same
%output) or \emph{balanced} (the two outputs differ).  
%Classically, determining which case holds requires two queries to~$f$.  
%Deutsch's algorithm decides this with a single query by exploiting quantum
%superposition and interference.
%
%Operationally, the oracle~$U_f$ acts as
%\(
%  U_f:\ket{xy}\mapsto \ket{x}\otimes\ket{y\oplus f(x)}
%\),
%where $\oplus$ is addition modulo~2.  
%The textbook circuit prepares the state~$\ket{+-}$, applies~$U_f$, and
%then applies a Hadamard on the first qubit before measuring it.  
%The outcome is~$\ket{0}$ if~$f$ is constant and~$\ket{1}$ if~$f$ is balanced.
%To implement this algorithm in $\lambdaB$, we first encode the gates 
%we will use:
%\begin{align*}
%  \Hd &:= \Lam{x}{\B}{\case{x}{\ket{0}}{\ket{1}}{\ket{+}}{\ket{-}}},\\
%  \mathsf{NOT} &:= \Lam{x}{\B}{\case{x}{\ket{0}}{\ket{1}}{\ket{1}}{\ket{0}}},\\
%  \mathsf{CNOT} &:= \Lam{x}{\B}{\Lam{y}{\B}{
%    \case{x}{\ket{0}}{\ket{1}}
%      {\Pair{\ket{0}}{y}}
%      {\Pair{\ket{1}}{\pauliX{y}}}}}.
%\end{align*}
%
%We model the four possible oracles $U_f$ (two constant and two balanced):
%\begin{align*}
%  O_{\mathrm{const}\,0} &:= \Lam{x}{\B}{\Lam{y}{\B}{\Pair{x}{y}}},
%  & O_{\mathrm{id}}       &:= \Lam{x}{\B}{\Lam{y}{\B}{\cnot{x}{y}}},\\
%  O_{\mathrm{const}\,1} &:= \Lam{x}{\B}{\Lam{y}{\B}{\Pair{x}{(\pauliX{y})}}},
%  & O_{\mathrm{flip}}     &:= \Lam{x}{\B}{\Lam{y}{\B}{\cnot{x}{(\pauliX{y})}}}.
%\end{align*}
%
%The standard Deutsch term prepares~$\ket{+-}$, calls the oracle, then
%applies~$H$ on the first qubit and returns the (classical) first component:
%\[
%  \mathsf{Deutsch}_{\mathrm{std}} :=
%  \Lam{f}{\AbsBasis}{
%    \LetP{x}{\B}{y}{\B}
%      {(f\,(\Hd\,\ket{0})\,(\Hd\,\ket{1}))}
%      {\Pair{\Hd\,x}{y}}
%  }.
%\]
%
%The oracles can be typed as 
%\(
%  \basis{\B}\Arr\basis{\B}\Arr(\basis{\B}\times\basis{\B})
%\).
%However, since we apply them to $\ket{+}$ and $\ket{-}$, 
%the derived judgement is
%\(
%  (\sharp\basis{\B}\Arr\sharp\basis{\B}\Arr(\sharp\basis{\B}\times\sharp\basis{\B}))
%   \Arr \sharp(\basis{\B}\times\basis{\B})
%\).
%This type is correct but coarse: it only ensures the result is a
%unitary combination of boolean pairs. Operationally, the
%first output is a basis state ($\ket{0}$ or~$\ket{1}$) indicating if
%$f$ is constant or balanced, while the second may be ignored.
%
%Then, we can consider a basis-aware implementation with tighter typing.
%The key observation is that the oracle is always called on the fixed state
%$\ket{+-}$. Thus it is natural to write the program in the
%$\XB=\{\ket{+},\ket{-}\}$ basis, letting the types track that we remain in a
%basis state at the oracle boundary. We first define the same gates in the $\XB$ basis:
%\begin{align*}
%  \pauliZXB &:= \Lam{x}{\XB}{\case{x}{\ket{+}}{\ket{-}}{\ket{-}}{\ket{+}}},\\
%  \pauliXXB &:= \Lam{x}{\XB}{\case{x}{\ket{+}}{\ket{-}}{\ket{+}}{-1\,\ket{-}}},\\
%  \cnotXB &:= \Lam{x}{\XB}{\Lam{y}{\XB}{
%    \case{y}{\ket{+}}{\ket{-}}
%      {\Pair{x}{\ket{+}}}
%      {\Pair{\pauliZXB{x}}{\ket{-}}}
%  }}.
%\end{align*}
%We then rewrite the program and the four oracles in~$\XB$:
%\[
%\mathsf{Deutsch} := \Lam{f}{\AbsBasis}{ \LetP{x}{\XB}{y}{\XB}{(f\,\ket{+}\,\ket{-})}
%{\case{x}{\ket{+}}{\ket{-}}{\ket{0}}{\ket{1}}}},\\
%\]
%\begin{align*}
%  O_{\mathrm{const}\,0}^{\XB} &:= \Lam{x}{\XB}{\Lam{y}{\XB}{\Pair{x}{y}}},
%  &O_{\mathrm{id}}^{\XB}\ \ &:= \Lam{x}{\XB}{\Lam{y}{\XB}{\cnotXB{x}{y}}},\\
%  O_{\mathrm{const}\,1}^{\XB} &:= \Lam{x}{\XB}{\Lam{y}{\XB}{\Pair{x}{(\pauliXXB{y})}}},
%  &O_{\mathrm{flip}}^{\XB} &:= \Lam{x}{\XB}{\Lam{y}{\XB}{\cnotXB{x}{(\pauliXXB{y})}}}.
%\end{align*}
%
%Now every oracle has the tight type
%\(
%  \basis{\XB}\Arr\basis{\XB}\Arr(\basis{\XB}\times\basis{\XB})
%\),
%and the program itself can be typed with
%\(
%(\basis{\XB}\Arr\basis{\XB}\Arr(\basis{\XB}\times\basis{\XB}))\Arr\basis{\B}.
%\)
%Intuitively, the oracle---when fed with $\ket{+-}$---produces a pair of
%basis states in~$\XB$ (up to a global phase).  Hence we can treat its output
%classically: a single $\mathsf{case}$ on the first component suffices to return
%a classical bit in the computational basis, with no need for a $\sharp$-type on
%the result.
%
%Both implementations are operationally equivalent: they compute a bit that
%decides whether $f$ is constant or balanced.  The difference lies in the
%\emph{precision} of their typings.  The standard version, expressed in~$\B$
%with explicit Hadamards, forces $\sharp$ on the oracle's interface and thus
%yields a result in~$\sharp(\basis{\B}\times\basis{\B})$.  The basis-aware
%version states, via types, that the oracle is used on a fixed $\XB$-input and
%therefore returns an $\XB$-basis pair; this lets us deterministically extract a
%$\B$-basis bit.  In the typing derivation, no variable carries a $\sharp$-type:
%(i) we may safely ignore the second qubit, and (ii) the first qubit is
%guaranteed to be classical in~$\B$, reflecting the determinism of Deutsch's
%algorithm.

\subsection{Quantum teleportation}\label{sec:teleportation}
The quantum teleportation protocol~\cite{PhysRevLett.70.1895} exemplifies transmitting quantum data via shared entanglement. In $\lambdaB$, this protocol illustrates how the \emph{deferred-measurement principle}~\cite[\S 4.4]{NC00} is internalised within a purely unitary framework. While measurement can be integrated into basis-sensitive calculi (as in~\cite{DiazcaroMonzonAPLAS25}), we focus on demonstrating that the $\mathsf{case}$ constructor suffices to reason about classically controlled gates via branching on basis states.

By delayining measurement, we define a term that explicitly describes the computation for each coherent branch of the superposition:
\begin{align*}
&\mathsf{Teleport}_{\Phi^+} :=
  \Lam{x}{\B}{
    \LetP{y_1}{\B}{y_2}{\B}{\Phi^{+}}{
      \mathsf{case}\; \Pair{x}{y_1}\; \mathsf{of}\; \{ \\
         &\Phi^{+} \mapsto \Pair{\Phi^{+}}{y_2},\, \Phi^{-} \mapsto \Pair{\Phi^{-}}{Z y_2},\,
	  \Psi^{+} \mapsto \Pair{\Psi^{+}}{X y_2},\, \Psi^{-} \mapsto \Pair{\Psi^{-}}{ZX y_2} \}
    }
  }.
\end{align*}
This term represents the protocol by encoding the coherent superposition of outcomes and their corresponding unitary corrections ($I, X, Z, ZX$). 

The strength of this approach lies in the type system's ability to verify the consistency of the entire superposition: each branch is statically checked as a unitary transformation. The overall term is assigned the precise type $\sharp\basis{\B}\Arr(\sharp\basis{\Bell}\times\sharp\basis{\B})$, where $\sharp\basis{\Bell}$ captures the entanglement of the first two qubits (relying on the subtypings 
$\basis{\Bell}\leq\sharp(\basis{\B}\times\sharp\basis{\B})$ and
$\sharp\basis{\B}\times\sharp\basis{\B}\leq\sharp\basis{\Bell}$).

%
%\subsection{Quantum teleportation}\label{sec:teleportation}
%We now turn to the \emph{quantum teleportation protocol}, a cornerstone of
%quantum information theory that exemplifies the manipulation of entangled
%states and the transmission of quantum data through classical communication.
%Within the $\lambdaB$ calculus, this protocol provides a natural setting to
%combine pattern matching, linear handling of qubits, and the
%\emph{deferred-measurement principle}.  
%Using the $\mathsf{case}$ constructor together with the expressive typing
%discipline introduced earlier, we can encode the teleportation process in a
%way that remains both syntactically compact and semantically faithful to its
%quantum-mechanical counterpart.
%
%The deferred-measurement principle (see, for example,~\cite[\S 4.4]{NC00})
%states that any quantum circuit can delay its measurements without altering
%the final outcome. More precisely, any gate classically controlled by the
%result of a measurement is equivalent to another circuit where the control
%qubit remains unmeasured, acting coherently on all branches of the superposition.  
%Although the $\lambdaB$ calculus has no primitive operation for measurement,
%the $\mathsf{case}$ constructor allows us to simulate such classically
%controlled gates by branching on basis states.
%
%A canonical example that exploits this principle is the \emph{quantum teleportation
%protocol}. Two agents (Alice and Bob) share an entangled pair of qubits forming a Bell
%state. Using this shared entanglement and two bits of classical information, Alice can
%transmit an unknown quantum state~$\ket{\psi}$ to Bob without physically sending the qubit itself.  
%
%The algorithm first creates the Bell state $\Phi^{+}$ between the second and third qubits,
%then performs a Bell-basis measurement on the first and second qubits.  
%Operationally, this measurement is implemented by applying a CNOT gate followed
%by a Hadamard on the first qubit (the adjoint of Bell-state preparation), and
%then measuring both qubits. Depending on the pair of classical outcomes, a Pauli correction
%($I$, $X$, $Z$, or $ZX$) is applied to Bob's qubit to recover the original state~$\ket{\psi}$
%(for a more thorough analysis of the algorithm, see \cite{PhysRevLett.70.1895}).
%
%We can simulate this protocol in $\lambdaB$ by
%defining a term that, instead of performing measurements, explicitly
%describes the computation corresponding to each branch:
%\begin{align*}
%&\mathsf{Teleport}_{\Phi^+} :=
%  \Lam{x}{\B}{
%    \LetP{y_1}{\B}{y_2}{\B}{\Phi^{+}}{
%      \mathsf{case}\;
%        \Pair{x}{y_1}\;
%        \mathsf{of}\\
%         &\{ \Phi^{+} \mapsto \Pair{\Phi^{+}}{y_2},\;
%          \Phi^{-} \mapsto \Pair{\Phi^{-}}{Z\,y_2},\;
%	  \Psi^{+} \mapsto \Pair{\Psi^{+}}{X\,y_2},\;
%	\Psi^{-} \mapsto \Pair{\Psi^{-}}{ZX\,y_2} \}
%    }
%  }.
%\end{align*}
%
%This term describes the quantum teleportation protocol in terms of
%corrections on the third qubit based on the result of a Bell basis measurement.
%It takes the input~$\ket{\psi}$ and pairs it with one half of an
%entangled Bell pair.  The $\mathsf{case}$ construct then matches the first two
%qubits (the input and Alice's entangled qubit) against the Bell basis, and in
%each branch applies the appropriate correction to Bob's qubit to
%recover~$\ket{\psi}$. The resulting term represents the same quantum
%transformation but expressed without any explicit measurement---instead,
%each branch encodes the coherent superposition of possible outcomes.
%
%The $\lambdaB$ calculus allows us to abstract both the encoding and decoding
%steps in the Bell basis, exploiting the deferred-measurement principle in a
%type-safe way.  
%Each branch of the $\mathsf{case}$ corresponds to a unitary transformation
%preserving linearity and orthogonality, and the overall term has type
%${\sharp\basis{\B}\Arr(\sharp\basis{\Bell}\times\sharp\basis{\B})}$
%\footnote{The type derivation relies on the subtypings 
%$\basis{\Bell}\leq\sharp(\basis{\B}\times\sharp\basis{\B})$ and
%$\sharp\basis{\B}\times\sharp\basis{\B}\leq\sharp\basis{\Bell}$.}.

\section{Conclusion}\label{sec:conclusion}
We have explored a quantum-control $\lambda$-calculus where abstractions are expressed relative to arbitrary bases. The central mechanism is the decoration of $\lambda$-abstractions and $\mathsf{let}$-constructors with basis annotations, supported by a basis-dependent substitution. While this does not increase raw expressivity, it offers a novel perspective for reasoning about quantum programs under basis changes. The reduction system ensures that evaluation commutes with the $\equiv$-relation, meaning that interpreting a value combination in a different basis does not affect the computational result, thus justifying reasoning modulo basis congruence.

The benefits of this design are manifest in the realisability model. Inclusion atomic types $\basis{X}$ enables a direct characterisation of abstractions representing unitary operators, generalising results from~\cite{DiazcaroGuillermoMiquelValironLICS19} via a simpler, more transparent proof. Based on this interpretation, we have established the validity of the typing rules in Table~\ref{tab:TypingRules}, ensuring their soundness and providing a principled foundation for a typed language based on the calculus.

Finally, we illustrated the system's advantages through two canonical examples: in \emph{Deutsch's algorithm}, basis-aware typing captures the algorithm's determinism by treating the result classically; in \emph{quantum teleportation}, the $\mathsf{case}$ construct simulates gates controlled by Bell-basis measurements, internalising the deferred-measurement principle. These results show that basis-sensitivity is a powerful tool for the static analysis of quantum protocols.

%\section{Conclusion}\label{sec:conclusion}
%
%In this paper we have explored a quantum-control $\lambda$-calculus equipped
%with the feature of allowing abstractions to be expressed relative
%to arbitrary bases, beyond the canonical one.  
%
%The central mechanism enabling this extension is the decoration of
%$\lambda$-ab\-strac\-tions and $\mathsf{let}$-constructors with basis annotations,
%alongside a modified substitution operation that governs how value
%combinations decompose onto different bases.  
%These additions do not increase the expressivity of the original calculus
%on which $\lambdaB$ builds, yet they offer a novel perspective for reasoning
%about quantum programs and their behaviour under basis changes.
%
%The reduction system coordinates computation through these extended syntactic
%constructs and substitutions.  
%A key property is that evaluation commutes with the $\equiv$-relation,
%ensuring that interpreting a value combination in a different basis does not
%affect the computational result. Then, it is meaningful to reason about terms modulo basis
%congruence.
%
%The benefit of this design becomes clear in the realisability model.
%The inclusion of atomic types~$\basis{X}$ enables a direct characterisation of
%abstractions representing unitary operators---our main semantic result,
%generalising the characterisation from~\cite{DiazcaroGuillermoMiquelValironLICS19}.  
%Here, the use of basis types yields a simpler and more transparent proof.  
%
%The second major result is the validity of the typing rules presented in
%Table~\ref{tab:TypingRules}.  
%By deriving these rules from the realisability interpretation, we ensure their
%soundness and obtain a principled foundation for a typed programming language
%based on the calculus.
%
%Finally, we have illustrated the expressive advantages of the system through two
%canonical examples.  
%In the case of \emph{Deutsch's algorithm}, the use of basis-aware typing allows
%the result to be treated classically, reflecting the algorithm's determinism.  
%In the case of \emph{quantum teleportation}, we show how the
%$\mathsf{case}$ construct can simulate gates controlled by Bell-basis
%measurements, effectively capturing the deferred-measurement principle within
%the calculus.

\bibliography{basisSensitive}

\appendix

\section{Omitted proofs from Section~\ref{sec:calculus}}\label{sec:appendixA}

\begin{restatetheorem}[Restatement of Theorem~\ref{thm:UniqueDecomposition}]
  \itshape
  If $B$ is an $2^n$-dimensional basis, then every $n$-qubit state has
  a unique decomposition over $B$.
\end{restatetheorem}
\begin{proof}
  Let $\vec{b}_{i}$ be the basis vectors of $B$. Suppose
  $\sum_{i=1}^{2^n} \alpha_i \vec{b}_{i}$ and $\sum_{i=1}^{2^n} \beta_i \vec{b}_{i}$
  are two decompositions of $\vec{v}$ over $B$. Then
  \[
    0=\vec{v}-\vec{v}=\sum_{i=1}^{2^n}(\alpha_i-\beta_i)\vec{b}_{i}.
  \]
  By linear independence, $\alpha_i=\beta_i$ for all $i$.
\end{proof}

\begin{restatecorollary}[Restatement of Corollary~\ref{cor:EquivalentDecomposition}]
  \itshape
  If $\vec{v}\equiv\vec{w}$, then they share the same decomposition over any
  basis $B$.
\end{restatecorollary}
\begin{proof}
  Since $\vec{v}-\vec{w}\equiv\vec{v}-\vec{v}\equiv\vec{w}-\vec{w}$, the same
  argument as in Theorem~\ref{thm:UniqueDecomposition} shows that $\vec{v}$ and $\vec{w}$ have the
  same decomposition over $B$.
\end{proof}

\begin{restatelemma}[Restatement of Lemma~\ref{lem:distributiveSubstitution}]
  \itshape
  For term combinations $\vec{t}_{i}$, a value combination $\vec{v}$, a
  variable $x$, coefficients $\alpha_i\in\C$, and a basis $B$ such that
  $\ansubst{\vec v/x}{B}$ is defined:
  \[
    \Bigl(\sum_i \alpha_i\vec{t}_{i}\Bigr)\ansubst{\vec v/x}{B}
    \equiv
    \sum_i \alpha_i\vec{t}_{i} \ansubst{\vec v/x}{B}.
  \]
\end{restatelemma}
\begin{proof}
  Let $B\neq\AbsBasis$ and
  $\vec{v}\equiv\sum_{j=1}^n\beta_j\vec{b}_{j}$ with each $\vec{b}_{j}\in B$.
  Then
  \begin{align*}
    \Bigl(\sum_i \alpha_i\vec{t}_{i}\Bigr)\ansubst{\vec v/x}{B}
    &= \sum_{j=1}^n \beta_j\Bigl(\sum_{i=1}^{n}\alpha_i t_i\Bigr)[\vec{b}_{j}/x]\\
    &\equiv \sum_{i=1}^{n}\alpha_i\Bigl(\sum_{j=1}^n\beta_j t_i[\vec{b}_{j}/x]\Bigr)
    = \sum_{i=1}^{n}\alpha_i\vec{t}_{i}\ansubst{\vec v/x}{B}.
  \end{align*}
  The case $B=\AbsBasis$ is analogous.
\end{proof}

\begin{restatelemma}[Restatement of Lemma~\ref{lem:EquivSubstitutions}]
  \itshape
  For value combinations $\vec{v},\vec{w}$, a term combination $\vec{t}$, and
  an orthonormal basis $B$ such that both
  $\ansubst{\vec{v}/x}{B}$ and $\ansubst{\vec{w}/x}{B}$ are defined:
  if $\vec{v}\equiv\vec{w}$, then
  $\vec{t}\ansubst{\vec{v}/x}{B}
  =\vec{t}\ansubst{\vec{w}/x}{B}$.
\end{restatelemma}
\begin{proof}
  Since $\vec{v}\equiv\vec{w}$, by
  Corollary~\ref{cor:EquivalentDecomposition},
  both can be written as
  $\vec{v}\equiv\vec{w}\equiv\sum_{i=1}^{n}\alpha_i\vec{b}_{i}$ with
  $\vec{b}_{i}\in B$. Hence
  \[
    \vec{t}\ansubst{\vec{v}/x}{B}
    = \sum_{i=1}^{n}\alpha_i\vec{t}[\vec{b}_{i}/x]
    = \vec{t}\ansubst{\vec{w}/x}{B}.
    \tag*{\qedhere}
  \]
\end{proof}

\section{Omitted proofs from Section~\ref{sec:reduction}}\label{sec:appendixB}

\begin{lemma}[Weak diamond property for $\lraneq$]\label{lem:SquigDiamond}
  Let $\vec{t}, \vec{s}_{1}, \vec{s}_{2}$ term combinations such that $\vec{t}\lraneq{s_1}$ and $\vec{t}\lraneq\vec{s}_{2}$. 
  Then, either there exists a term combination $\vec{r}$ such that $\vec{s}_{1}\lraneq\vec{r}$ and 
  $\vec{s}_{2}\lraneq \vec{r}$. Or, $\vec{s}_{1}=\vec{s}_{2}$. Diagrammatically:
  \[
    \begin{tikzcd}
      & \vec{t}
        \arrow[ld,decorate,decoration={snake, amplitude=0.8, segment length=6pt}, ->]
        \arrow[rd,decorate,decoration={snake, amplitude=0.8, segment length=6pt}, ->]
      &\\
      \vec{s}_{1}\arrow[dr,dashed,decorate,decoration={snake, amplitude=0.8, segment length=6pt}, ->] & &
      \vec{s}_{2}\arrow[ld,dashed,decorate,decoration={snake, amplitude=0.8, segment length=6pt}, ->] \\
      & \vec{r} &
    \end{tikzcd}
    \quad\text{ Or }\quad
    \begin{tikzcd}
      &[-2em] \vec{t}
        \arrow[ld,decorate,decoration={snake, amplitude=0.8, segment length=6pt}, ->]
        \arrow[rd,decorate,decoration={snake, amplitude=0.8, segment length=6pt}, ->]
      &[-2em]\\
      \vec{s}_{1}& = &\vec{s}_{2}\\
    \end{tikzcd}
  \]
\end{lemma}

\begin{proof}
  The proof follows from the fact that the $\lraneq$ reduction is deterministic over basic values. 
  And, in case of term combinations, we only need to match the reduction on the corresponding sub-terms. 
  Let $\vec{t}=\sum\limits_{i=1}^n \alpha_i \vec{t}_{i}$, $\vec{s}_{1}=\sum\limits_{i=1; i\neq j}^n \alpha_i \vec{t}_{i} + \alpha_j\vec{s}_{j}$, and $\vec{s}_{2}=\sum\limits_{i=1; i\neq k}^n \alpha_i \vec{t}_{i}+\alpha_k\vec{s}_{k}$. Where $\vec{t}_{j}\lraneq\vec{s}_{j}$ and $\vec{t}_{k}\lraneq\vec{s}_{k}$. If $j=k$ we are done, so we consider the case where $j\neq k$. Diagrammatically:
  \[
    \begin{tikzcd}
      &[-1em] \sum\limits_{i=1}^n \alpha_i \vec{t}_{i}
        \arrow[ld,decorate,decoration={snake, amplitude=0.8, segment length=6pt}, ->]
        \arrow[rd,decorate,decoration={snake, amplitude=0.8, segment length=6pt}, ->]
      &[-1em]\\
      \sum\limits_{i=1; i\neq j}^n \alpha_i \vec{t}_{i} + \alpha_j\vec{s}_{j}
      \arrow[dr,dashed,decorate,decoration={snake, amplitude=0.8, segment length=6pt}, ->] & &
      \sum\limits_{i=1; i\neq k}^n \alpha_i \vec{t}_{i} + \alpha_k\vec{s}_{k}
      \arrow[ld,dashed,decorate,decoration={snake, amplitude=0.8, segment length=6pt}, ->] \\
      & \sum\limits_{i=1; i\neq j,k}^n \alpha_i \vec{t}_{i} + \alpha_j\vec{s}_{j}+\alpha_k\vec{s}_{k} &
    \end{tikzcd}
  \]
\end{proof}

\begin{restatetheorem}[Restatement of Theorem~\ref{thm:confluence}]
  \itshape
  Let $\vec{t}$ and $\vec{s}$ be closed term combinations with
  $\vec{t}\equiv\vec{s}$. If $\vec{t}\lraneq\vv{t'}$ and $\vec{s}\lraneq\vv{s'}$,
  then there exist term combinations $\vec{r}_{1}$ and $\vec{r}_{2}$ such that
  $\vv{t'}\eval\vec{r}_{1}$, $\vv{s'}\eval\vec{r}_{2}$, and
  $\vec{r}_{1}\equiv\vec{r}_{2}$.
  Diagrammatically:
  \[
    \begin{tikzcd}
      & \vec{t}
        \arrow[ld,decorate,decoration={snake, amplitude=0.8, segment length=6pt}, ->]
        &[-2.5em] \equiv
        &[-2.5em] \vec{s}
        \arrow[rd,decorate,decoration={snake, amplitude=0.8, segment length=6pt}, ->]
        &\\
      \vv{t'}\arrow[dr,"*",pos=0.9] & & & &
      \vv{s'}\arrow[ld,"*"',pos=0.9] \\
      & \vec{r}_{1} & \equiv & \vec{r}_{2} &
    \end{tikzcd}
  \]
\end{restatetheorem}
\begin{proof}
  We do a case-by-case analysis over the relation $\vec{t}\equiv\vec{s}$.
  \begin{description}
    \item[$\vec{t}_{1} + 0\vec{t}_{2}\equiv\vec{t}_{1}$:] This case follows from Lemma \ref{lem:SquigDiamond} since the reductions can only be performed in $\vec{t}_{1}$.
    
    \item[$0\vec{t}\equiv\vec{0}$:] The term combinations cannot reduce on either side of the equivalence.
    
    \item[$1\vec{t}\equiv\vec{t}$:] This case follows from Lemma \ref{lem:SquigDiamond}.
    
    \item[$\alpha(\beta \vec{t})\equiv\delta\vec{t}$:] This case follows from Lemma \ref{lem:SquigDiamond}.
    
    \item[$\vec{t}_{1}+\vec{t}_{2}\equiv\vec{t}_{2}+\vec{t}_{1}$:] This case follows from Lemma \ref{lem:SquigDiamond}. We just have to match the reductions on both sides of the equivalence.
    
    \item[$\vec{t}_{1}+(\vec{t}_{2}+\vec{t}_{3})\equiv(\vec{t}_{1}+\vec{t}_{2})+\vec{t}_{3}$:] This case follows from Lemma \ref{lem:SquigDiamond}. We just have to match the reductions on both sides of the equivalence.
    
    \item[$(\alpha+\beta)\vec{t}\equiv\vec{t}$:] We start analyzing the coefficients. If $\alpha+\beta = 0$, then there cannot be a reduction on the left hand-side. If $(\alpha + \beta)\neq 0$ and either $\alpha=0$ or $\beta=0$, then we are on a particular case of $\vec{t}_{1} + 0\vec{t}_{2}\equiv\vec{t}_{1}$ with $\vec{t}_{1}=\vec{t}_{2}$. Otherwise, we match the reductions on both sides of the equivalence with Lemma \ref{lem:SquigDiamond}.
    
    \item[$\alpha(\vec{t}_{1}+\vec{t}_{2})\equiv\alpha\vec{t}_{1}+\alpha\vec{t}_{2}$:] If $\alpha=0$, then the term combinations cannot reduce on either side of the equivalence. Otherwise, we match the reductions on both sides of the equivalence with Lemma \ref{lem:SquigDiamond}.
    
    \item[$\vec{t} (\alpha\vec{s})\equiv\alpha(\vec{t}\vec{s})$:] If $\alpha=0$, then there is no reduction possible on the right-hand side. If there is an internal reduction on either $\vec{s}$ or $\vec{t}$, then we match the reductions on both sides of the equivalence with Lemma \ref{lem:SquigDiamond}.
    
    If $\vec{t} = (\Lam{x}{B}{\vec{t}_{1}})$ and $\vec{s}=\vec{v}$ and $\vec{t}_{1}\ansubst{\vec{v}/x}{B}$, is defined then (we consider the case $B\neq\AbsBasis$):
    \begin{align*}
      (\Lam{x}{B}{\vec{t}_{i}}) (\alpha\vec{v}) &\lraneq \vec{t}_{1}\ansubst{\alpha\vec{v}/x}{B}\\
      &= \sum_{i=1}^{n} \alpha\beta_i \vec{t}_{1}[\vec{b}_{i}/x]\quad \text{with }\vec{v}\equiv\sum_{i=1}^{n} \beta_i \vec{b}_{i} \text{ with } \vec{b}_{i}\in B
    \end{align*}
    On the other side:
    \begin{align*}
      \alpha ((\Lam{x}{B}{\vec{t}_{1}}) \vec{v}) &\lraneq \alpha(\vec{t}_{1}\ansubst{\vec{v}/x}{B})\\
      &=\alpha(\sum_{i=1}^{n} \beta_i \vec{t}_{1}[\vec{b}_{i}/x])\quad \text{with }\vec{v}\equiv\sum_{i=1}^{n} \beta_i \vec{b}_{i} \text{ with } \vec{b}_{i}\in B
    \end{align*}

    And we have that both terms are equivalent. The case for $B=\AbsBasis$ is similar.

    \item[$(\alpha\vec{t})\vec{s}\equiv\alpha(\vec{t}\vec{s})$:] If $\alpha=0$, then there is no reduction possible on the right-hand side. If there is an internal reduction on either $\vec{s}$ or $\vec{t}$, then we match the reductions on both sides of the equivalence with Lemma \ref{lem:SquigDiamond}. There are no other possible redexes since the abstraction must be a basic value to reduce on the left hand-side.
    
    \item[$(\vec{t}+\vec{s})\vec{r}\equiv \vec{t}\vec{s} + \vec{t}\vec{r}$:] If there is an internal reduction on either $\vec{t}$,$\vec{s}$ or $\vec{r}$, then we match the reductions on both sides of the equivalence with Lemma \ref{lem:SquigDiamond}. There are no other possible redexes since the abstraction must be a basic value to reduce on the left hand-side.
    
    \item[$\vec{t}(\vec{s}+\vec{r})\equiv\vec{t}\vec{s} + \vec{t}\vec{r}$:] If there is an internal reduction on either $\vec{t}$, $\vec{s}$ or $\vec{r}$, then we match the reductions on both sides of the equivalence with Lemma \ref{lem:SquigDiamond}.
    
    If $\vec{t} = (\Lam{x}{B}{\vec{t}_{1}})$, $\vec{s}=\vec{v}$, and $\vec{r}=\vec{w}$ with
    $\vec{t}_{1}\ansubst{\vec{v}/x}{B}$ and $\vec{t}_{1}\ansubst{\vec{w}/x}{B}$ defined then (we consider the case where $B\neq\AbsBasis$):
    \begin{align*}
      (\Lam{x}{B}{\vec{t}_{1}}) (\vec{v}+\vec{w}) &\lraneq \vec{t}_{1}\ansubst{\vec{v}+\vec{w}/x}{B}\\
      &=\sum_{i=1}^{n}(\alpha_i+\beta_i)\vec{t}_{1}[\vec{b}_{i}/x]\\
      &\quad\text{where }\vec{v}\equiv\sum_{i=1}^{n}\alpha_i\vec{b}_{i},\vec{w}\equiv\sum_{i=1}^{n}\beta_i\vec{b}_{i}\text{ with }\vec{b}_{i}\in B 
    \end{align*}

    On the other side:
    \begin{align*}
      (\Lam{x}{B}{\vec{t}_{1}}) \vec{v} + (\Lam{x}{B}{\vec{t}_{1}}) \vec{w} &\lraneq \vec{t}_{1}\ansubst{\vec{v}/x}{B} + (\Lam{x}{B}{\vec{t}_{1}}) \vec{w}\\
      &\lraneq \vec{t}_{1}\ansubst{\vec{v}/x}{B} + \vec{t}_{1}\ansubst{\vec{w}/x}{B}\\
      &=\sum_{i=1}^{n}(\alpha_i)\vec{t}_{1}[\vec{b}_{i}/x] + =\sum_{i=1}^{n}(\alpha_i)\vec{t}_{1}[\vec{b}_{i}/x]\\
      &\quad\text{where }\vec{v}\equiv\sum_{i=1}^{n}\alpha_i\vec{b}_{i},\vec{w}\equiv\sum_{i=1}^{n}\beta_i\vec{b}_{i}\text{ with }\vec{b}_{i}\in B 
    \end{align*}

    And we have that both terms are equivalent. The case for $B=\AbsBasis$ is similar.

    \item[$\LetP{x_1}{B_1}{x_2}{B_2}{(\alpha \vec{t})}{\vec{s}}\equiv\alpha(\LetP{x_1}{A}{x_2}{B}{\vec{t}}{\vec{s}})$:] If $\alpha=0$, then there is no reduction possible on the right-hand side. If there is an internal reduction on either $\vec{s}$ or $\vec{t}$, then we match the reductions on both sides of the equivalence with Lemma \ref{lem:SquigDiamond}.
    
    If $\vec{t}=\vec{v}$ and $\vec{s}\ansubst{\vec{v}/x_1\otimes x_2}{B_1\otimes B_2}$ is defined, then (we consider $B_1,B_2\neq\AbsBasis$):
    \begin{align*}
    \LetP{x_1}{B_1}{x_2}{B_2}{(\alpha \vec{v})}{&\vec{s}}\lraneq \vec{s}\ansubst{\alpha\vec{v}/x_1\otimes x_2}{B_1\otimes B_2}\\
    &=\sum_{i=1}^{n}\alpha\beta_i \vec{s}[\vec{v}_{i}/x_1][\vec{w}_{i}/x_2]\\ 
    &\text{where: }\vec{v}\equiv\sum_{i=1}^n\beta_i\Pair{\vec{v}_{i}}{\vec{w}_{i}} \text{ with } \vec{v}_{i}\in B_1, \vec{w}_{i}\in B_2
    \end{align*}
    On the other side;
    \begin{align*}
    \alpha(\LetP{x_1}{B_1}{x_2}{B_2}{\vec{v}}{&\vec{s}})\lraneq \alpha(\vec{s}\ansubst{\vec{v}/x_1\otimes x_2}{B_1\otimes B_2})\\
    &=\alpha\sum_{i=1}^{n}\beta_i \vec{s}[\vec{v}_{i}/x_1][\vec{w}_{i}/x_2]\\ 
    &\text{where: } \vec{v}\equiv\sum_{i=1}^n\beta_i\Pair{\vec{v}_{i}}{\vec{w}_{i}} \text{ with } \vec{v}_{i}\in B_1, \vec{w}_{i}\in B_2
    \end{align*}

    And we have that both terms are equivalent. The case for $B_1,B_2=\AbsBasis$ are similar.

    \item[\parbox{.75\linewidth}{\begin{align*}
      &\LetP{x_1}{B_1}{x_2}{B_2}{\vec{t}+\vec{s}}{\vec{r}}\equiv\\
      &(\LetP{x_1}{B_1}{x_2}{B_2}{\vec{t}}{\vec{r}}) +
      (\LetP{x_1}{B_2}{x_2}{B_2}{\vec{s}}{\vec{r}})
      \end{align*}}:]\hfill\\
      If there is an internal reduction on either $\vec{t}, \vec{s}$ or $\vec{r}$, then we match the reductions on both sides of the equivalence with Lemma \ref{lem:SquigDiamond}.

      If $\vec{t}=\vec{v}$ and $\vec{s}=\vec{w}$ with $\vec{r}\ansubst{\vec{v}/x_1\otimes x_2}{B_1\otimes B_2}$ and $\vec{r}\ansubst{\vec{v}/x_1\otimes x_2}{B_1\otimes B_2}$ defined, then (we consider $B_1,B_2\neq\AbsBasis$):
      \begin{align*}
      \LetP{x_1}{B_1}{x_2}{B_2}{\vec{v}+\vec{w}}{\vec{r}}&\lraneq\vec{r}\ansubst{\vec{v}+\vec{w}/x_1\otimes x_2}{B_1\otimes B_2}\\
      &=\sum_{i=1}^n (\alpha_i+\beta_i) \vec{r}[\vec{v}_{i}/x_1][\vec{w}_{i}/x_2]\\
      &\text{where: }\vec{v}\equiv\sum_{i=1}^n\alpha_i\Pair{\vec{v}_{i}}{\vec{w}_{i}} \text{ with } \vec{v}_{i}\in B_1, \vec{w}_{i}\in B_2\\
      &\text{and: }\vec{w}\equiv\sum_{i=1}^n\beta_i\Pair{\vec{v}_{i}}{\vec{w}_{i}} \text{ with } \vec{v}_{i}\in B_1, \vec{w}_{i}\in B_2
      \end{align*}

      On the other side:
      \begin{align*}
      &(\LetP{x_1}{B_1}{x_2}{B_2}{\vec{t}}{\vec{r}}) + (\LetP{x_1}{B_2}{x_2}{B_2}{\vec{s}}{\vec{r}})\\
      &\lraneq\vec{r}\ansubst{\vec{v}/x_1\otimes x_2}{B_1\otimes B_2} + (\LetP{x_1}{B_2}{x_2}{B_2}{\vec{s}}{\vec{r}})\\
      &\lraneq\vec{r}\ansubst{\vec{v}/x_1\otimes x_2}{B_1\otimes B_2} + \vec{r}\ansubst{\vec{w}/x_1\otimes x_2}{B_1\otimes B_2}\\
      &=\sum_{i=1}^n\alpha_i\vec{r}[\vec{v}_{i}/x_1][\vec{w}_{i}/x_2] + \sum_{i=1}^n \beta_i \vec{r}[\vec{v}_{i}/x_1][\vec{w}_{i}/x_2]\\
      &\text{where: }\vec{v}\equiv\sum_{i=1}^n\alpha_i\Pair{\vec{v}_{i}}{\vec{w}_{i}} \text{ with } \vec{v}_{i}\in B_1, \vec{w}_{i}\in B_2\\
      &\text{and: }\vec{w}\equiv\sum_{i=1}^n\beta_i\Pair{\vec{v}_{i}}{\vec{w}_{i}} \text{ with } \vec{v}_{i}\in B_1, \vec{w}_{i}\in B_2
      \end{align*}
        
      And we have that both terms are equivalent. The case for $B_1,B_2=\AbsBasis$ are similar.

    \item[\parbox{.55\linewidth}{\begin{align*}
      &\gencase{\alpha \vec{t}}{\vec{v}_{1}}{\vec{v}_{n}}{\vec{s}_{1}}{\vec{s}_{n}}\equiv\\
      &\alpha(\gencase{\vec{t}}{\vec{v}_{1}}{\vec{v}_{n}}{\vec{s}_{1}}{\vec{s}_{n}})
    \end{align*}}:] \hfill\\
    
    If $\alpha=0$, then there is no reduction possible on the right hand-side. If there are internal reductions on $\vec{t}$, then we match on both sides of the equivalence with Lemma \ref{lem:SquigDiamond}.
    
    If $\vec{t}=\vec{v}\equiv\sum_{i=1}^n\beta_i \vec{v}_{i}$. Then:
    \begin{align*}
      \gencase{\alpha \vec{t}}{\vec{v}_{1}}{\vec{v}_{n}}{\vec{s}_{1}}{\vec{s}_{n}}\lraneq \sum_{i=1}^n\alpha\beta_i \vec{s}_{i}
    \end{align*}
    On the other side:
    \begin{align*}
      \alpha(\gencase{\vec{t}}{\vec{v}_{1}}{\vec{v}_{n}}{\vec{s}_{1}}{\vec{s}_{n}})\lraneq \alpha\sum_{i=1}^n\beta_i \vec{s}_{i}
    \end{align*}
    And we have that both terms are equivalent.

    \item[\parbox{.55\linewidth}{\begin{align*}
      &\gencase{(\vec{t}+\vec{s})}{\vec{v}_{1}}{\vec{v}_{n}}{\vec{r}_{1}}{\vec{r}_{n}}\equiv\\ 
      &\gencase{\vec{t}}{\vec{v}_{1}}{\vec{v}_{n}}{\vec{r}_{1}}{\vec{r}_{n}}+\\
      &\gencase{\vec{s}}{\vec{v}_{1}}{\vec{v}_{n}}{\vec{r}_{1}}{\vec{r}_{n}}  
    \end{align*}}:]\hfill\\

    If there is an internal reduction on either $\vec{t}$ or $\vec{s}$, then we match the reductions on both sides of the equivalence with Lemma \ref{lem:SquigDiamond}.

    If $\vec{t}=\vec{v}\equiv\sum_{i=1}^n\alpha_i \vec{v}_{i}$, and $\vec{s}=\vec{w}\equiv\sum_{i=1}^n\beta_i \vec{v}_{i}$. Then:
    \begin{align*}
      \gencase{(\vec{t}+\vec{s})}{\vec{v}_{1}}{\vec{v}_{n}}{\vec{r}_{1}}{\vec{r}_{n}}\lraneq \sum_{i=1}^n (\alpha_i+\beta_i) \vec{r}_{i}
    \end{align*}
    On the other side:
    \begin{align*}
      &\gencase{\vec{t}}{\vec{v}_{1}}{\vec{v}_{n}}{\vec{r}_{1}}{\vec{r}_{n}}+\gencase{\vec{s}}{\vec{v}_{1}}{\vec{v}_{n}}{\vec{r}_{1}}{\vec{r}_{n}}\\
      &\lraneq \sum_{i=1}^n \alpha_i+ \vec{r}_{i} + \gencase{\vec{s}}{\vec{v}_{1}}{\vec{v}_{n}}{\vec{r}_{1}}{\vec{r}_{n}}\\
      &\lraneq \sum_{i=1}^n \alpha_i \vec{r}_{i} + \sum_{i=1}^n \beta_i \vec{r}_{i}
    \end{align*}

    And we have that both terms are equivalent.
    \qedhere
  \end{description}
\end{proof}


\section{Omitted proofs from Section~\ref{sec:model}}\label{sec:appendixC}

\begin{restatetheorem}[Restatement of Theorem~\ref{thm:SharpCharacterization}]
  \itshape
  The interpretation of a type~$\sharp A$ contains precisely the
  norm-$1$ linear combinations of values in~$\sem{A}$:
  \[
    \sem{\sharp A}
    = (\sem{A}^\bot)^\bot
    = \Span(\sem{A}) \cap \Sph.
  \]
\end{restatetheorem}
\begin{proof}
  Proof by double inclusion.
  \begin{description}
    \item[$\Span(\sem{A})\cap\Sph\subseteq (\sem{A}^\bot)^\bot$:] Let $\vec{v}\in\Span(\sem{A})\cap\Sph$. Then $\vec{v}$ is of the form $\sum_{i=1}^{n}\alpha_i \vec{v}_{i}$ with $\vec{v}_{i}\in\sem{A}$. Taking $\vec{w}\in\sem{A}^\bot$, we examine the inner product:
    \begin{align*}
    \scal{\vec{v}}{\vec{w}} &= \scal{\sum_{i=1}^{n}\alpha_i \vec{v}_{i}}{\vec{w}}\\
    &= \sum_{i=1}^{n}\overline{\alpha_i}\scal{\vec{v}_{i}}{\vec{w}}=0
    \end{align*}

    Then $\vec{v}\in(\sem{A}^\bot)^\bot$.

    \item[$(\sem{A}^\bot)^\bot\subseteq \Span(\sem{A})\cap\Sph$:] Reasoning by contradiction, we assume that there is a $\vec{v}\in(\sem{A}^\bot)^\bot$ such that $v\not\in\Span(\sem{A})\cap\Sph$. Since $\vec{v}\not\in\Span(\sem{A})$, $\vec{v}=\vec{w}_{1} + \vec{w}_{2}$ such that $\vec{w}_{1}\in\Span{\sem{A}}$ and $\vec{w}_{2}$ is a non-null vector which cannot be written as a linear combination of elements of $\sem{A}$. In other words, $\vec{w}_{2}\in\sem{A}^\bot$. Taking the inner product:
    \[
    \scal{\vec{v}}{\vec{w}_{2}} = \scal{\vec{w}_{1}+\vec{w}_{2}}{\vec{w}_{2}} = \|\vec{w}_{2}\|\neq 0
    \]
    Then $\vec{v}\not\in(\sem{A}^\bot)^\bot$. The contradiction stems from assuming $\vec{v}\not\in\Span{\sem{A}}\cap\Sph$.\qedhere
  \end{description}
\end{proof}

\begin{restatetheorem}[Restatement of Theorem~\ref{thm:IdempotentSharp}]
  \itshape
  The~$\sharp$ operator is idempotent; that is,
  $\sem{\sharp A} = \sem{\sharp(\sharp A)}$.
\end{restatetheorem}
\begin{proof}
  We want to prove that $(((\comp{\sem{A}})^\bot)^\bot)^\bot = (\comp{\sem{A}})^\bot$. For ease of reading, we will write $\comp[n]{A}$ for $n$ successive applications of the operation $\bot$.

  \begin{description}
    \item[$A\subseteq A^{\bot^2}$:] Let $\vec{v}\in A$. Then, for all $\vec{w}\in\comp{A}$, $\scal{\vec{v}}{\vec{w}} = 0$. Then $\vec{v}\in\comp[2]{A}$. With this we have $A\subseteq\comp[2]{A}$.
    
    \item[$A^{\bot^3}\subseteq \comp{A}$:] Let $\vec{u}\in \comp[3]{A}$. Then, for all $\vec{v}\in\comp[2]{A}$, $\scal{\vec u}{\vec v} = 0$. Since we have shown that $A\subseteq \comp[2]{A}$, we have that for all $\vec{w}\in A$, $\scal{\vec u}{\vec w} = 0$. Then $\vec u\in\comp{A}$. With this we have $\comp[3]{A}\subseteq \comp{A}$.
  \end{description}

  With these two inclusions we have that $\comp{A}=\comp[3]{A}$. So we conclude that: $\sem{\sharp(\sharp A)} = \comp[4]{A} = \comp[2]{A} = \sem{\sharp A}$ \qedhere
\end{proof}

\begin{restatetheorem}[Restatement of Theorem~\ref{prop:UnitaryTypes}]
  \itshape
  For every type~$A$, $\sem{A}\subseteq\Sph$.
\end{restatetheorem}

\begin{proof}
  Proof by induction on the shape of $A$. Since by definition, $\sem{\basis{X}}$, $\sem{A\Arr B}$ and $\sem{\sharp{A}}$ are built from values in $\Sph$ the only case we need to examine is $\sem{A\times B}$.
  
  Let $\vec v = \sum_{i=0}^{n} \alpha_i v_i \in\sem{A}$ and $\vec w = \sum_{j=0}^{m} \beta_j w_j$ where every $v_i$ are pairwise orthogonal, same for $w_j$. Then:
     
  \[(\vec v, \vec w) = \sum_{i=0}^{n} \sum_{j=0}^{m} \alpha_i\beta_j (v_i,w_j)\]
  
  So we have: 
  \[\|\Pair{\vec v}{\vec w}\| = \sqrt{\sum_{i=1}^n\sum_{j=1}^{m} |\alpha_i\beta_j|^2} = \sqrt{\sum_{i=1}^n |\alpha_i|^2 \sum_{j=1}^{m} |\beta_j|^2}\]

  Since both $\vec v\in\sem{A}$ and $\vec w\in\sem{B}$, by inductive hypothesis, we have that $\|\vec v\| = \| \vec w \| = 1$. Which is to say $\sum_{i=1}^{n} |\alpha_i|^2 = \sum_{j=1}^{m} |\beta_j| = 1$. So we conclude $\|\Pair{\vec{v}}{\vec{w}}\| = 1$.
  \qedhere
\end{proof}

\begin{restatelemma}[Restatement of Lemma~\ref{lem:BasesIso}]
  Let $X$ and $Y$ be orthonormal bases of the same finite
  dimension, and let $\Lam{x}{{X}}{\vec t}$ be a closed $\lambda$-abstraction.
  Then $\Lam{x}{{X}}{\vec t}\in\sem{\sharp\basis{X}\Arr\sharp\basis{Y}}$
  if and only if 
  for all $\vec{v}_{i},\vec{v}_{j}\in\sem{\basis{X}}$,
  there exist value combinations
  $\vec{w}_{i},\vec{w}_{j}\in\sem{\sharp\basis{Y}}$ such that,
  \[
    \vec{t}[\vec{v}_{i}/x]\eval\vec{w}_{i}
    \quad\text{and}\quad
    \vec{t}[\vec{v}_{j}/x]\eval\vec{w}_{j},
    \quad\text{with } 
    \vec{w}_{i}\perp\vec{w}_{j}\text{ whenever }i\neq j.
  \]
\end{restatelemma}
\begin{proof}
  \textit{The condition is necessary:} Suppose that $\Lam{x}{{X}}{\vec{t}_{k}}\in\sem{\sharp\basis{X}\Arr\sharp\basis{Y}}$, thus $\forall \vec{v}_{i}\in\sem{\sharp\basis{X}},\ \vec{t}\ansubst{\vec{v}_{i}/x}{X}\eval\vec{w}_{i}\in\sem{\sharp\basis{Y}}$. It remains to be seen that $\vec{w}_{i} \perp \vec{w}_{j}$ if $i\neq j$. For that, we consider $\alpha_i\in\C$ such that $\sum_{i=1}^n |\alpha_i|^2 = 1$. By linear application on the basis $X$ we observe that:
  \begin{align*}
    (\Lam{x}{{X}}{\vec{t}})(\sum_{i=1}^n \alpha_i \vec{v}_{i}) &\lra \vec t\ansubst{\sum_{i=1}^n \alpha_i \vec{v}_{i}/x}{X}
    = \sum_{i=1}^{n} \alpha_i \vec{t}[\vec{v}_{i}/x] 
    \eval \sum_{i=1}^n \alpha_i \vec{w}_{i}
  \end{align*}

  But since $\sum_{i=1}^n \alpha_i \vec{v}_{i}\in\sem{\sharp A}$, then $\sum_{i=1}^n \alpha_i \vec{w}_{i}\in\sem{\sharp B}$ too. Which implies $\|\sum_{i=1}^n \alpha_i \vec{w}_{i}\|=1$. Therefore:
  \begin{align*}
    1 = \|\sum_{i=1}^n \alpha_i \vec{w}_{i}\| &= \scal{\sum_{i=1}^n \alpha_i \vec{w}_{i}}{\sum_{j=1}^n \alpha_j \vec{w}_{j}}\\
    &=\sum_{i=1}^n |\alpha_i|^2 \scal{\vec{w}_{i}}{\vec{w}_{i} } + \sum_{i,j=1; i\neq j}^n \bar{\alpha_i}\alpha_j \scal{\vec{w}_{i}}{\vec{w}_{j}}\\
    &=\sum_{i=1}^n |\alpha_i|^2 \scal{\vec{w}_{i}}{\vec{w}_{i} } + \sum_{i,j=1; i<j}^n 2~\Rpart{\bar{\alpha_i}\alpha_j \scal{\vec{w}_{i}}{\vec{w}_{j}}}\\
    &=\sum_{i=1}^n |\alpha_i|^2 \|\vec{w}_{i}\|^2 + 2\sum_{i,j=1; i<j}^n \Rpart{\bar{\alpha_i}\alpha_j \scal{\vec{w}_{i}}{\vec{w}_{j}}}\\
    &=\sum_{i=1}^n |\alpha_i|^2 + 2\sum_{i,j=1; i<j}^n\Rpart{\bar{\alpha_i}\alpha_j \scal{\vec{w}_{i}}{\vec{w}_{j}}}\\
    &= 1 + 2\sum_{i,j=1; i<j}^n \Rpart{\bar{\alpha_i}\alpha_j \scal{\vec{w}_{i}}{\vec{w}_{j}}}
  \end{align*}

  And thus we are left with $\sum_{i,j=1; i<j}^n \Rpart{\bar{\alpha_i}\alpha_j \scal{\vec{w}_{i}}{\vec{w}_{j}}} = 0$. Taking $\alpha_{i'} = \alpha_{j'} = \frac{1}{\sqrt{2}}$ with $0$ for the rest of coefficients, we have $\Rpart{\scal{\vv{w_{i'}}}{\vv{w_{j'}}}} = 0$ for any two arbitrary $i'$ and $j'$. In the same way, taking $\alpha_{i'} = \frac{1}{\sqrt{2}}$ and $\alpha_{j'}=\frac{i}{\sqrt{2}}$ with $0$ for the rest of the coefficients, we have $\Ipart{\scal{\vv{w_{i'}}}{\vv{w_{j'}}}} = 0$ for any two arbitrary $i'$ and $j'$. Finally, we can conclude that $\scal{\vec{w}_{i}}{\vec{w}_{j}}=0$ if $i\neq j$.

  \textit{The condition is sufficient:} Suppose that there are $\vec{w}_{i}\in\sem{\sharp\basis{Y}}$ such that for every $\vec{v}_{i}\in\sem{\basis{X}}$:
  \[
    \vec t[\vec{v}_{i}/x] \eval \vec{w}_{i} \perp \vec{w}_{j} \lave \vec t[\vec{v}_{j}/x]\qquad \text{If } i\neq j
  \]
  Given any $\vec u\in\sem{\sharp\basis{X}}$ we have that $\vec u = \sum_{i=1}^n \alpha_i \vec{v}_{i}$ with $\sum_{i=1}^n |\alpha_i|^2 = 1$ and $\vec{v}_{i}\in\sem{\basis{X}}$. Then 
  \[
    (\Lam{x}{{X}}{\vec t}) \vec u \lra \vec{t}_{k}\ansubst{\vec u/x}{X}=\sum_{i=1}^{n}\alpha_i \vec{t}[\vec{v}_{i}/x]\eval\sum_{i=1}^n \alpha_i\vec{w}_{i}
  \]

  We have that for each $i$, $\vec{w}_{i}\in\sem{\sharp\basis{Y}}$. In order to show that $(\Lam{x}{A}{\vec t})\vec u\real\sharp\basis{Y}$ we still have to prove that $\|\sum_{i=1}^n \alpha_i \vec{w}_{i}\| = 1$
  \begin{align*}
    \|\sum_{i=1}^n \alpha_i \vec{w}_{i}\|^2 &= \scal{\sum_{i=1}^n \alpha_i \vec{w}_{i}}{\sum_{j=1}^n \alpha_j \vec{w}_{j}}\\
    &=\sum_{i=1}^n |\alpha_i|^2 \scal{\vec{w}_{i}}{\vec{w}_{i} } + \sum_{i,j=1; i\neq j}^n \bar{\alpha_i}\alpha_j \scal{\vec{w}_{i}}{\vec{w}_{j}}\\
    &=\sum_{i=1}^n |\alpha_i|^2 + 0\\
    &= 1
  \end{align*}

  Then $\sum_{i=1}^n \alpha_i \vec{w}_{i}\in\sem{\sharp(\sharp\basis{Y})}=\sem{\sharp\basis{Y}}$ by Theorem~\ref{thm:IdempotentSharp}. Since for every $\vec u\in\sem{\sharp A}$, $(\Lam{x}{A}{\vec t}) \vec u\real\sharp B$, we can conclude that $\Lam{x}{A}{\vec t}\in\sem{\sharp A\Arr\sharp B}$.\qedhere
\end{proof}

\begin{restatetheorem}[Restatement of Theorem~\ref{thm:characterisation}]
  Let $X$ and $Y$ be orthonormal bases of size~$2^n$.
  A closed $\lambda$-abstraction
  $\Lam{x}{{X}}{\vec{t}}\in\sem{\sharp\basis{X}\Arr\sharp\basis{Y}}$
  if and only if it represents a unitary operator
  $F:\C^{2^n}\to\C^{2^n}$.
\end{restatetheorem}

\begin{proof}
  \textit{Necessity.}
  Suppose that $\Lam{x}{{X}}{\vec{t}}\in\sem{\sharp\basis{X}\Arr\sharp\basis{Y}}$.
  Then, by Lemma~\ref{lem:BasesIso}, for every
  $\vec{v}_{i}\in\sem{\basis{X}}$ there exists
  $\vec{w}_{i}\in\sem{\sharp\basis{Y}}$ such that
  $\vec{t}[\vec{v}_{i}/x]\eval\vec{w}_{i}$ and
  $\vec{w}_{i}\perp\vec{w}_{j}$ whenever $i\neq j$.
  Let $F:\C^{2^n}\to\C^{2^n}$ be the operator defined by
  $F(\vec{v}_{i})=\vec{w}_{i}$.
  By linearity over~$X$, it follows that
  $\Lam{x}{{X}}{\vec{t}}$ represents~$F$.
  Then, $F$ is unitary since
  $\|\vec{w}_{i}\|_{\C^{2^n}}=\|\vec{w}_{j}\|_{\C^{2^n}}=1$ and
  $\scal{\vec{w}_{i}}{\vec{w}_{j}}_{\C^{2^n}}=0$.

  \smallskip
  \textit{Sufficiency.}
  Conversely, suppose that $\Lam{x}{{X}}{\vec{t}}$ represents a
  unitary operator $F:\C^{2^n}\to\C^{2^n}$.
  Then, for each $\vec{v}_{i}\in\sem{\basis{X}}$, there exists
  $\vec{w}_{i}\in\sem{\sharp\basis{Y}}$ such that
  $F(\vec{v}_{i}) = \vec{w}_{i}$ and
  $(\Lam{x}{{X}}{\vec{t}})\vec{v}_{i}\eval\vec{w}_{i}$.
  Hence,
  \(
    (\Lam{x}{{X}}{\vec{t}})\vec{v}_{i}
    \lraneq
    \vec{t}\ansubst{\vec{v}_{i}/x}{X}
    = \vec{t}[\vec{v}_{i}/x]
    \eval \vec{w}_{i}
    \in \sem{\sharp\basis{Y}}
  \),
  as $\|\vec{w}_{i}\|=\|F(\vec{v}_{i})\|_{\C^n}=1$. Since
  \(
    \scal{\vec{w}_{i}}{\vec{w}_{j}}
    = \scal{F(\vec{v}_{i})}{F(\vec{v}_{j})}_{\C^n}
    = 0
  \),
  from Lemma~\ref{lem:BasesIso} it follows that
  $\Lam{x}{{X}}{\vec{t}}\in\sem{\sharp\basis{X}\Arr\sharp\basis{Y}}$.
  \qedhere
\end{proof}


Before proving the soundness of the typing rules (Theorem~\ref{thm:TypingRulesValidity}), we need the following results.

\begin{theorem}\label{prop:InnerProdPairs} For all value combinations $\vec{v}_{1}, \vec{v}_{2}, \vec{w}_{1}, \vec{w}_{2}$ we have:
\[
\scal{\Pair{\vec{v}_{1}}{\vec{w}_{1}}}{\Pair{\vec{v}_{2}}{\vec{w}_{2}}} = \scal{\vec{v}_{1}}{\vec{v}_{2}}\scal{\vec{w}_{1}}{\vec{w}_{2}}
\]
\begin{proof}
    Let us write $\vec{v}_{1}=\sum_{i_1=1}^{n_1}\alpha_{i_1} v_{i_1}$, $\vec{v}_{2}=\sum_{i_2=1}^{n_2}\alpha'_{i_2} v_{i_2}$, $\vec{w}_{1}=\sum_{j_1=1}^{m_1}\beta_{j_1} w_{j_1}$ and $\vec{w}_{2}=\sum_{j_2=1}^{m_2}\beta'_{j_2} w_{j_2}$. Then we have:
    \begin{align*}
        &\scal{\Pair{\vec{v}_{1}}{\vec{w}_{1}}}{\Pair{\vec{v}_{2}}{\vec{w}_{2}}}\\
        &=\scal{\sum_{i_1=1}^{n_1}\sum_{j_1=1}^{m_1} \alpha_{i_1}\beta'_{j_1}\Pair{v_{i_1}}{w_{j_1}}}{\sum_{i_2=1}^{n_2}\sum_{j_2=1}^{m_2} \alpha_{i_2}\beta'_{j_2}\Pair{v_{i_2}}{w_{j_2}}}\\
        &=\sum_{i_1}^{n_1}\sum_{j_1}^{m_1}\sum_{i_2}^{n_2}\sum_{j_2}^{m_2} \overline{\alpha_{i_1}\beta_{j_1}} \alpha'_{i_2}\beta'_{j_2} \scal{\Pair{v_{i_1}}{w_{j_1}}}{\Pair{v_{i_2}}{w_{j_2}}}\\
        &=\sum_{i_1}^{n_1}\sum_{j_1}^{m_1}\sum_{i_2}^{n_2}\sum_{j_2}^{m_2} \overline{\alpha_{i_1}\beta_{j_1}} \alpha'_{i_2}\beta'_{j_2} \Kron{\Pair{v_{i_1}}{w_{j_1}}}{\Pair{v_{i_2}}{w_{j_2}}}\\
        &=\sum_{i_1}^{n_1}\sum_{j_1}^{m_1}\sum_{i_2}^{n_2}\sum_{j_2}^{m_2} \overline{\alpha_{i_1}\beta_{j_1}} \alpha'_{i_2}\beta'_{j_2} \Kron{v_{i_1}}{v_{i_2}}\Kron{w_{j_1}}{w_{j_2}}\\
        &=(\sum_{i_1}^{n_1}\sum_{j_1}^{m_1}\overline{\alpha_{i_1}}\alpha'_{i_2}\Kron{v_{i_1}}{v_{i_2}})(\sum_{i_2}^{n_2}\sum_{j_2}^{m_2} \overline{\beta_{j_1}} \beta'_{j_2} \Kron{w_{j_1}}{w_{j_2}})\\
        &=(\sum_{i_1}^{n_1}\sum_{j_1}^{m_1}\overline{\alpha_{i_1}}\alpha'_{i_2}\Pair{v_{i_1}}{v_{i_2}})(\sum_{i_2}^{n_2}\sum_{j_2}^{m_2} \overline{\beta_{j_1}} \beta'_{j_2} \Pair{w_{j_1}}{w_{j_2}})\\
        &=\scal{\vec{v}_{1}}{\vec{v}_{2}}\scal{\vec{w}_{1}}{\vec{w}_{2}}
	\tag*{\qedhere}
    \end{align*}
\end{proof}  

\end{theorem}

\begin{lemma}\label{lem:VecRewrite}%A5 en el paper de LICS
Given a type $A$, two vectors $\vec{u}_{1},\vec{u}_{2}\in\sem{\sharp A}$ and a scalar $\alpha\in\C$, there exists a vector $\vec{u}_{0}\in\sem{\sharp A}$ and a scalar $\lambda\in\C$ such that:
\[
\vec{u}_{1} + \alpha\vec{u}_{2} = \lambda \vec{u}_{0} 
\]
\end{lemma}
\begin{proof}
    Let $\lambda:=\|\vec{u}_{1}+\alpha\vec{u}_{2}\|$. When $\lambda\neq 0$, we take $\vec{u}_{0}=\frac{1}{\lambda}(\vec{u}_{1}+\alpha\vec{u}_{2})\in\sem{\sharp A}$, and we are done.

    When $\lambda=0$, we first observe that $\alpha\neq 0$ since it would mean that $\|\vec{u}_{1}\|=0$ which is absurd since $\|\vec{u}_{1}\|=1$. Moreover, since $\lambda=\|\vec{u}_{1}+\alpha\vec{u}_{2}\|=0$, we observe that all the coefficients of the combination $\vec{u}_{1}+\alpha\vec{u}_{2}$ are zeroes when written in canonical form which implies that:
    \[
    \vec{u}_{1}+\alpha\vec{u}_{2} = 0(\vec{u}_{1}+\alpha\vec{u}_{2}) = 0\vec{u}_{1}+0\vec{u}_{2}
    \]
    Using the triangular inequality we observe that:
    \begin{align*}
    0 &< 2|\alpha|\\
    &= \|2\alpha\vec{u}_{2}\|\\
    &\leq\|\vec{u}_{1}+\alpha\vec{u}_{2}\| + \|\vv{u_1 }+ (-\alpha)\vec{u}_{2}\|\\
    &= \|\vec{u}_{1}+(-\alpha)\vec{u}_{2}\|
    \end{align*}
    Hence $\lambda' := \|\vec{u}_{1}+(-\alpha)\vec{u}_{2}\|>0$. Taking $\vec{u}_{0}:= \frac{1}{\lambda'}(\vec{u}_{1}+ (-\alpha)\vec{u}_{2})\in\sem{\sharp A}$, we easily see that:
    \[
    \vec{u}_{1}+\alpha\vec{u}_{2} = 0\vec{u}_{1} + 0\vec{u}_{2} = 0(\frac{1}{\lambda'} (\vec{u}_{1} + (-\alpha) \vec{u}_{2})) = \lambda \vec{u}_{0}
    \tag*{\qedhere}
    \]
\end{proof}

\begin{theorem}[Polarization identity]\label{prop:Polarization} %A6
For all values $\vec{v}$ and $\vec{w}$ we have:
\[
  \scal{\vec{v}}{\vec{w}}=
  \frac{1}{4} (\|\vec{v}+\vec{w}\|^2 - \|\vec{v} + (-1) \vec{w}\|^2 - i\|\vec{v} + i\vec{w}\|^2 + i\|\vec{v}+ (-i)\vec{w}\|^2)
\]
\end{theorem}

\begin{lemma}\label{lem:InnerProdSingleVar} %A7
Given a valid typing judgement of the term $\TYP{\Delta,x_B:\sharp A}{\vec{s}}{C}$, a substitution $\sigma\in\sem{\Delta}$ and value combinations $\vec{u}_{1},\vec{u}_{2}\in\sem{\sharp A}$, there are value combinations $\vec{w}_{1}, \vec{w}_{2}\in\sem{C}$ such that:
\[
\begin{array}{c}
    \vec{s}\ansubst{\sigma}{}\ansubst{\vec{u}_{1}/x}{B_1}{\ansubst{\vec{v}_{1}/y}{B_2}}\eval\vec{w}_{1}\\
    \vec{s}\ansubst{\sigma}{}\ansubst{\vec{u}_{2}/x}{B_1}{\ansubst{\vec{v}_{2}/y}{B_2}}\eval\vec{w}_{2}
\end{array}
\]

And, $\scal{\vec{w}_{1}}{\vec{w}_{2}} = \scal{\vec{u}_{1}}{\vec{u}_{2}}$.
\end{lemma}

\begin{proof}
    From the validity of the judgement of the form $\TYP{\Delta, x_A:\sharp A}{\vec{s}}{C}$, a substitution $\sigma\in\sem{\Delta}$, and value combinations $\vec{w}_{1},\vec{w}_{2}\in\sem{C}$ such that $\vec{s}\ansubst{\sigma}{}\ansubst{\vec{u}_{1}/x}{A}\eval\vec{w}_{1}$ and $\vec{s}\ansubst{\sigma}{}\ansubst{\vec{u}_{2}/x}{A}\eval\vec{w}_{2}$. In particular, we have that $\|\vec{w}_{1}\| = \|\vec{w}_{2}\|=1$. Applying Lemma~\ref{lem:VecRewrite}~four times, we know there are vectors $\vv{u_{01}},\vv{u_{02}},\vv{u_{03}},\vv{u_{04}}\in\sem{\sharp A}$ and scalars $\lambda_1,\lambda_2,\lambda_3,\lambda_4$ such that:
    \begin{align*}
        \vec{u}_{1} + \vec{u}_{2} = \lambda_1 \vv{u_{01}} & \vec{u}_{1} + i \vec{u}_{2} = \lambda_3 \vv{u_{03}} \\
        \vec{u}_{1} + (-1) \vec{u}_{2} = \lambda_2 \vv{u_{02}} & \vec{u}_{1} + (-i) \vec{u}_{2} = \lambda_4 \vv{u_{04}} 
    \end{align*}

    From the validity of the judgement  $\TYP{\Delta, x_A:\sharp A}{\vec{s}}{C}$, we also know that there are value combinations $\vv{w_{01}},\vv{w_{02}},\vv{w_{03}},\vv{w_{04}}\in\sem{C}$ such that $\vec{s}\ansubst{\sigma}{}\ansubst{\vv{u_{0j}}}{}\eval\vv{w_{oj}}$ for all $f\in\{1\dotsb 4\}$. Combining the linearity of evaluation on the basis $A$ with the uniqueness of normal forms we deduce from what precedes that:
    \begin{align*}
        \vec{w}_{1} + \vec{w}_{2} = \lambda_1 \vv{w_{01}} & \vec{w}_{1} + i \vec{w}_{2} = \lambda_3 \vv{w_{03}} \\
        \vec{w}_{1} + (-1) \vec{w}_{2} = \lambda_2 \vv{w_{02}} & \vec{w}_{1} + (-i) \vec{w}_{2} = \lambda_4 \vv{w_{04}} 
    \end{align*}

    Using the polarization identity (Theorem~\ref{prop:Polarization}), we conclude that:
    \begin{align*}
        &\scal{\vec{w}_{1}}{\vec{w}_{2}}\\
        &= \frac{1}{4}(\|\vec{w}_{1}+\vec{w}_{2}\| - \|\vec{w}_{1} + (-1)\vec{w}_{2}\| - i \|\vec{v}_{1} + i \vec{v}_{2}\| + i \|\vec{v}_{1} + (-i) \vec{v}_{2}\|)\\
        &= \frac{1}{4}((\lambda_1)^2\|\vv{w_{01}}\| - (\lambda_2)^2\|\vv{w_{02}}\| - i (\lambda_)^2 \|\vv{w_{03}}\| + i (\lambda_)^2\|\vv{w_{04}}\|)\\
        &= \frac{1}{4}((\lambda_1)^2\|\vv{u_{01}}\| - (\lambda_2)^2\|\vv{u_{02}}\| - i (\lambda_)^2 \|\vv{u_{03}}\| + i (\lambda_)^2\|\vv{u_{04}}\|)\\
        &= \frac{1}{4}(\|\vec{u}_{1}+\vec{u}_{2}\| - \|\vec{u}_{1} + (-1)\vec{u}_{2}\| - i \|\vec{u}_{1} + i \vec{u}_{2}\| + i \|\vec{u}_{1} + (-i) \vec{u}_{2}\|)\\
        &=\scal{u_1}{u_2}
	\tag*{\qedhere}
    \end{align*}
\end{proof}

\begin{lemma}\label{lem:OrthogonalSubstitution} %A8
Given a valid typing judgement of the form $\TYP{\Delta, x_{B_1}:\sharp A_1, y_{B_2}: \sharp A_2}{\vec{s}}{C}$, a substitution $\sigma\in\sem{\Delta}$ and value combinations $\vec{u}_{1},\vec{u}_{2}\in\sem{\sharp A}$, there are value combinations $\vec{w}_{1},\vec{w}_{2}\in\sem{C}$ such that:
\[
\begin{array}{c}
    \vec{s}\ansubst{\sigma}{}\ansubst{\vec{u}_{1}/x}{B_1}{\ansubst{\vec{v}_{1}/y}{B_2}}\eval\vec{w}_{1}\\
    \vec{s}\ansubst{\sigma}{}\ansubst{\vec{u}_{2}/x}{B_1}{\ansubst{\vec{v}_{2}/y}{B_2}}\eval\vec{w}_{2}\\
\end{array}
\]
And, $\scal{\vec{w}_{1}}{\vec{w}_{2}} = 0$.
\end{lemma}

\begin{proof}
    From Lemma~\ref{lem:VecRewrite} we know that there are $\vec{u}_{0}\in\sem{\sharp A}, \vec{v}_{0}\in\sem{\sharp B}$ and $\lambda,\mu\in\C$ such that:
    \[
    \vec{u}_{2} + (-1) \vec{u}_{1} = \lambda\vec{u}_{0}\quad\text{and}\quad\vec{v}_{2} + (-1) \vec{v}_{1} = \mu \vec{v}_{0}
    \]
    For all $j,k\in\{0,1,2\}$, we have $\vec{s}\ansubst{\sigma}{}\ansubst{\vec{u}_{j}/x}{B_1}\ansubst{\vec{v}_{k}/y}{B_2}\eval\vv{w_{jk}}$. In particular, we can take $\vec{w}_{1}=\vv{w_{11}}$ and $\vec{w}_{2}=\vv{w_{22}}$. Now we observe that:
    \begin{enumerate}
        \item\label{A8:it1} $\vec{u}_{1}+\lambda\vec{u}_{0}= \vec{u}_{1} + \vec{u}_{2} + (-1) \vec{u}_{1}= \vec{u}_{2} + 0\vec{u}_{1}$, so that from linearity of substitution, linearity of evaluation and uniqueness of normal forms, we get:
        \[
        \begin{array}{c c}
            \begin{array}{c}
                \vv{w_{1k}} + \lambda\vv{w_{0k}} = \vv{w_{2k}} + 0 \vv{w_{1k}}\\
                \vv{w_{2k}} + (-\lambda)\vv{w_{0k}} = \vv{w_{1k}} + 0 \vv{w_{2k}}
            \end{array}&
            (\text{for all }k\in\{0,1,2\})
        \end{array}
        \]
        
        \item\label{A8:it2} $\vec{v}_{1}+\mu\vec{v}_{0}= \vec{v}_{1} + \vec{v}_{2} + (-1) \vec{v}_{1}= \vec{v}_{2} + 0\vec{v}_{1}$, so that from linearity of substitution, linearity of evaluation and uniqueness of normal forms, we get:
        \[
        \begin{array}{c c}
            \begin{array}{c}
                \vv{w_{j1}} + \mu\vv{w_{j0}} = \vv{w_{j2}} + 0 \vv{w_{j1}}\\
                \vv{w_{j2}} + (-\mu)\vv{w_{j0}} = \vv{w_{j1}} + 0 \vv{w_{j2}}
            \end{array}&
            (\text{for all }j\in\{0,1,2\})
        \end{array}
        \]
        
        \item\label{A8:it3} $\scal{\vec{u}_{1}}{\vec{u}_{2}}=0$, so that from Lemma~\ref{lem:InnerProdSingleVar}~we get $\scal{\vv{w_{1k}}}{\vv{w_{2k}}}=0$ (for all $k\in\{0,1,2\}$).
        
        \item\label{A8:it4} $\scal{\vec{v}_{1}}{\vec{v}_{2}}=0$, so that from Lemma~\ref{lem:InnerProdSingleVar}~we get $\scal{\vv{w_{j1}}}{\vv{w_{j2}}}=0$ (for all $j\in\{0,1,2\}$).
    \end{enumerate}

    From the above, we get:
    \begin{align*}
        \scal{\vec{w}_{1}}{\vec{w}_{2}} &= \scal{\vv{w_{11}}}{\vv{w_{22}}} = \scal{\vv{w_{11}}}{\vv{w_{22}}+0\vv{w_{12}}} & \\
        &=\scal{\vv{w_{11}}}{\vv{w_{12}}+ \lambda\vv{w_{02}}} & (\text{from Item~\ref{A8:it1}, } k=2)\\
        &=\scal{\vv{w_{11}}}{\vv{w_{12}}} + \lambda \scal{\vv{w_{11}}}{\vv{w_{02}}} &\\
        &= 0 + \lambda \scal{\vv{w_{11}}}{\vv{w_{02}}} & (\text{from Item~\ref{A8:it4}, } j=1)\\
        &= \lambda \scal{\vv{w_{11}} + 0\vv{w_{21}}}{\vv{w_{02}}} & \\
        &= \lambda \scal{\vv{w_{21}} + (-\lambda)\vv{w_{01}}}{\vv{w_{02}}} & (\text{from Item~\ref{A8:it1}, } k=1)\\
        &= \lambda \scal{\vv{w_{21}}}{\vv{w_{02}}} - |\lambda|^2 \scal{\vv{w_{01}}}{\vv{w_{02}}} & \\
        &= \lambda \scal{\vv{w_{21}}}{\vv{w_{02}}} - 0 & (\text{from Item~\ref{A8:it4}, } j=0)\\
        &=\scal{\vv{w_{21}}}{\vv{w_{22}}- \vv{w_{12}}} & \\
        &=\scal{\vv{w_{21}}}{\vv{22}} - \scal{\vv{w_{21}}}{\vv{w_12}} & \\
        &= 0 - \scal{\vv{w_{21}}}{\vv{w_{12}}} & (\text{from Item~\ref{A8:it4}, } j=2)
    \end{align*}

    Hence $\scal{\vec{w}_{1}}{\vec{w}_{2}} = \scal{\vv{w_{11}}}{\vv{w_{22}}} = - \scal{\vv{w_{21}}}{\vv{w_{12}}}$. Exchanging the indices in the previous reasoning, we also get 
    \[
    \scal{\vec{w}_{1}}{\vec{w}_{2}}=-\scal{\vv{w_{21}}}{\vv{w_{12}}}=-\scal{\vv{w_{12}}}{\vv{w_{21}}}
    \]
    So that we have:
    \[
        \scal{\vec{w}_{1}}{\vec{w}_{2}}=-\scal{\vv{w_{21}}}{\vv{w_{12}}}=-\overline{\scal{\vv{w_{21}}}{\vv{w_{12}}}}\in\R
    \]
    If we now replace $\vec{u}_{2}\in\sem{\sharp A}$ with $i\vec{u}_{2}\in\sem{\sharp A}$, the very same technique allows us to prove that $i\scal{\vec{w}_{1}}{\vec{w}_{2}}=\scal{\vec{w}_{1}}{i \vec{w}_{2}}\in\R$. Therefore, $\scal{\vec{w}_{1}}{\vec{w}_{2}}=0$.
\end{proof}

\begin{lemma}\label{lem:UnitPreserTens} %A9
Given a valid typing judgement of the form $\TYP{\Delta,x_{B_1}:\sharp A_1, y_{B_2}:\sharp A_2}{\vec{s}}{C}$, a substitution $\sigma\in\sem{\Delta}$, and value combinations $\vec{u}_{1},\vec{u}_{2}\in\sem{\sharp A}$ and $\vec{v}_{1},\vec{v}_{2}\in\sem{\sharp B}$, there are value combinations $\vec{w}_{1},\vec{w}_{2}\in\sem{C}$ such that:
\[
\begin{array}{c}
    \vec{s}\ansubst{\sigma}{}\ansubst{\vec{u}_{1}/x}{B_1}{\ansubst{\vec{v}_{1}/y}{B_2}}\eval\vec{w}_{1}\\
    \vec{s}\ansubst{\sigma}{}\ansubst{\vec{u}_{2}/x}{B_1}{\ansubst{\vec{v}_{2}/y}{B_2}}\eval\vec{w}_{2}\\
\end{array}
\]

And, $\scal{\vec{w}_{1}}{\vec{w}_{2}} = \scal{\vec{u}_{1}}{\vec{u}_{2}} \scal{\vec{v}_{1}}{\vec{v}_{2}}$.

\begin{proof}
    Let $\alpha=\scal{\vec{u}_{1}}{\vec{u}_{2}}$ and $\beta=\scal{\vec{v}_{1}}{\vec{v}_{2}}$. We observe that:
    \[
    \scal{\vec{u}_{1}}{\vec{u}_{2}+(-\alpha)\vec{u}_{1}} = \scal{\vec{u}_{1}}{\vec{u}_{2}} - \alpha \scal{\vec{u}_{1}}{\vec{u}_{1}} = \alpha - \alpha = 0
    \]
    And similarly that, $\scal{\vec{v}_{1}}{\vec{v}_{2}+ (-\beta) \vec{v}_{1}} = 0$. From Lemma~\ref{lem:VecRewrite}, we know that there are $\vec{u}_{0}\in\sem{\sharp A}$, $\vec{v}_{0}\in\sem{\sharp B}$ and $\lambda,\mu\in\C$ such that:
    \begin{align*}
        \vec{u}_{2} +(-\alpha)\vec{u}_{1} = \lambda\vec{u}_{0}& \text{ and } & \vec{v}_{2} + (-\beta)\vec{v}_{1} = \mu\vec{v}_{0} 
    \end{align*}
    For all $j,k\in\{0,1,2\}$, we have$\ansubst{\sigma}{}\ansubst{\vec{u}_{j}/x}{B_1}\ansubst{\vec{v}_{k}/y}{B_2}\in\sem{\Delta,x_{B_1}:\sharp A_1, y_{B_2}:\sharp A_2}$, hence there is $\vv{w_{jk}}\in\sem{C}$ such that:
    \[
    \vec{s}\ansubst{\sigma}{}\ansubst{u_j/x}{B_1}\ansubst{\vec{v}_{k}/y}{B_2}\eval\vv{w_{jk}}
    \]
    In particular, we can take $\vec{w}_{1}=\vv{w_{11}}$ and $\vec{w}_{2}=\vv{w_{22}}$. Now we observe that:
    \begin{enumerate}
        \item\label{A9:it1} $\lambda \vec{u}_{0} + \alpha\vec{u}_{1}=\vec{u}_{2} + (-\alpha) \vec{u}_{1} + \alpha \vec{u}_{1} = \vec{u}_{2} + 0 \vec{u}_{1}$, so that from the linearity of the substitution, linearity of evaluation and uniqueness of normal forms, we get:
        \[
        \lambda\vv{w_{0k}} + \alpha \vv{w_{1k}} = \vv{w_{2k}} + 0 \vv{w_{1k}} \qquad(\text{for all }k\in\{0,1,2\})
        \]
        
        \item\label{A9:it2} $\mu\vec{v}_{0} + \beta\vec{v}_{1}=\vec{v}_{2} + (-\beta) \vec{v}_{1} + \beta \vec{v}_{1} = \vec{v}_{2} + 0 \vec{v}_{1}$, so that from the linearity of the substitution, linearity of evaluation and uniqueness of normal forms, we get:
        \[
        \mu\vv{w_{j0}} + \beta \vv{w_{j1}} = \vv{w_{j2}} + 0 \vv{w_{j1}} \qquad(\text{for all }j\in\{0,1,2\})
        \]
        
        \item\label{A9:it3} $\scal{\vec{u}_{1}}{\lambda\vec{u}_{0}}=\scal{\vec{u}_{1}}{\vec{u}_{2}+(- \alpha)\vec{u}_{1}}=0$, so that from Lemma~\ref{lem:InnerProdSingleVar} we get:
        \[
        \scal{\vv{w_{1k}}}{\lambda\vv{w_{0k}}}= 0 \qquad(\text{for all }k\in\{0,1,2\})
        \]
        
        \item\label{A9:it4} $\scal{\vec{v}_{1}}{\mu\vec{v}_{0}}=\scal{\vec{v}_{1}}{\vec{v}_{2} + (-\beta)}\vec{v}_{1}=0$, so that from Lemma~\ref{lem:InnerProdSingleVar} we get:
        \[
        \scal{\vv{w_{j1}}}{\mu\vv{w_{j0}}}=0
        \]
                
        \item\label{A9:it5} $\scal{\vec{u}_{1}}{\lambda\vec{u}_{0}}=\scal{\vec{v}_{1}}{\mu\vec{v}_{0}}=0$ so that from Lemma~\ref{lem:OrthogonalSubstitution} we get:
        \[
        \scal{\vv{w_{11}}}{\lambda\mu\vv{w_{00}}}=0
        \]
        (Again the equality $\scal{\vv{w_{11}}}{\lambda\mu\vv{w_{00}}}$ is trivial when $\lambda=0$ or $\mu=0$. When $\lambda ,\mu\neq 0$ we deduce from the above that $\scal{\vec{u}_{1}}{\vec{u}_{0}}=\scal{\vec{v}_{1}}{\vec{v}_{0}}=0$, from which we get $\scal{\vv{w_{11}}}{\vv{w_{00}}}=0$ by Lemma~\ref{lem:OrthogonalSubstitution})
    \end{enumerate}

    From above, we get:
    \begin{align*}
        \vv{w_{22}} + 0\vv{w_{12}} + &0\vv{w_{01}} + 0\vv{w_{11}} \\
        &= \lambda\vv{w_{02}} + \alpha\vv{w_{12}} + 0\vv{w_{01}} + 0\vv{w_{11}} & (\text{from Item~\ref{A9:it1}}, k =1)\\
        &= \lambda(\vv{w_{02}}+0\vv{w_{01}}) + \alpha(\vv{w_{12}}+0\vv{w_{11}})&\\
        &= \lambda(\mu\vv{w_{00}} + \beta\vv{w_{01}}) + \alpha(\mu\vv{w_{01}}+\beta\vv{w_{11}}) & (\text{from Item~\ref{A9:it2}}, j=0,1)\\
        &= \lambda\mu\vv{w_{00}} + \lambda\beta\vv{w_{01}} + \alpha\mu\vv{w_{10}} + \alpha\beta\vv{w_{11}}
    \end{align*}
    Therefore:
    \begin{align*}
        &\scal{\vec{w}_{1}}{\vec{w}_{2}} \\
        &= \scal{\vv{w_{11}}}{\vv{w_{22}} + 0 \vv{w_{12}} + 0 \vv{w_{01}} + 0 \vv{w_{11}}}\\
        &= \scal{\vv{w_{11}}}{\lambda\mu\vv{w_{00}} + \lambda\beta\vv{w_{01}} + \alpha\mu\vv{w_{10}} + \alpha\beta\vv{w_{11}}}\\
        &=\scal{\vv{w_{11}}}{\lambda\mu\vv{w_{00}}} + \scal{\vv{w_{11}}}{\lambda\beta\vv{w_{01}}} + \scal{\vv{w_{11}}}{\alpha\mu\vv{w_{10}}} + \scal{\vv{w_{11}}}{\alpha\beta\vv{w_{11}}}\\
        &=\lambda\mu\scal{\vv{w_{11}}}{\vv{w_{00}}} + \lambda\beta\scal{\vv{w_{11}}}{\vv{w_{01}}} + \alpha\mu\scal{\vv{w_{11}}}{\vv{w_{10}}} + \alpha\beta\scal{\vv{w_{11}}}{\vv{w_{11}}}\\
        &= 0 + 0 + 0 + \alpha\beta = \scal{\vec{u}_{1}}{\vec{u}_{2}}\scal{\vec{v}_{1}}{\vec{v}_{2}}
    \end{align*}
  From Items~\ref{A9:it5}, \ref{A9:it3}, \ref{A9:it4} with $j=1$ and concluding with the definition of $\alpha$ and $\beta$.
\end{proof}
\end{lemma}

Now, we can restate and prove Theorem~\ref{thm:TypingRulesValidity}.
\begin{restatetheorem}[Restatement of Theorem~\ref{thm:TypingRulesValidity}]
  \itshape
  All the typing rules in Table~\ref{tab:TypingRules} are valid.
\end{restatetheorem}
\begin{proof}
    For each typing rule in Table~\ref{tab:TypingRules}~we have to show the typing judgement is valid starting from the premises:
    \begin{description}
    \item[Axiom] It is clear that $\sdom{x:A}\subseteq\{x\}=\dom{x:A}$. Moreover, given $\sigma\in\sem{x^B:A}$, we have $\sigma=\ansubst{\vec v/x}{B}$ for some $\vec{v}\in\sem{A}$. Therefore, $x\ansubst{\sigma}{}=x\ansubst{\vec v}{B}=\vec{v}\real A$.
    
    \item[UnitLam] If the hypothesis is valid, $\sdom{\Gamma,x^X:A}\subseteq \FV{\sum_{i=1}^{n}\alpha_i \vec{t}_{i}}\subseteq \dom{\Gamma,x^X:A}$. It follows that $\sdom{\Gamma}\subseteq \FV{\sum_{i=1}^{n}\alpha_i (\Lam{x}{X}{\vec{t}_{i}})}\subseteq \dom{\Gamma}$. Given $\sigma\in\sem{\Gamma}$, we want to show that $(\sum_{i=1}^{n}\alpha_i (\Lam{x}{X}{\vec{t}_{i}}))\ansubst{\sigma}{}\real A\Arr B$. Let $\vec v\in\sem{A}$, then:
    \begin{align*}
        (\sum_{i=1}^{n} \alpha_i(\Lam{x}{X}{\vec{t}_{i}}))\ansubst{\sigma}{} \vec v&= (\sum_{j=1}^{m} \beta_j (\sum_{i=1}^{n} \alpha_i (\Lam{x}{X}{\vec{t}_{i}}) [\sigma_i])) \vec{v} \\
        &= (\sum_{i=1}^{n}\sum_{j=1}^{m}\alpha_i\beta_j (\Lam{x}{X}{\vec{t}_{i}[\sigma_j]}))\vec v\\
        &\lra \sum_{i=1}^{n}\sum_{j=1}^{m}\alpha_i\beta_j \vec{t}_{i}[\sigma_j]\ansubst{\vec v/x}{X}\\
        &=\sum_{i=1}^{n}\alpha_i \vec{t}_{i}\ansubst{\sigma}{}\ansubst{\vec v/x}{X}\\
        &=(\sum_{i=1}^{n}\alpha_i \vec{t}_{i})\ansubst{\sigma}{}\ansubst{\vec v/x}{X}\qquad{\text{By Lemma~\ref{lem:distributiveSubstitution}}}
    \end{align*}
    
    Considering that $\ansubst{\sigma}{}\in\sem{\Gamma}$, then we have that $\ansubst{\sigma}{}\ansubst{\vec v/x}{X}\in\sem{\Gamma,x^X:A}$. Since we assume $\TYP{\Gamma, x^X:A}{\sum_{i=1}^{n}\alpha_i\vec{t}_{i}}{B}$, then $\vec{t}_{i}\ansubst{\sigma}{}\ansubst{\vec v/x}{X}\real B$. Finally, we can conclude that the combination: $\sum_{i=1}^{n}\alpha_i (\Lam{x}{X}{\vec{t}_{i}})\in\sem{A\Arr B}$.

    \item[App] If the hypotheses are valid, then:
    \begin{itemize}
        \item $\sdom{\Gamma}\subseteq \FV{\vec s}\subseteq \dom{\Gamma}$ and $\vec s \ansubst{\sigma_\Gamma}{}\Vdash A\Arr B\ \forall \sigma_\Gamma\in\sem{\Gamma}$.
        \item $\sdom{\Delta}\subseteq \FV{\vec t}\subseteq \dom{\Delta}$ and $\vec t\ansubst{\sigma_\Delta}{}\Vdash A\ \forall\sigma_\Delta\in\sem{\Delta}$.
    \end{itemize}
    
    From this, we can conclude that $\sdom{\Gamma,\Delta}\subseteq \FV{\vec s \vec t}\subseteq \dom{\Gamma,\Delta}$. Given $\sigma\in\sem{\Gamma,\Delta}$, we can observe that $\sigma=\sigma_\Gamma,\sigma_\Delta$ for some $\sigma_\Gamma\in\sem{\Gamma}$ and $\sigma_\Delta\in\sem{\Delta}$. Then we have:
    \begin{align*}
        (\vec{t}\vec{s})\ansubst{\sigma}{} &= (\vec{t}\vec{s})\ansubst{\sigma_\Gamma}{}\ansubst{\sigma_\Delta}{}\\
        &=(\sum_{i=i}^{n}\alpha_i (\vec{t}\vec{s})[\sigma_{\Gamma i}])\ansubst{\sigma_\Delta}{}\\
        &=\sum_{j=1}^{m} \beta_j (\sum_{i=1}^{n} \alpha_i (\vec{t} \vec{s})[\sigma_{\Gamma i}])[\sigma_{\Delta j}]\\
        &=\sum_{i=1}^{n}\sum_{j=1}^{m} \alpha_i \beta_j \vec{t}\,[\sigma_{\Gamma i}][\sigma_{\Delta j}] \vec{s}\,[\sigma_{\Gamma i}][\sigma_{\Delta j}]\\
        &=\sum_{i=1}^{n}\sum_{j=1}^{m} \alpha_i \beta_j \vec{t}\,[\sigma_{\Gamma i}]\vec{s}\,[\sigma_{\Delta j}]\\
        &\equiv (\sum_{i=1}^{n}\alpha_i\vec{t}[\sigma_{\Gamma i}])(\sum_{j=1}^{m} \beta_j \vec{s}[\sigma_{\Delta j}])\\
        &=\vec{t}\ansubst{\sigma_\Gamma}{} \vec{s}\ansubst{\sigma_\Delta}{}\\
        &\eval (e^{i\theta_{1}} \vec{v}) (e^{i\theta_{2}} \vec{w})\qquad\text{Where: } \vec{v}\in\sem{A\Arr B}, \vec{w}\in\sem{A}\\
        &\equiv e^{i\theta} (\vec{v} \vec{w})\qquad\text{with: }\theta=\theta_1 + \theta_2\\
        &\lraneq e^{i\theta}\vec r\qquad\text{where: } \vec{r}\real B
    \end{align*}
    
    Then we can conclude that $(\vec{t}\vec{s})\ansubst{\sigma}{}\real B$.

    \item[Pair] If the hypotheses are valid, then:

    \begin{itemize}
        \item $\sdom{\Gamma}\subseteq \FV{\vec s}\subseteq \dom{\Gamma}$ and $\vec s \ansubst{\sigma_\Gamma}{}\Vdash A\ \forall \sigma_\Gamma\in\sem{\Gamma}$.
        \item $\sdom{\Delta}\subseteq \FV{\vec t}\subseteq \dom{\Delta}$ and $\vec t\ansubst{\sigma_\Delta}{}\Vdash B\ \forall \sigma_\Delta\in\sem{\Delta}$.
    \end{itemize}
    
    From this, we can conclude that $\sdom{\Gamma,\Delta}\subseteq \FV{(\vec s, \vec t)}\subseteq \dom{\Gamma,\Delta}$. Given $\sigma\in\sem{\Gamma,\Delta}$, we can observe that $\sigma=\sigma_\Gamma,\sigma_\Delta$ for some  $\sigma_\Gamma\in\sem{\Gamma}$ and $\sigma_\Delta\in\sem{\Delta}$. Then we have:
    \begin{align*}
        \Pair{\vec{t}}{\vec{s}}\ansubst{\sigma}{} &= \Pair{\vec{t}}{\vec{s}}\ansubst{\sigma_\Gamma}{}\ansubst{\sigma_\Delta}{}\\
        &=\sum_{j=1}^{m} \beta_j (\sum_{i=1}^{n} \alpha_i \Pair{\vec{t}}{\vec{s}}[\sigma_{\Gamma i}])[\sigma_{\Delta j}]\\
        &\equiv\sum_{i=1}^{n}\sum_{j=1}^{m} \alpha_i \beta_j \Pair{\vec{t}\,[\sigma_{\Gamma i}][\sigma_{\Delta j}]}{\vec{s}\,[\sigma_{\Gamma i}][\sigma_{\Delta j}]}\\
        &=\sum_{i=1}^{n}\sum_{j=1}^{m} \alpha_i \beta_j \Pair{\vec{t}\,[\sigma_{\Gamma i}]}{\vec{s}\,[\sigma_{\Delta j}]}\\
        &=\Pair{\sum_{i=1}^{n} \alpha_i \vec{t}\, [\sigma_{\Gamma i}]}{\sum_{j=1}^{m} \beta_j \vec{s}\, [\sigma_{\Delta j}]}\\
        &=\Pair{\vec{t}\ansubst{\sigma_\Gamma}{}}{\vec{s}\ansubst{\sigma_\Delta}{}}\\
        &\eval \Pair{e^{i\theta_1}\vec v}{e^{i\theta_2}\vec w}\qquad\text{where: }\vec{v}\in\sem{A}, \vec{w}\in\sem{B}\\
        &= e^{i\theta} \Pair{\vec{v}}{\vec{w}}\qquad\text{where: }\theta=\theta_1 + \theta_2
    \end{align*}
    
    From this we can conclude that $\Pair{\vec t}{\vec{s}}\ansubst{\sigma}{}\real A\times B$.
    
    \item[LetPair] If the hypotheses are valid, then:
    \begin{itemize}
        \item $\sdom{\Gamma}\subseteq \FV{\vec t} \subseteq \dom{\Gamma}$ and $\vec t \ansubst{\sigma_\Gamma}{}\Vdash A\times B\ \forall \sigma_\Gamma\in\sem\Gamma$
        \item $\sdom{\Delta, {x}^{X}:A, {y}^{Y}:B}\subseteq\FV{\vec s}$
        \item $\FV{\vec{s}}\subseteq \dom{\Delta, {x}^{X}:A, {y}^{Y}:B}$
        \item $\vec s \ansubst{\sigma_\Delta}{}\Vdash C\ \forall \sigma_\Delta\in\sem{\Delta, {x}^{X}:A, {y}^{Y}:B}$
    \end{itemize}
    From this, we can conclude that:
    \begin{itemize}
        \item $\sdom{\Gamma,\Delta}\subseteq\FV{\LetP{x}{X}{y}{Y}{\vec{s}}{\vec{t}}}$
        \item $\FV{\LetP{x}{X}{y}{Y}{\vec{s}}{\vec{t}}}\subseteq\dom{\Gamma,\Delta}$
    \end{itemize}
    
    Given $\sigma\in\sem{\Gamma,\Delta}$, we have that $\ansubst{\sigma}{}=\ansubst{\sigma_\Gamma}{},\ansubst{\sigma_\Delta}{}$ for some $\sigma_\Gamma\in\sem\Gamma$ and $\sigma_\Delta\in\sem\Delta$. Then we have:
    \begin{align*}
        (&\LetP{x}{X}{y}{Y}{\vec{t}}{\vec{s}})\ansubst{\sigma}{} = \\
        &(\LetP{x}{X}{y}{Y}{\vec{t}}{\vec{s}})\ansubst{\sigma_\Gamma}{}\ansubst{\sigma_\Delta}{}\\
        &= (\sum_{i=1}^{n}\alpha_i(\LetP{x}{X}{y}{Y}{\vec{t}}{\vec{s}})[\sigma_{\Gamma i}])\ansubst{\sigma_{\Delta j}}{}\\
        &\equiv (\LetP{x}{X}{y}{Y}{\sum_{i=1}^{n}\alpha_i[\sigma_{\Gamma i}]\vec{t}}{\vec{s}})\ansubst{\sigma_\Delta}{}\\
        &= (\LetP{x}{X}{y}{Y}{\vec{t}\ansubst{\sigma_\Gamma}{}}{\vec{s}})\ansubst{\sigma_\Delta}{}\\
        &\eval (\LetP{x}{X}{y}{Y}{e^{i\theta} \Pair{\vec{v}}{\vec{w}}}{\vec{s}})\ansubst{\sigma_\Delta}{}\\
        &\hspace*{4cm}{\text{Where: }}\vec{v}\in\sem{A},\vec{w}\in\sem{B}\\
        &\lra e^{i\theta_1} (\vec{s}\ansubst{\sigma_\Delta}{}\ansubst{\Pair{\vec{v}}{\vec{w}}/x\otimes y}{X\otimes Y})\\
        &= e^{i\theta_1} (\vec{s}\ansubst{\sigma_\Delta}{}\ansubst{\vec{v}/x}{X}\ansubst{\vec{w}/y}{Y})\\
        &\eval e^{i\theta_1}  (e^{i\theta_2}  \vec{u})\qquad\text{where: }\vec{u}\in\sem{C}\\
        &\equiv e^{i\theta}  \vec{u}\qquad\text{where: }\theta=\theta_1 + \theta_2
    \end{align*}
    
    Since $\ansubst{\sigma_\Delta}{}\ansubst{\vec{v}/x}{X}\ansubst{\vec{w}/y}{Y}\in\sem{\Delta,{x}^{X}:A,{y}^Y:B}$, then we can conclude that $(\LetP{x}{X}{y}{Y}{\vec{t}}{\vec{s}})\ansubst{\sigma}{}\real C$.

    \item[LetTens] If the hypotheses are valid then:
    \begin{itemize}
        \item $\sdom{\Gamma}\subseteq \FV{\vec t} \subseteq \dom{\Gamma}$ and $\vec t \ansubst{\sigma}{}\Vdash \sharp(A\times B)\ \forall \sigma\in\sem\Gamma$
        \item $\sdom{\Delta, {x}^{X}:\sharp A, {y}^{Y}:\sharp B}\subseteq \FV{\vec s}$
        \item $\subseteq \dom{\Delta, {x}^{X}:\sharp A, {y}^{Y}:\sharp B}$
        \item $\vec s \ansubst{\sigma}{}\Vdash \sharp C\ \forall \sigma\in\sem{\Delta, {x}^{X}:\sharp A, {y}^{Y}:\sharp B}$
    \end{itemize}
    
    From this we can conclude that:
    \begin{itemize}
        \item $\sdom{\Gamma,\Delta}\subseteq\FV{\LetP{x}{X}{y}{Y|}{\vec{t}}{\vec{s}}}$
        \item $\FV{\LetP{x}{X}{y}{Y}{\vec{t}}{\vec{s}}}\subseteq\dom{\Gamma,\Delta}$
    \end{itemize}
    
    Given $\sigma\in\sem{\Gamma,\Delta}$, we have that $\ansubst{\sigma}{}=\ansubst{\sigma_\Gamma}{},\ansubst{\sigma_\Delta}{}$ for some $\sigma_\Gamma\in\sem\Gamma$ and $\sigma_\Delta\in\sem\Delta$. Using the first hypothesis we have that, $\vec t\ansubst{\sigma_\Gamma}{}\real \sharp(A\times B)$, from Theorem~\ref{thm:SharpCharacterization} we have that:
    
    \[\vec t\ansubst{\sigma_\Gamma}{}\eval e^{i\theta_1}\vec{u}=e^{i\theta_1}(\sum_{k=1}^{l} \gamma_k \Pair{\vec{v}_{k}}{\vec{u}_{k}})\] 
    
    With:
    \begin{itemize}
        \item $\sum_{k=1}^{l} |\gamma_k|^2 = 1$
        \item $\forall k,\ \vec{v}_{k}\in\sem{A},\ \vec{u}_{k}\in\sem{B}$
        \item $\forall k\neq l, \scal{\Pair{\vec{v}_{k}}{\vec{u}_{k}}}{\Pair{\vec{v}_{l}}{\vec{u}_{l}}}= 0$
    \end{itemize}
    
    Then:
    \begin{align*}
        (&\LetP{x}{X}{y}{Y}{\vec{t}}{\vec{s}})\ansubst{\sigma}{} \\
        &= \LetP{x}{X}{y}{Y}{\vec{t}}{\vec{s}}\ansubst{\sigma_\Gamma}{}\ansubst{\sigma_\Delta}{}\\
        &=(\sum_{i=1}^{n}\alpha_i\LetP{x}{X}{y}{Y}{\vec{t}}{\vec{s}}\ [\sigma_{\Gamma i}])\ansubst{\sigma_{\Delta}}{}\\
        &\equiv (\LetP{x}{X}{y}{Y}{\sum_{i=1}^{n}\alpha_i\vec{t}\ [\sigma_{\Gamma i}]}{\vec{s}})\ansubst{\sigma_\Delta}{}\\
        &=(\LetP{x}{X}{y}{Y}{\vec{t}\ansubst{\sigma_\Gamma}{}}{\vec{s}})\ansubst{\sigma_\Delta}{}\\
        &\eval(\LetP{x}{X}{y}{Y}{e^{i\theta_1}\vec{u}}{\vec{s}})\ansubst{\sigma_\Delta}{}\\
        &\lra e^{i\theta_1}(\vec{s}\ansubst{\sigma_\Delta}{}\ansubst{\vec{u}/x\otimes y}{X\otimes Y})\\
        &=e^{i\theta_1} (\sum_{k=1}^{l}\gamma_k\vec{s}\ansubst{\sigma_\Delta}{}\ansubst{\vec{v}_{k}/x}{X}\ansubst{\vec{u}_{k}/y}{Y})\\
        &\eval e^{i\theta_1} (\sum_{k=1}^{l}\gamma_k e^{i\rho_k} \vec{w}_{k})\quad\text{where: }\vec{w}_{k}\in\sem{C}
    \end{align*}

    It remains to be seen that the term has norm-$1$, $\|\sum_{k=1}^{l}\gamma_k e^{i\rho_k} \vec{w}_{k}\|=1$. For that, we observe:
    \begin{align*}
        \|&\sum_{k=1}^{l}\gamma_k e^{i\rho_k} \vec{w}_{k}\| \\
        &= \scal{\sum_{k=1}^{l}\alpha_i e^{i\rho_k} \vec{w}_{k}}{\sum_{k'=1}^{l}\gamma_{k'} e^{i\rho_{k'}} \vv{w_{k'}}}\\
        &= \sum_{k=1}^{l}\sum_{k'}^{l}\overline{\gamma_k e^{i\rho_k}}\  \gamma_{k'} e^{i\rho_{k'}}\scal{\vec{w}_{k}}{\vv{w_{k'}}}\\
        &=\sum_{k=1}^{l}\sum_{k'=1}^{l}\overline{\gamma_k e^{i\rho_k}}\ \gamma_{k'} e^{i\rho_{k'}} \scal{\vec{v}_{k}}{\vv{v_{k'}}}\scal{\vec{u}_{k}}{\vv{u_{k'}}}\quad(\text{from Lemma~\ref{lem:UnitPreserTens}})\\
        &= \sum_{k=1}^{k}\sum_{k'=1}^{l}\overline{\gamma_k e^{i\rho_k}}\  \gamma_{k'} e^{i\rho_{k'}} \scal{\Pair{\vec{u}_{k}}{\vec{v}_{k}}}{\Pair{\vv{u_{k'}}}{\vv{v_{k'}}}}\quad(\text{from Theorem~\ref{prop:InnerProdPairs}})\\
        &=\sum_{k=1}^n \overline{\gamma_k e^{i\rho_k}}\ \gamma_k e^{i\rho_k} \scal{\Pair{\vec{v}_{k}}{\vec{u}_{k}}}{\Pair{\vec{v}_{k}}{\vec{u}_{k}}} \\
        & \quad + \sum_{k,k'=1; k\neq k'}^n \overline{\gamma_k e^{i\rho_k}}\  \gamma_{k'} e^{i\rho_{k'}} \scal{\Pair{\vec{v}_{k}}{\vec{u}_{k}}}{\Pair{\vv{v_{k'}}}{\vv{u_{k'}}}}\\
        &= \sum_{k=1}^n \overline{\gamma_k e^{i\rho_k}}\ \gamma_k e^{i\rho_k} + 0 \\
        &= \sum_{k=1}^{l} |\gamma_k|^2 |e^{i\rho_k}|^2 = 1
    \end{align*}

    Then $\sum_{k=1}^{l}\gamma_k e^{i\rho_k} \vec{w}_{k}\in\sem{\sharp C}$. Finally, we can conclude that: $(\LetP{x}{X}{y}{Y}{\vec{t}}{\vec{s}})\ansubst{\sigma}{}\real{\sharp C}$

    \item[Case] If the hypotheses are valid then:
    \begin{itemize}
        \item $\sdom{\Gamma}\subseteq \FV{\vec{t}}\subseteq \dom{\Gamma}$
        \item For every $\sigma_\Gamma\in\sem{\Gamma}$, $\vec{t}\ansubst{\sigma_\Gamma}{}\real\genbasis{\vec{v}_{i}}{i=1}{n}$
        \item For every $i\in\{0,\dotsb ,n\}, \sdom{\Delta}\subseteq \FV{\vec{s}_{i}}\subseteq \dom{\Delta}$
        \item For every $i\in\{0,\dotsb ,n\}, \sigma_\Delta\in\sem{\Delta}$, $\vec{s}_{i}\ansubst{\sigma_\Delta}{}\real A$
    \end{itemize}

    From this we can conclude that:

    \begin{itemize}
        \item $\sdom{\Gamma,\Delta}\subseteq \FV{\gencase{\vec{t}}{\vec{v}_{1}}{\vec {v_n}}{\vec{s}_{1}}{\vec{s}_{n}}}$
        \item $\FV{\gencase{\vec{t}}{\vec{v}_{1}}{\vec {v_n}}{\vec{s}_{1}}{\vec{s}_{n}}}\subseteq \dom{\Gamma,\Delta}$
    \end{itemize}
    
    Then, given $\sigma\in\sem{\Gamma,\Delta}$, we have that $\ansubst{\sigma}{}=\ansubst{\sigma_\Gamma}{}\ansubst{\sigma_\Delta}{}$ for some $\sigma_\Gamma\in\sem{\Gamma}$ and $\sigma_\Delta\in\sem{\Delta}$. Using the first hypothesis we have that, $\vec{t}\ansubst{\sigma_\Gamma}{}\eval e^{i\theta_1}\vec{v}_{k}$ for some $k\in\{1,\dotsb ,n\}$. From the second hypothesis we have that $\vec{s}_{i}\ansubst{\sigma_\Delta}{}\eval e^{i\rho_i}\vec{u}_{i}\in\sem{A}$ for $i\in\{1,\dotsb , n\}$. Therefore:

    \begin{align*}
        (&\gencase{\vec{t}}{\vec{v}_{1}}{\vec{v}_{n}}{\vec{s}_{1}}{\vec{s}_{n}})\ansubst{\sigma}{}\\ 
        &= (\gencase{\vec{t}}{\vec{v}_{1}}{\vec{v}_{n}}{\vec{s}_{1}}{\vec{s}_{n}})\ansubst{\sigma_\Gamma}{}\ansubst{\sigma_\Delta}{}\\
        &= (\sum_{i=1}^{n}\alpha_i \gencase{\vec{t}[\sigma_{\Gamma i}]}{\vec{v}_{1}}{\vec{v}_{n}}{\vec{s}_{1}}{\vec{s}_{n}})\ansubst{\sigma_\Delta}{} \\
        &\equiv (\gencase{\sum_{i=1}^{n} \alpha_i \vec{t}[\sigma_{\Gamma i}]}{\vec{v}_{1}}{\vec{v}_{n}}{\vec{s}_{1}}{\vec{s}_{n}})\ansubst{\sigma_\Delta}{}\\
        &=(\gencase{\vec{t}\ansubst{\sigma_\Gamma}{}}{\vec{v}_{1}}{\vec{v}_{n}}{\vec{s}_{1}}{\vec{s}_{n}})\ansubst{\sigma_\Delta}{}\\
        &\eval(\gencase{e^{i\theta_1}\vec{v}_{k}}{\vec{v}_{1}}{\vec{v}_{n}}{\vec{s}_{1}}{\vec{s}_{n}})\ansubst{\sigma_\Delta}{}\\
        &\lraneq e^{i\theta_1}(\vec{s}_{k}\ansubst{\sigma_\Delta}{})\\
        &\eval e^{i\theta_1}(e^{i\rho_k} \vec{u}_{k})\qquad\text{Where: }\vec{u}_{k}\in\sem{A}\\
        &\equiv e^{i\theta} \vec{u}_{k}\qquad\text{With: }\theta=\theta_1 +\theta_2
    \end{align*}
    
    Since we pose no restriction on $k$, we can conclude that: $(\gencase{\vec{t}}{\vec{v}_{1}}{\vec{v}_{n}}{\vec{s}_{1}}{\vec{s}_{n}})\ansubst{\sigma}{}\real A$

    \item[UnitCase] If the hypotheses are valid, then:
    \begin{itemize}
        \item $\sdom{\Gamma}\subseteq \FV{\vec{t}}\subseteq \dom{\Gamma}$
        \item For every $\sigma_\Gamma\in\sem{\Gamma}$, $\vec{t}\ansubst{\sigma_\Gamma}{}\real\sharp\genbasis{\vec{v}_{i}}{i=1}{n}$
        \item For every $i\in\{0,\dotsb ,n\}$, $\sdom{\Delta}\subseteq \FV{\vec{s}_{i}}\subseteq \dom{\Delta}$
        \item For every $i\in\{0,\dotsb ,n\}, \sigma_\Delta\in\sem{\Delta}$, $\vec{s}_{i}\ansubst{\sigma_\Delta}{}\real A$
    \end{itemize}
    
    From this we can conclude that:
    
    \begin{itemize}
        \item $\sdom{\Gamma,\Delta}\subseteq \FV{\gencase{\vec{t}}{\vec{v}_{1}}{\vec{v}_{n}}{\vec{s}_{1}}{\vec{s}_{n}}}$
        \item $\FV{\gencase{\vec{t}}{\vec{v}_{1}}{\vec{v}_{n}}{\vec{s}_{1}}{\vec{s}_{n}}}\subseteq \dom{\Gamma,\Delta}$
    \end{itemize}
    
    Then, given $\sigma\in\sem{\Gamma,\Delta}$, we have that $\ansubst{\sigma}{}=\ansubst{\sigma_\Gamma}{}\ansubst{\sigma_\Delta}{}$ for some $\sigma_\Gamma\in\sem{\Gamma}$ and $\sigma_\Delta\in\sem{\Delta}$. Using the first hypothesis we have that, $\vec{t}\ansubst{\sigma_\Gamma}{}\real\sharp\genbasis{\vec{v}_{i}}{i=1}{n}$, then $\vec{t}\ansubst{\sigma_\Gamma}{}\eval e^{i\theta_1} \vec{u}\equiv e^{i\theta_1} (\sum_{i=1}^{n}\beta_i \vec{v}_{i})$ where $\sum_{i=1}^{n}|\beta_i|^2=1$. From the second hypothesis we have that $\vec{s}_{i}\ansubst{\sigma_\Delta}{}\eval e^{i\rho_i} \vec{u}_{i}\in\sem{A}$ for $i\in\{1,\dotsb ,n\}$ and $u_i\perp u_j$ if $i\neq j$. Therefore:

    \begin{align*}
        (&\gencase{\vec{t}}{\vec{v}_{1}}{\vec{v}_{n}}{\vec{s}_{1}}{\vec{s}_{n}})\ansubst{\sigma}{}\\ 
        &= (\gencase{\vec{t}}{\vec{v}_{1}}{\vec{v}_{n}}{\vec{s}_{1}}{\vec{s}_{n}})\ansubst{\sigma_\Gamma}{}\ansubst{\sigma_\Delta}{}\\
        &=(\sum_{i=1}^{n}\alpha_i \gencase{\vec{t}[\sigma_{\Gamma i}]}{\vec{v}_{1}}{\vec{v}_{n}}{\vec{s}_{1}}{\vec{s}_{n}})\ansubst{\sigma_\Delta}{} \\
        &\equiv (\gencase{\sum_{i=1}^{n} \alpha_i \vec{t}[\sigma_{\Gamma i}]}{\vec {v_1}}{\vec{v}_{n}}{\vec{s}_{1}}{\vec{s}_{n}})\ansubst{\sigma_\Delta}{}\\
        &=(\gencase{\vec{t}\ansubst{\sigma_\Gamma}{}}{\vec{v}_{1}}{\vec{v}_{n}}{\vec{s}_{1}}{\vec{s}_{n}})\ansubst{\sigma_\Delta}{}\\
        &\eval(\gencase{e^{i\theta_1} \vec{u}}{\vec{v}}{\vec{w}}{\vec{s}_{1}}{\vec{s}_{2}})\ansubst{\sigma_\Delta}{}\\
        &\lra e^{i\theta_1} (\sum_{i=1}^{n}\beta_i s_i)\ansubst{\sigma_\Delta}{}\\
        &= e^{i\theta_1} (\sum_{j=1}^{n}\delta_j (\sum_{i=1}^{n}\beta_i \vec{s}_{i})[\sigma_{\Delta j}])\\
        &\equiv e^{i\theta_1} (\sum_{i,j=1}^{n}\beta_i\delta_j\vec{s}_{i}[\sigma_{\Delta j}])\\
        &= e^{i\theta_1} (\sum_{i=1}^{n}\beta_i \vec{s}_{i}\ansubst{\sigma_\Delta}{})\\
        &\eval e^{i\theta_1} (\sum_{i=1}^{n}\beta_i e^{i\rho_i} \vec{u}_{i})
    \end{align*}
    
    It remains to be seen that: $\|\sum_{i=1}^{n}\beta_i e^{i\rho_i} \vec{u}_{i}\|=1$:
    \begin{align*}
        \|\sum_{i=1}^{n}\beta_i e^{i\rho_i} \vec{u}_{i}\| &= \scal{\sum_{i=1}^{n}\beta_i e^{i\rho_i} \vec{u}_{i}}{\sum_{i=1}^{n}\beta_i e^{i\rho_i} \vec{u}_{i}}\\
        &= \sum_{i,j=1}^{n}\overline{\beta_i e^{i\rho_i}}\beta_j e^{i\rho_j} \scal{\vec{u}_{i}}{\vec{u}_{j}}\\
        &= \sum_{i=1}^{n}\overline{\beta_i e^{i\rho_i}}\beta_i e^{i\rho_i} \scal{\vec{u}_{i}}{\vec{u}_{i}}\\
        &\qquad + \sum_{i,j=1; i\neq j}^{n}\overline{\beta_i e^{i\rho_i}}\beta_j e^{i\rho_j} \scal{\vec{u}_{i}}{\vec{u}_{j}}\\
        &= \sum_{i=1}^{n}|\beta_i|^2 |e^{i\rho_i}|^2  + 0\\
        &= \sum_{i=1}^{n}|\beta_i|^2 = 1
    \end{align*}

    Then we can conclude that $\sum_{i=1}^{n}\beta_i e^{i\rho_i}\vec{u}_{i}\in\sem{\sharp A}$ and finally: $(\gencase{\vec{t}}{\vec{v}_{1}}{\vec{v}_{n}}{\vec{s}_{1}}{\vec{s}_{n}})\ansubst{\sigma}{}\real\sharp A $
  
    \item[Sum] If the hypothesis is valid then for every $i$, $\sdom{\Gamma}\subseteq\FV{\vec{t}_{i}}\subseteq\dom{\Gamma}$.
    
    From this we can conclude that $\sdom{\Gamma}\subseteq\sum_{i=1}^{n}\alpha_i \vec{t}_{i}\subseteq\dom{\Gamma}$. Given $\sigma\in\sem{\Gamma}$, we have for every $i$, $\vec{t}_{i}\ansubst{\sigma}{}\eval e^{i\rho_i} \vec{v}_{i}$ where $\vec{v}_{i}\in\sem{A}$. Moreover, for every $i\neq j$, $\vec{v}_{i}\perp\vec{v}_{j}$ and $\sum_{i=1}^{n}|\alpha_i|^2=1$. Then:
    \begin{align*}
    (\sum_{i=1}^{n}\alpha_i\vec{t}_{i})\ansubst{\sigma}{} 
    &= \sum_{j=1}^{m}\beta_j(\sum_{i=1}^{n}\alpha_i \vec{t}_{i})[\sigma_j]\\
    &\equiv \sum_{i=1}^{n} \alpha_i \sum_{j=1}^{m} \beta_j \vec{t}_{i}[\sigma_j]\\
    &=\sum_{i=1}^{n} \alpha_i \vec{t}_{i}\ansubst{\sigma}{}\\
    &\eval \sum_{i=1}^{n} \alpha_i e^{i\rho_i} \vec{v}_{i}\\
    \end{align*}

    It remains to be seen that $\|\sum_{i=1}^{n} \alpha_i e^{i\rho_i} \vec{v}_{i}\|=1$:
    \begin{align*}
    &\|\sum_{i=1}^{n} \alpha_i e^{i\rho_i} \vec{v}_{i}\| \\
    &=\scal{\sum_{i=1}^{n} \alpha_i e^{i\rho_i}\vec{v}_{i}}{\sum_{i=1}^{n} \alpha_i e^{i\rho_i} \vec{v}_{i}}\\
    &= \sum_{i=i}^{n}\sum_{j=1}^{n} \overline{\alpha_i e^{i\rho_i}}\alpha_j e^{i\rho_j} \scal{\vec{v}_{i}}{\vec{v}_{j}}\\
    &=\sum_{i=1}^{n} \overline{\alpha_i e^{i\rho_i}}\alpha_i e^{i\rho_i} \scal{\vec{v}_{i}}{\vec{v}_{i}} + \sum_{\substack{i,j=1\\i\neq j}}^{n} \overline{\alpha_i e^{i\rho_i}}\alpha_j e^{i\rho_j} \scal{\vec{v}_{i}}{\vec{v}_{j}}\\
    &=\sum_{i=1}^{n}|\alpha_i|^2 |e^{i\rho_i}|^2+ 0\\
    &=\sum_{i=1}^{n}|\alpha_i|^2 = 1\\
    \end{align*}

    Then we can conclude that $\sum_{i=1}^{n}\alpha_i e^{i\rho_i}\vec{v}_{i}\in\sem{\sharp A}$ and finally $(\sum_{i=1}^{n}\alpha_i\vec{t}_{i})\ansubst{\sigma}{}\real\sharp A$.

    \item[Contr] If the hypothesis is valid, we have that $\sdom{\Gamma, x^X:\basis{X}, y^X:\basis{X}}\subseteq\FV{\vec{t}}\subseteq\dom{\Gamma,x^X:\basis{X}, y^X:\basis{X}}$ and given $\sigma\in\sem{\Gamma, x^X:\basis{X}, y^X:\basis{X}}$, then $\vec{t}\ansubst{\sigma}{}\in\sem{B}$. Then, we have that $\sdom{\Gamma, x^X:\basis{X}, y^X:\basis{X}}=\sdom{\Gamma, x^X:\basis{X}}$. Therefore:
    
    \[
    \sdom{\Gamma, x^X:\basis{X}}\subseteq\FV{\vec{t}}[x/y]\subseteq\dom{\Gamma, x^X:\basis{X}}
    \]

    Given $\sigma\in\sem{\Gamma, x^X:\basis{X}}$, we observe that $\ansubst{\sigma}{}=\ansubst{\vec{v}/x}{X}\ansubst{\sigma_\Gamma}{}$ with $\sigma_\Gamma\in\sem{\Gamma}$ and $\vec{v}\in\sem{\basis{X}}$. Since $\vec{v}\in\sem{\basis{X}}$, we know that $\vec{t}[\vec v/z] =\vec{t}\ansubst{\vec{v}/z}{X}$ for any variable $z$. Then we have:
    \begin{align*}
        \vec{t}[x/y]\ansubst{\sigma}{} &= \vec{t}[x/y]\ansubst{\vec{v}/x}{X}\ansubst{\sigma_\Gamma}{}\\
        &=\vec{t}[x/y][\vec{v}/x]\ansubst{\sigma_\Gamma}{}\\
        &=\vec{t}[\vec{v}/y][\vec{v}/x]\ansubst{\sigma_\Gamma}{}\\
        &=\vec{t}\ansubst{\vec{v}/y}{X}\ansubst{\vec{v}/x}{X}\ansubst{\sigma_\Gamma}{}    
    \end{align*}
    
    Since $\ansubst{\vec{v}/y}{X}\ansubst{\vec{v}/x}{X}\ansubst{\sigma}{}\in\sem{\Gamma, x^X:\basis{X}, y^X:\basis{X}}$, we get: $\vec{t}\ansubst{\vec{v}/y}{X}\ansubst{\vec{v}/x}{X}\\\ansubst{\sigma_\Gamma}{}\eval e^{i\theta}\vec{w}\in\sem{B}$
    Then we can finally conclude that $\vec{t}[x/y]\ansubst{\sigma}{}\real B$.

    \item[Weak] Given $\sigma\in\sem{\Gamma,x^X:\basis{X}}$, we observe that $\ansubst{\sigma}=\ansubst{\sigma_\Gamma}{}\ansubst{\vec{v}/x}{X}$ for some $\sigma_\Gamma\in\sem{\Gamma}$ and $\vec{v}\in\sem{\basis{X}}$. Using the first hypothesis, we know that $\vec{t}\ansubst{\sigma_\Gamma}{}\eval e^{i\theta}\vec{w}$ where $\vec{w}\in\sem{B}$. Then we have:
    \[
    \vec{t}\ansubst{\sigma}{}=\vec{t}\ansubst{\sigma_\Gamma}{}\ansubst{\vec{v}/x}{X}\eval e^{i\theta}\vec{w}\ansubst{\vec{v}/x}{X}
    \]
    Since $\vec{v}\in\sem{\basis{X}}$, $\vec{w}\ansubst{\vec{v}/x}{X}=\vec{w}[\vec{v}/x]=\vec{w}$ and $\vec{w}\in\sem{B}$, then we can finally conclude that $\vec{t}\ansubst{\sigma}{}\real B$.
    
    \item[Sub] Trivial since the set of realizers of ${A}$ is included in the set of realizers of ${B}$. 

    \item[Equiv] It follows from definition and the fact that the reduction commutes with the congruence relation.
    
    \item[Phase] It follows from the definition of type realizers.
      \qedhere
    \end{description}
\end{proof}

\begin{lemma}[Substitution]\label{lem:Substitution}
  Let $\Gamma$, $\Delta$ be contexts, $A$ and $B$ types, and $X$ an orthonormal basis.
  If
  \(
    \TYP{\Gamma,\,x^{X}\!:\!A}{\vec{t}}{B}
    \quad\text{and}\quad
    \TYP{\Delta}{\vec{v}}{A}
  \) can be derived using the set of rules in Table~\ref{tab:TypingRules},
  and the substitution $\vec{t}\ansubst{\vec{v}/x}{X}$ is defined,
  then
  \(
    \TYP{\Gamma,\,\Delta}{\vec{t}\ansubst{\vec{v}/x}{X}}{B}
  \) can also be derived by the same set of rules.
\end{lemma}
\begin{proof}
  By induction on $\vec{t}$.
  \begin{description}
    \item[$x\not\in\FV{\vec{t}}$:] Then $x\not\in\sdom{\Gamma,\,x^{X}\!:\!A}$. By a straightforward generation lemma we have that $\TYP{\Gamma}{\vec{t}}{B}$, and we can derive $\TYP{\Gamma,\Delta}{\vec{t}}{B}$ via the $\textsc{Weak}$ rule. Notice that if $x\not\in\FV{\vec{t}}$, every variable in $\Delta$ has a non-linear type and $\vec{t}\ansubst{\vec{v}/x}{X}=\vec{t}$.
    
    \item[$\vec{t}=x$:] Then every variable in $\Gamma$ has a non-linear type and $A\leq B$. This means we can derive $\TYP{\Gamma,\Delta}{\vec{v}}{B}$, by rules $\textsc{Weak}$ and $\textsc{Sub}$.
    
    \item[$\vec{t}=(\Lam{y}{Y}{\vec{s}})$:] Then $\TYP{\Gamma,\,y^{Y}\!:\!C,\,x^{X}\!:\!A}{\vec{s}}{D}$, with $C\Arr D\leq B$. By induction hypothesis, $\TYP{\Gamma,\,y^{Y}\!:\!C}{\vec{s}\ansubst{\vec{v}/x}{X}}{D}$. This means we can derive $\TYP{\Gamma,\Delta}{(\Lam{y}{Y}{\vec{s}})\ansubst{\vec{v}/X}{X}}{B}$, by rules $\textsc{Sub}$ and $\textsc{UnitLam}$.
    
    \item[$\vec{t}=\vec{s}_{1}\,\vec{s}_{2}$:] Then we have that $\TYP{\Gamma_1}{\vec{s}_{1}}{C\Arr D}$, and $\TYP{\Gamma_2}{\vec{s}_{2}}{C}$ with $D\leq B$ and $\Gamma_1,\Gamma_2 = \Gamma,x^{X}\!:\!A$. We consider the case where $\Gamma_1 = \Gamma_1',\,x^{X}\!:\!A$. Then, by inductive hypothesis $\TYP{\Gamma_1'\Delta}{\vec{s}_{1}\ansubst{\vec{v}/x}{X}}{C\Arr D}$. This means we can derive $\TYP{\Gamma_1',\Gamma_2,\Delta}{(\vec{s}_{1}\,\vec{s}_{2})\ansubst{\vec{v}/x}{X}}{B}$, by rules $\textsc{Sub}$ and $\textsc{App}$. The case where $\Gamma_2 = \Gamma_2',\,x^{X}\!:\!A$ is analogous.
    
    \item[$\vec{t}=\Pair{s_1}{s_2}$:] Then we have that $\TYP{\Gamma_1}{s_1}{C}$, and $\TYP{\Gamma_2}{s_2}{D}$ with $C\times D\leq B$ and $\Gamma_1,\,\Gamma_2 = \Gamma,x^{X}\!:\!A$. We consider the case where $\Gamma_1 = \Gamma_1',\,x^{X}\!:\!A$. Then, by inductive hypothesis $\TYP{\Gamma_1',\,\Delta}{s_1}{C}$. This means we can derive $\TYP{\Gamma_1',\,\Gamma_2,\,\Delta}{\Pair{s_1}{s_2}\ansubst{\vec{v}/x}{X}}{B}$, by rules $\textsc{Sub}$ and $\textsc{Pair}$. The case where $\Gamma_2 = \Gamma_2',\,x^{X}\!:\!A$ is analogous.
    
    \item[$\vec{t}=\LetP{y}{Y}{z}{Z}{\vec{s}_{1}}{\vec{s}_{2}}$:] Then we have two possibilities:
      \begin{enumerate}
        \item\label{case:subs_letp_1} Either, $\TYP{\Gamma_1}{\vec{s}_{1}}{C\times D}$, and $\TYP{\Gamma_2,\,y^{Y}\!:\!C,\,z^{Z}\!:\!D}{\vec{s}_{2}}{E}$ with $E\leq B$.
        \item\label{case:subs_letp_2} Or, $\TYP{\Gamma_1}{\vec{s}_{1}}{\sharp(C\times D)}$, and $\TYP{\Gamma_2,\,y^{Y}\!:\!\sharp{C},\,z^{Z}\!:\!\sharp{D}}{\vec{s}_{2}}{E}$ with $\sharp E\leq B$.
      \end{enumerate}
      In either way, we consider the case where $\Gamma_1 = \Gamma_1',\,x^{X}\!:\!A$. Then, by inductive hypothesis $\TYP{\Gamma_1'\,\Delta}{\vec{s}_{1}\ansubst{\vec{v}/x}{}}{C\times D}$ in case \ref{case:subs_letp_1} ($\sharp (C\times D)$ in case \ref{case:subs_letp_2}). This means we can derive $\TYP{\Gamma_1',\,\Gamma_2,\Delta}{(\LetP{y}{Y}{z}{Z}{\vec{s}_{1}}{\vec{s}_{2}})\ansubst{\vec{v}/x}{X}}{B}$ by rules $\textsc{Sub}$ and $\textsc{LetPair}$ (or, $\textsc{LetTens}$). The case where $\Gamma_2 = \Gamma_2',\,x^{X}\!:\!A$ is analogous.

    \item[$\vec{t}=\gencase{\vec{s}}{\vec{w}_{1}}{\vec{w}_{n}}{\vec{r}_{1}}{\vec{r}_{n}}$:] Then we have two possibilities:
    \begin{enumerate}
      \item\label{case:subs_case_1} Either, $\TYP{\Gamma_1}{\vec{s}_{1}}{\basis{\{\vec{v}_{i}\}_{i=1}^n}}$, and for all $i\in\{0,\dotsb,n\}$, $\TYP{\Gamma_2}{\vec{s}_{i}}{C}$ with $C\leq B$.
      \item\label{case:subs_case_2} Or, $\TYP{\Gamma_1}{\vec{s}_{1}}{\sharp\basis{\{\vec{v}_{i}\}_{i=1}^n}}$, and for all $i,j\in\{0,\dotsb,n\}$, with $i\neq j$ $\SORTH{\Gamma_2}{\vec{s}_{i}}{\vec{s}_{j}}{C}$ with $\sharp C\leq B$.
    \end{enumerate}
    In either way, we consider the case where $\Gamma_1 = \Gamma_1',\,x^{X}\!:\!A$. Then, by inductive hypothesis $\TYP{\Gamma_1'\,\Delta}{\vec{s}_{1}\ansubst{\vec{v}/x}{}}{\basis{\{\vec{v}_{i}\}_{i=1}^n}}$ in case \ref{case:subs_case_1} ($\sharp \basis{\{\vec{v}_{i}\}_{i=1}^n}$ in case \ref{case:subs_case_2}). This means we can derive $\TYP{\Gamma_1',\,\Gamma_2,\Delta}{(\gencase{\vec{s}}{\vec{w}_{1}}{\vec{w}_{n}}{\vec{r}_{1}}{\vec{r}_{n}})\ansubst{\vec{v}/x}{X}}{B}$ by rules $\textsc{Sub}$ and $\textsc{Case}$ (or, $\textsc{UnitCase}$). Since orthogonality is preserved by substitutions in $\sem{\Gamma_2}$, the case where $\Gamma_2 = \Gamma_2',\,x^{X}\!:\!A$ is analogous.

    \item[$\vec{t}=\sum_{i=1}^{n}\alpha_i \vec{s}_{i}$:] Then we have two posibilities:
    \begin{enumerate}
      \item\label{case:subs_sum_1} For all $i\in\{0,\dotsb, n\}, \vec{s}_{i}=(\Lam{y}{Y}{\vec{r}_{i}})$, and:\\
      $\TYP{\Gamma,\,x^{X}\!:\!A}{\sum_{i=1}^{n}\alpha_i(\Lam{y}{Y}{\vec{r}_{i}})}{C\Arr D}$ with $C\Arr D\leq B$.
      \item\label{case:subs_sum_2} For all $i\in\{0,\dotsb, n\}$, with $i\neq j$, $\SORTH{\Gamma,\,x^{X}\!:\!A}{\vec{s}_{i}}{\vec{s}_{j}}{C}$, $\sum_{i=1}^{n}|\alpha_i|^2=1$, and $\sharp C\leq B$.
    \end{enumerate}
    
    In case \ref{case:subs_sum_1}, by inductive hypothesis we have that:\\
    $\TYP{\Gamma,\,\Delta,\,y^{Y}\!:\!C}{\sum_{i=1}^{n}\alpha_i \vec{r}_{i}\ansubst{\vec{v}/x}{X}}{D}$. This means we can derive \\$\TYP{\Gamma,\Delta}{\sum_{i=1}^{n}\alpha_i (\Lam{y}{Y}{\vec{r}_{i}})\ansubst{\vec{v}/x}{X}}{B}$ by rules $\textsc{Sub}$ and $\textsc{UnitLam}$.

    In case \ref{case:subs_sum_2}, by inductive hypothesis we have that for all $i\in\{0,\dotsb,n\}$, and $i\neq j$ $\SORTH{\Gamma,\Delta}{\vec{s}_{i}\ansubst{\vec{v}/x}{X}}{\vec{s}_{j}\ansubst{\vec{v}/x}{X}}{C}$. Since orthogonality is preserved by substitutions in $\sem{\Gamma,\Delta}$, this means we can derive $\TYP{\Gamma,\Delta}{\sum_{i=1}^n\alpha_i \vec{s}_{i}\ansubst{\vec{v}/x}{X}}{B}$ by rules $\textsc{Sub}$ and $\textsc{Sum}$.

    \item[$\vec{t}=\alpha\vec{s}$:] Then $\alpha=e^{i\theta}$, and $\TYP{\Gamma,\,x^{X}\!:\!A}{\vec{s}}{C}$ with $C\leq B$. By induction hypothesis, $\TYP{\Gamma}{\vec{s}\ansubst{\vec{v}/x}{X}}{C}$. This means we can derive $\TYP{\Gamma,\Delta}{e^{i\theta}\vec{s}\ansubst{\vec{v}/X}{X}}{\\B}$, by rules $\textsc{Sub}$ and $\textsc{Phase}$.
      \qedhere
  \end{description}
\end{proof}

\begin{restatetheorem}[Restatement of Theorem~\ref{thm:SubjectReduction}]
  If $\TYP{\Gamma}{\vec{t}}{A}$ can be derived using the set of rules in
  Table~\ref{tab:TypingRules} and $\vec{t}\to\vec{u}$, then
  $\TYP{\Gamma}{\vec{u}}{A}$ can also be derived by the same set of rules.
\end{restatetheorem}
\begin{proof}
  We proceed by induction on the derivation of the elementary reduction
  $\lraneq$. The congruence closure to obtain $\lra$ is handled trivially via
  the \textsc{Equiv} rule. We only give the basis cases as the inductive cases (the contextual cases) are straightforward.
  \begin{itemize}
    \item Let $(\Lam{x}{X}{\vec t})\vec v \lraneq \vec t\ansubst{\vec v/x}{X}$:
      Assume $\TYP{\Gamma,\Delta}{(\Lam{x}{X}{\vec t})\,\vec v}{B}$. By
      \textsc{App}, we have $\TYP{\Gamma}{\Lam{x}{X}{\vec t}}{A\Arr B}$ and
      $\TYP{\Delta}{\vec v}{A}$.  From $\TYP{\Gamma}{\Lam{x}{X}{\vec t}}{A\Arr B}$
      and \textsc{UnitLam}, we get $\TYP{\Gamma,x^{X}:A}{\vec t}{B}$. By
      Lemma~\ref{lem:Substitution},
      using $\TYP{\Delta}{\vec v}{A}$ and the fact that $\vec t\ansubst{\vv
      v/x}{X}$ is defined, we conclude $\TYP{\Gamma,\Delta}{\vec t\ansubst{\vv
      v/x}{X}}{B}$, as required.

    \item Let $\LetP{x}{X}{y}{Y}{\vec v}{\vec t} \lraneq \vec t\ansubst{\vv
      v/x\otimes y}{X\otimes Y}$: Assume
      $\TYP{\Gamma,\Delta}{\LetP{x}{X}{y}{Y}{\vec v}{\vec t}}{C}$.  By
      \textsc{LetPair} we have $\TYP{\Gamma}{\vec v}{A\times B}$ and
      $\TYP{\Delta,\,x^{X}:A,\,y^{Y}:B}{\vec t}{C}$.  Decompose $\vec v$ (in
      $X\otimes Y$) as required by the definition of $\ansubst{\vec v/x\otimes
      y}{X\otimes Y}$; by Lemma~\ref{lem:Substitution}
      applied twice
      (first for $x$, then for $y$), we conclude $\TYP{\Gamma,\Delta}{\vv
      t\ansubst{\vec v/x\otimes y}{X\otimes Y}}{C}$.

    \item Let $\gencase{\vec{v}_{k}}{\vec{v}_{1}}{\vec{v}_{n}}{\vec{t}_{1}}{\vec{t}_{n}}
      \lraneq \vec{t}_{k}$: Assume
      $\TYP{\Gamma,\Delta}{\gencase{\vec{v}_{k}}{\vec{v}_{1}}{\vec{v}_{n}}{\vec{t}_{1}}{\vec{t}_{n}}}{A}$.
      By \textsc{Case} we have $\TYP{\Gamma}{{\vec{v}_{k}}}{\basis{\{\vv
      v_i\}_{i=1}^n}}$ and $\TYP{\Delta}{\vec t_i}{A}$ for all $i$.  Thus, we are
      done.

    \item Let
      $\gencase{\sum_{i=1}^n\alpha_i\vec{v}_{i}}{\vec{v}_{1}}{\vec{v}_{n}}{\vec{t}_{1}}{\vec{t}_{n}}
      \lraneq \sum_{i=1}^n\alpha_i\vec{t}_{i}$: \\ Assume 
      $\TYP{\Gamma,\Delta}{\gencase{\sum_{i=1}^n\alpha_i\vec{v}_{i}}{\vec{v}_{1}}{\vec{v}_{n}}{\vec{t}_{1}}{\vec{t}_{n}}}{A}$.
      By \textsc{UnitCase} we have $\TYP{\Gamma}{\vec t}{\sharp\basis{\{\vv
      v_i\}}}$ and, for all $i\neq j$, $\SORTH{\Delta}{\vec t_i}{\vec t_j}{A}$.
      Then the reduct $\sum_{i=1}^{n}\alpha_i \vec{t}_{i}$ is typed by \textsc{Sum}
      as $\sharp A$ (using the orthogonality premises and the normalisation
      condition ensured by the semantics of $\sharp$), hence
      $\TYP{\Gamma,\Delta}{\sum_{i=1}^{n}\alpha_i \vec{t}_{i}}{\sharp A}$.
      \qedhere
  \end{itemize}
\end{proof}

\section{Gates for the Deutsch's algorithm example of Section~\ref{subsec:deutsch}}\label{sec:appendixD}
\subsection{The four possible oracles \texorpdfstring{$U_f$}{Uf} in the computational basis \texorpdfstring{$\B$}{Z}}
\begin{align*}
  O_{\mathrm{const}\,0} &:= \Lam{x}{\B}{\Lam{y}{\B}{\Pair{x}{y}}}\\
  O_{\mathrm{const}\,1} &:= \Lam{x}{\B}{\Lam{y}{\B}{\Pair{x}{(\pauliX{y})}}}\\
  O_{\mathrm{id}}       &:= \Lam{x}{\B}{\Lam{y}{\B}{\cnot{x}{y}}}\\
  O_{\mathrm{flip}}     &:= \Lam{x}{\B}{\Lam{y}{\B}{\cnot{x}{(\pauliX{y})}}}
\end{align*}

\subsection{Gates in the Hadamard basis \texorpdfstring{$\XB$}{X}}
\begin{align*}
  \pauliZXB &:= \Lam{x}{\XB}{\case{x}{\ket{+}}{\ket{-}}{\ket{-}}{\ket{+}}}\\
  \pauliXXB &:= \Lam{x}{\XB}{\case{x}{\ket{+}}{\ket{-}}{\ket{+}}{-1\,\ket{-}}}\\
  \cnotXB &:= \Lam{x}{\XB}{\Lam{y}{\XB}{
      \case{y}{\ket{+}}{\ket{-}}
      {\Pair{x}{\ket{+}}}
      {\Pair{\pauliZXB{x}}{\ket{-}}}
  }}.
\end{align*}

\subsection{The four oracles rewritten in the \texorpdfstring{$\XB$}{X} basis}
\begin{align*}
  O_{\mathrm{const}\,0}^{\XB} &:= \Lam{x}{\XB}{\Lam{y}{\XB}{\Pair{x}{y}}}\\
  O_{\mathrm{const}\,1}^{\XB} &:= \Lam{x}{\XB}{\Lam{y}{\XB}{\Pair{x}{(\pauliXXB{y})}}}\\
  O_{\mathrm{id}}^{\XB} &:= \Lam{x}{\XB}{\Lam{y}{\XB}{\cnotXB{x}{y}}}\\
  O_{\mathrm{flip}}^{\XB} &:= \Lam{x}{\XB}{\Lam{y}{\XB}{\cnotXB{x}{(\pauliXXB{y})}}}
\end{align*}

\end{document}
